% =====================================================================
%  R. Nielsen, POPULAR LECTURES, Given at the Invitation of Lecture Societies in Kolding, Fredericia and Aagaard.
%  Copenhagen: Gyldendal (F. Hegel), 1875.  164 pp.  (Folkelige Foredrag efter Indbydelse af Foredragsforeninger i
%  Kolding, Fredericia og Aagaard.)
%  ENGLISH TRANSLATION, from transcription.tex (Danish).
%  GENERATED by ebook-method/folkelige-foredrag/tr_build.py --write from this template and the ten translators'
%  parts in ebook-method/folkelige-foredrag/tr/ (TRANSLATOR-BRIEF.md; terms: GLOSSARY.md, the title and four key
%  terms approved by Hans 2026-10-12).  Conventions: TRANSLATION-PLAYBOOK.md.  Page numbers of the 1875 print in the
%  margin.
% =====================================================================
\documentclass[12pt,a4paper,openany]{book}
\usepackage[utf8]{inputenc}
\usepackage[T1]{fontenc}
\usepackage{libertinus}
\usepackage{libertinust1math}
\usepackage[english]{babel}
\usepackage[protrusion=true,expansion=true]{microtype}
\usepackage{setspace}
\usepackage[margin=1.3in]{geometry}
\usepackage{fancyhdr}
\usepackage{hyperref}
\hypersetup{pdftitle={Popular Lectures}, pdfauthor={Rasmus Nielsen}, hidelinks}
\setlength{\headheight}{15pt}
\pagestyle{fancy}
\fancyhf{}
\fancyhead[LE,RO]{\thepage}
\fancyhead[C]{\textit{R. Nielsen: Popular Lectures (1875)}}
\setstretch{1.3}
\usepackage{marginnote}
\renewcommand{\thefootnote}{\ifcase\value{footnote}\or*)\or**)\or***)\fi}   % as the print: *) and **), restarting each page
\newcommand{\tocentry}[3]{\noindent\hangindent=2em #1 #2\dotfill\ \textbf{#3}\par}
\newcommand{\opage}[1]{\marginnote{\footnotesize\textit{[#1]}}}
% a lecture opens a new page: its heading and place line, centred, then a short rule; the PDF outline takes the title
% from the Contents (the print has none at the head).
\newcounter{lect}
\expandafter\def\csname lecturetitle1\endcsname{The Conditions of the Times: An Old Breach of the Fundamental Law}
\expandafter\def\csname lecturetitle2\endcsname{Christ as the Teacher and Leader of the Peoples}
\expandafter\def\csname lecturetitle3\endcsname{What Christ Teaches about the Historical Future}
\expandafter\def\csname lecturetitle4\endcsname{The Present Age: Light and Shadow}
\expandafter\def\csname lecturetitle5\endcsname{Church and Religion: Fiat and Authority}
\expandafter\def\csname lecturetitle6\endcsname{``Judgments of God'' in History}
\expandafter\def\csname lecturetitle7\endcsname{The Ordering of the Danish People's Church by Law}
\expandafter\def\csname lecturetitle8\endcsname{The University and the Folk High School}
\expandafter\def\csname lecturetitle9\endcsname{A Small People}
\expandafter\def\csname lecturetitle10\endcsname{In Memory of the Battle of Isted, a Festival Address}
\newcommand{\lecture}[2]{\clearpage\stepcounter{lect}\phantomsection\addcontentsline{toc}{chapter}{#1 \csname lecturetitle\arabic{lect}\endcsname}\begin{center}{\large #1}\\[0.4em]{\normalsize\bfseries #2}\\[0.4em]\rule{3em}{0.4pt}\end{center}\medskip\noindent\ignorespaces}

\begin{document}

% split page 20, the title page
\begin{titlepage}
\begin{center}
\vspace*{2em}
{\Huge POPULAR LECTURES.}\\[1.5em]
{\small GIVEN AT THE INVITATION OF LECTURE SOCIETIES IN KOLDING,}\\[0.3em]
{\small FREDERICIA AND AAGAARD}\\[1.2em]
{\small BY}\\[0.6em]
{\large R. NIELSEN.}\\[5em]
{\large\bfseries COPENHAGEN.}\\[0.5em]
PUBLISHED BY THE GYLDENDAL BOOKSHOP (F. HEGEL).\\[0.5em]
{\small PRINTED BY J. JØRGENSEN \& CO.}\\[0.5em]
{\small 1875.}\\[3em]
{\small Translated from the Danish}
\end{center}
\end{titlepage}

% split page 22, the print's Indhold
\begin{center}{\large Contents.}\\[0.4em]\rule{3em}{0.4pt}\end{center}
\hfill{\small\textbf{Page}}\par
\tocentry{First Lecture.}{The Conditions of the Times: An Old Breach of the Fundamental Law}{1}
\tocentry{Second Lecture,}{Christ as the Teacher and Leader of the Peoples}{18}
\tocentry{Third Lecture.}{What Christ Teaches about the Historical Future}{35}
\tocentry{Fourth Lecture.}{The Present Age: Light and Shadow}{54}
\tocentry{Fifth Lecture.}{Church and Religion: Fiat and Authority}{70}
\tocentry{Sixth Lecture.}{``Judgments of God'' in History}{86}
\tocentry{Seventh Lecture.}{The Ordering of the Danish People's Church by Law}{103}
\tocentry{Eighth Lecture.}{The University and the Folk High School}{120}
\tocentry{Ninth Lecture.}{A Small People}{139}
\tocentry{}{In Memory of the Battle of Isted, a Festival Address}{156}
\clearpage

% --- p. 1 ---
\setcounter{footnote}{0}%
\lecture{First Lecture,} {given in Kolding.}
\opage{1}The present age is full of questions, for the world
is full of doubt and unrest. Even if it is quiet on the surface,
yet it seethes in the depths; the stillness is like a
calm at sea. The mighty speak of peace, but behold! they stand
armed to the teeth, they measure one another with their
eyes: a wrinkled brow, a word of threat, a grey cloud
on the forehead of power at once takes shape as a drawn
sword, an omen that awakens fear and sets mind and
thoughts astir. Forebodings weigh heavily, the tension grows;
one explanation takes refuge in another, one
question calls to another: whence does it come?
what does it aim at? whither does it lead?

We look at what lies nearest to us; but this
nearest of ours points to something that is farther off,
this again to something that is yet more distant. Whence
proceeds the first impulse? Where is the source of the movements,
the great driving-wheel that keeps the whole in motion?
We are borne along by a mighty current, but where does it
issue? Where is our goal, our anchorage, our safe
harbor, when the storms come?

To give an all-sided, complete, in every respect reliable
answer to such great questions
% --- p. 2 ---
\setcounter{footnote}{0}%
\opage{2}surpasses human powers, and no mortal shall presume
to read such dark runes. Yet we are not without
guidance. There is in all that is uncertain something
that is unfailingly certain. History is an open
book; we have a revealed word of truth to light
our way: there are guiding stars in the heaven of the spirit.
Everyone who has looked into the deep river and at the same time given heed
to the guiding stars has an intimation of the direction
of the movements and will be able, from his own station, to contribute,
be it only a slight contribution, to the answering of the
questions. I shall endeavor to make a contribution.
What I have to communicate does not directly concern
the burning question of the moment, not such
practical affairs as are otherwise
discussed in popular assemblies and ought to be discussed by men of insight and expert knowledge,
men who with an unbiased eye for the real
whole and for all the particulars know the means at hand
and the nearest ways that lead to the goal.
Outside the political movements of the present
--- but yet by no means indifferent to them ---, I am, by the bent of my life,
referred chiefly to a contemplative apprehension,
a setting up of general points of view, an observation of
the great laws according to which the movements proceed. I
shall set forth my view as plainly and clearly, as
openly and freely as lies in my power. What I then
have to say on this occasion I commend to
the honored assembly's kind attention and
serious examination.

The Preacher says that all is vanity, ``vanity of vanities. One generation goes and one
generation comes, but the earth abides forever. The sun
rises and the sun goes down and hastens to the place
% --- p. 3 ---
\setcounter{footnote}{0}%
\opage{3}from which it is to rise again. All the brooks run to the
sea, but the sea is not full; for the place to which
the rivers go, from there they return again. All
words grow weary; no one can utter it, the eye is not
weary of seeing nor the ear filled with hearing. What
has been, that will be, what has happened, that will
happen, and there is nothing new under the sun.''

If this were the explanation; if the changing
times brought nothing new and distinctive, but
the same and the same in endless repetitions;
if the present were only a ghost from the graves of the past,
how hollow and ghostlike
would the many-colored, richly stirring life of man then
not be! It is the Preacher's last thought that \emph{the fear of God is the beginning of wisdom}; if the thought of the fear of God is taken away,
then the whole is a mere up and down.
To establish this short, pithy thought the
view of the cycle that has been set forth is to be understood as
a prolix introduction; it is the faded experience of the great ones of the world,
the moth-eaten worldly wisdom, the life-weary
wisdom, that speaks.

``What has been, that will be; what has
happened, that will happen;'' that is undeniably true, so
long as the talk is only of the general laws of existence.
There are and there always have been sin and frailty
in the course of the generations' lives; there is and there always has
been a most dubious, here and there revolting discord
between the mighty and the lowly, rich and poor, those who enjoy
and those who suffer, the sorrowful and the glad. What has happened,
that will happen, that the strong overpower the weak, that
power is courted, that its center of gravity is now above
and now below, now with a few and now
% --- p. 4 ---
\setcounter{footnote}{0}%
\opage{4}with the multitude. It has happened and it will happen that those
who come to power misuse it, whether in the name of the individual
or of the multitude, if no counterweight is set up.
It has happened and it will happen that men
envy one another the use of power in the same degree as
they flatter those in power. If power is with the individual,
and the individual is absolute, then the absolute ruler is
hated as a devil and hailed as a god.
If power is with the multitude, then the voice of the majority is the voice of
God, then cunning is afoot to spy out
the passions of the mob, to use its dim instincts and
to make the ignorant believe that they really have a will:
then one flatters and crawls before the multitude. If this is the only
sight on which we are to dwell, then the Preacher is
right: vanity of vanities; there is nothing
new under the sun.

But the sight ought not to stop at the undulating
surface of the movements, at the monotonous rocking up and
down or at the shadows on the broad river
which passing clouds cast; deeper at the bottom stands an
image, the image of wisdom, a finger-post of a Providence
that steers the movements toward a goal. Every age,
when we look more closely, has its labors, its
peculiar tasks and receives in proportion to these its stamp and
its distinctive character. It happens thus, that
one great endeavor gains power and takes up the others
into itself, one thought of life rises and rules the many
thoughts, one direction of the will works its way forward and gains the
upper hand over the many wills, and in this thought of life
and the corresponding will the Governance has laid a
goal.

In the Middle Ages it was the Church that became the
% --- p. 5 ---
\setcounter{footnote}{0}%
\opage{5}highest ruling power. Men might turn whither
they would, try whatever they liked, private
or public, in family life or in civil society,
everywhere they met the ever-watchful eye of ecclesiastical authority:
only what received the Church's blessing
could grow and thrive; what was struck by its
curse had to wither and die. The Middle Ages were
the age of the Church.

Between the Reformation and the French Revolution
(1517---1789) lies a period in which the relations of society
were gradually transformed in such
a manner that the secular power gained a decided preponderance
over the ecclesiastical. It was the State with its absolutism
that, in the form of absolute monarchy, extended
its scepter over high and low. From the Crown
the light of life radiated; the small herbs found shelter, and the
tall plants grew in the shadow of the Throne. It was no
invention of the arbitrariness of an individual, but a natural
consequence of the course of things. The State was the end, all else
only means. The great enjoyed their privileges, in this
there was no wrong; it was, after all, as the supports of the State,
for the State's sake, for the glorification of the State. The multitude
toiled and sighed, in this there was no wrong; it was, after all, in
the service of the State, for the State's sake, for the good of the whole.
Thus the absolutism of the State is the opposite of the Church's; no one
can mistake the two ages.

But now our own age, the present in which we
live, what character has it, and what name will
it bear? It is still in its beginning, in its early
spring. But in the early spring the weather is
apt to be unsettled: sunshine alternates with drizzle and sleet.
Many seed-corns and many kinds of seed lie hidden in
% --- p. 6 ---
\setcounter{footnote}{0}%
\opage{6}the earth, and no one can see what it is until it has
sprouted and come up. Much of what comes up
first is not always the best, for where the soil
is not cultivated and cleared, the weeds grow
as luxuriantly as the corn. None the less the age has
already such firm and distinctive features that
we cannot be in doubt as to its character and
nature. Its birth and descent are no accident;
it is the true child of its past, born of the womb of the dispensations
of Providence and baptized in the stream of events. It was
in '48, a stirring year, that the thread snapped, and the newborn
stirred vigorously in its swaddling-clothes. On a stormy day
it was held over the font; black, lightning-laden clouds
darkened the heavens, but the clouds parted, and the sun broke
forth, when the age received the new, yet age-old name: \emph{The Time of the People, the Time of the Peoples!} Great expectations,
rich hopes attach themselves to this
name; may the young age but grasp the greatness
of its task; may it keep its coat of arms unspotted
and bright and remain true to its mission; may its baptism of the people
be more than a baptism of water, may it do
honor to its name!

A difficulty at once presents itself when the talk
is of the Time of the People, for by what is it actually to be
known? When the highest governance has its seat in
the Church, then it is easy to say what the Church is. It
is like the city on the mountain, a grand institution
with a distinctive doctrine and laws. It presents itself to
view as a divine order of things; it knows
its friends and it knows its foes: it can
be stated exactly how far its dominion extends.
The same holds of the State, when it has received its
% --- p. 7 ---
\setcounter{footnote}{0}%
\opage{7}proper form as the sole ruling power on earth
But to show exactly where the people is, and what it
is, to make due distinction between its friends and
foes, is by no means an easy task. Where is
the people? In the whole mass of the population. The multitude must
be included, if the people is to be included; and yet the multitude is not
the people: the people as a people stands high above the multitude.
The multitude must be kept in check and ruled; but the people must rule
itself and its multitude. Without popular self-government
no true popular life; but how is popular self-government
possible?

Under the absolutism of the Church popular life could not
thrive; for self-government lay exclusively in the Church's
hand. The Church alone had to command, the peoples had to
believe and obey; with that the matter was settled. In such a
state of blind submission the people had, at least
in part, to sink down to being a mere multitude.
Likewise under the absolutism of the State. In this respect it is
instructive to note the likeness between
the two great powers, the Church and the State, as regards the conception
of the people.

When the multitude, the great mob, sinks on its knees before
the Pope, believes in his miracles and goes devoutly in
the procession, then it is not the multitude that kneels,
no, it is the people, the pious, believing people; and the Church
is popular, it loves the people. But if a deeper religious
need awakens in the people, if a feeling stirs in its
mind that the inward is the goal, that the
divine is profaned by outward mummery: then
the Church's eye grows dark; it does not see the people at all then,
it sees only the multitude, the unbelieving mob, damnable
heretics, who are to be rooted out.

% --- p. 8 ---
\setcounter{footnote}{0}%
\opage{8}When, in the age of the State, the ruler showed himself
in public, and the mob shouted Hurrah, and the multitude rejoiced,
it was not the mob, not the multitude, that shouted
Hurrah, no! it was the good citizens, the dear, faithful
people; and the sole ruler was popular, he loved
his people, after all. But if there awoke in the people a question
of right, a feeling of oppression and revolting
inequality, and an uncanny murmuring was heard, then the
police had to step forward. It was naturally not the people that
murmured, it was only the multitude, the rude mob, and
the multitude, the rude mob, was to be chastised.

If it were only such self-governing disciplinarians
as the absolute Church or the absolute
State that the people had against it, it would long since have
outgrown the rod; but it has an adversary to fight
who is far more dangerous, because he is not merely
outward, but also inward. He does not announce himself
as a foe, but as a friend and brother; his nature is
so naturally grown together with the existence of the people that
at any moment, when occasion offers, he can step forward and
say: it is I who am the people. This false friend,
this dangerous brother is --- \emph{the multitude.} To wish to let
the multitude rule, or to set about introducing popular
self-government after the mind of the multitude, is the most absurd
of all things, for the multitude cannot rule itself: its
nature is to be unruly. Let it be well observed
that the multitude has precisely the same conception of right
and might as the most revolting tyranny of a single man.
The same battle-cry, the same watchword; We are strong, for we
are many; the watchword is: \emph{the right of the stronger.}

From this point of view, history gives a sorrowful
lesson. Every time a popular will is in
% --- p. 9 ---
\setcounter{footnote}{0}%
\opage{9}the process of forming itself, a will that equally embraces
all, small and great, a will that wills the good and
right, because it truly wills what is genuinely human,
there is at once a craftiness in the game, a
cunning deceit, which disrupts the whole. Now it is
a gifted man of power who in the name of the people wins the multitude
to his side and brings order after his own fashion by
taking the lion's share and doing violence to the people. Now it is
the multitude which, singing hymns of freedom,
tears asunder the bonds of order and at last, despite cries and
shouts for law and right, ends in a disorder from which
tyranny is born. It is a fate that has weighed heavily on peoples
and generations for millennia on end. Whence
can it come? From an old sin, from inner
discord between will and will, between natural will, self-willed
instinctual will, and the will of the spirit.

There is an old account outstanding between the spirit
and the race, a breach of the fundamental law from time immemorial, a rebellion
against the laws of the spirit, with which history begins.
The events of the world form a web that conceals this;
but the revealed word draws the veil aside
and shows us what the spirit wills. In the morning of the ages
judgment was held upon the egoism of the men of power as well as
upon the multitude.

In the sixth chapter of the first book of Moses the following is related.
``And it came to pass, when men began to multiply
on the face of the earth, and daughters were born to them, that
the sons of God saw the daughters of men, that they were
fair, and they took them wives of all whom they
chose. The giants were on the earth in those days, and likewise
afterwards, after the sons of God had gone in to
% --- p. 10 ---
\setcounter{footnote}{0}%
\opage{10}the daughters of men. These are the mighty ones, who from
ancient times were men of renown.''

There are now those who have thought that there was in this something % print: en (misprint for er)
veiled and mythical, something obscure and heathen, that
the sons of God (the sons of Elohim) must be either
sons of lower gods or perhaps fallen angels,
demons. These ruminations are superfluous; what the Word
reveals is plain enough. The sons and daughters
of whom it is here spoken are of common descent, all, naturally
understood, children of the same race, all, spiritually understood,
children of God, created in God's image. Adam is called
the son of God (Luke 3:38); the mighty on earth, princes
and judges, are called gods. ``God judges
among the gods'' (Ps. 82:1). ``I said: Ye are
gods, all of you are children of the Most High, but now
ye shall die like men'' (Ps. 82:6---7; John
10:34---35). The same human nature is here
presented from two sides: in the strength of the man is seen the
divine, in the beauty of the woman the human. % print: Mennneskelige (misprint)
Strong men and fair women, that is the race with
its rich endowment, a vigorous race, which neither the toilsomeness of
labor nor the yoke of bondage has yet
bent.

What comes next and is more surprising is the look
with which Jehovah regards this race; for it is not
in mildness and grace, but in indignation and wrath.
In this men have above all wished to glimpse a remnant of heathenism,
a mythical, still half idolatrous conception of
the nature of the spirit. Jehovah, it is said, according to the whole
narrative has evidently become envious and afraid. He
is envious of this surpassing fullness of strength and
beauty; he is afraid that these men of renown,
% --- p. 11 ---
\setcounter{footnote}{0}%
\opage{11}these mighty giants, will --- like the Titans
in the Greek myth --- try to storm his
heaven and overthrow his throne. That is a small
example of the supposedly higher knowledge of the present age about
the divine, but for the rest an enormous misunderstanding.
The wrath of the spirit of which it is here spoken is not directed against
the rich gifts of nature with which the Creator
himself from the beginning has endowed the human race;
no, it is directed against the misuse of these gifts. Jehovah sees
that the richly endowed race, which is now to try its
powers and tread the path of freedom, will not use
the gifts in the service of the spirit; he sees these giants
give themselves up to selfish natural impulses, he sees them
wildly working for their own renown and using
the right of the stronger; and the Judge's voice is heard:
``My spirit shall not rule among men forever,
since he is flesh, and his days shall be
one hundred and twenty years.''

When the powers of spirit in the gifted are transformed
into powers of flesh, when an unbridled self-will employs
these powers only to sate rapacity,
the hunger of ambition and the lust of enjoyment, then the
spirit turns away with loathing from the strong; their
boundary is set, their days are numbered, their doom is
pronounced. They are terrible words, the words of Revelation about
the giants' stock with its great deeds: ``Jehovah saw
that the wickedness of man was great on the earth, and all
the devisings of his heart only evil every day.'' The
holy wrath of the spirit cannot be depicted in stronger colors
than when it is said: ``it repented Jehovah that he had
made man on the earth, and it grieved him in
his heart.'' The strife between the will of the spirit and the
% --- p. 12 ---
\setcounter{footnote}{0}%
\opage{12}self-willed instinctual will of man, strong in nature, is already
here at its height. We may be assured
that what the spirit already at the beginning of history
judged of the brutal renown won by
violence and wrong, the same it judges even to this very day.
Just as those first attempts of the giants, by which
the earth was ``corrupt before God'' and ``was full
of violence,'' dissolved and were washed away in the Flood,
so history in its continued course has been a
flood of destruction for all great works which were not
built upon truth and right, and had not the consecration
of the spirit.

And now the great work of the multitude. The passage is found in the eleventh
chapter of the first book of Moses. ``And the whole earth was one
tongue and one language. And it came to pass, when the peoples
journeyed from the east, that they found a valley in the land of Shinar,
and they dwelt there. And they said, one to another,
Come, let us make bricks, and burn a burning;
and brick served them for stone and earth-pitch for
mortar. And they said: Come! let us build a city and a
tower, and its top shall reach to heaven, and let us
make us a name, lest we be scattered over the
face of the earth.''

In this event there is something which at first
glance must awaken our highest astonishment. What
people is it of whom it is here spoken? ``The families of Noah's sons;'' thus Jehovah's,
the spirit's, own people, for Jehovah
had made a covenant with Noah. What evil can
there then be in this, that these many set about working
together and building a city and a tower? Here
are no giants striking down among dwarfs,
no men of violence overawing the weak, no
% --- p. 13 ---
\setcounter{footnote}{0}%
\opage{13}slave-drivers, who whip the bondmen to gather straw
and make bricks. On the contrary! All goes so evenly, so
briskly and peacefully, that the immense multitude of busy hands
makes the impression of a great, harmonious workers' association,
in which the one unceasingly encourages the other with
a: Come, come! It is, after all, a fair sight, a grand
employment of the forces of the people, who work, not for
dispersion, but for unity and brotherly cohesion:
what can be objected to this? And yet the spirit is
wroth; Jehovah comes down and confounds the workers
and disrupts their work. How is it possible
that confusion and disruption can proceed from the spirit?
``And Jehovah said: Behold, it is one people and one tongue
for all, and this they have begun to do, and
now nothing will be able to be denied them of what they purpose
to do.'' As dark, yes, as meaningless as these
words would be, if we understood them in a manner so
childish that the Almighty were really minded to scatter the peoples and weaken
the popular will for fear of the high tower; so clear and
intelligible are they, when we grasp them as an expression of
the truth in the thought of the all-governing will. The tower,
whose top was to reach to heaven, was by no means
intended to pierce the vault of heaven; it
was a great work which was to have significance on earth alone,
a towering point of union, which, by
defying the changes of the times, was to be an eloquent
testimony to what human powers can accomplish.

The working tribes on this occasion think
not at all of either insulting heaven or
defiantly challenging the spirit; their thought is much rather
this, that they can best do the whole thing themselves; so
% --- p. 14 ---
\setcounter{footnote}{0}%
\opage{14}they let heaven be heaven, and the spirit spirit, and set about
the work. Age-old and yet as of yesterday!
The spirit is not, if we may so say, angered
on its own account, because it is overlooked by men
and does not enjoy the honor of being included. It is
for men's sake, on behalf of the people and of popular life,
that the spirit is wroth, for without the cooperation of the spirit
popular life is impossible.

The busy people who build in the valley of Shinar have
renounced what in a deeper sense can
recall the claim of the spirit, and are on the point of
sinking down to a mere multitude. They have forgotten
what is original, their forefather's relation to Jehovah
and the sign of the covenant that was set in the cloud. What is
a rainbow, after all? A fleeting gleam of light without reality,
without reliability; so it stands, in
their opinion, with the spirit also. No, a tower is
something quite otherwise real and lasting; and then
it is besides the work of their own hands, the
immortal work by which they ``make themselves a name.'' Jehovah's
word, his sign of omnipotence and his governing wisdom,
of that one did not know what it was; it was something one
could not get hold of at all, it was all fleeting
and airy like the rainbow. Bricks, clay and
earth-pitch, of these one knew what they were, they were things of
value, something one could feel and lay hold of. Thus
this teeming multitude with its bustle, its
capacity for work and colossal striving has denied the deepest
presuppositions of human life, broken with the spirit,
broken the holy covenant with Jehovah.

This is the condition of life and freedom: if the original is denied,
if the spirit is not allowed to be included
% --- p. 15 ---
\setcounter{footnote}{0}%
\opage{15}in living faith, in inner appropriation and free obedience, then
the Invisible One gathers the cloud of his wrath over the anthill
of spiritless self-will, disrupts its work and
scatters what it has gathered. It is impossible that there
in a human society can be true unity
without a deeper ground of unity; but the deeper
ground of unity is a ground of will. The ground of unity is in
the recognition of the good and right; the true decisive
will is the will of the spirit. Where the ground and the will are lacking,
the unity is imagined and transitory; minds
change, instinctual will expresses itself in shifting outbreaks of selfish
inclinations and passions: disunity comes
suddenly upon the united from behind, when they least expect it.

By the ground of unity and the kind of will one is to
know the difference between the people and the multitude. If
the great plan with the great tower in the valley of Shinar
had really had a true will-demand deeply rooted
in human nature, then we should have to admit without denial
that here was more than a multitude, that the many
tribes were one people, an enterprising, capable, aspiring
people, which worked at something exalted and noble;
the God who out of arbitrary caprice wished to disrupt a
noble human work we should have to declare outright
an idol. But the Babylonian tower-builders have no
such ground of unity; had they such a ground,
then they must necessarily have had the spirit along with them, but
it is precisely the spirit that is shut out. It is not
the people, it is the multitude, that has here banded together
to complete an outward extraordinary thing, which could
really catch the eye. Let us look at the matter naturally and
ask: how could the unity about so immense a
joint undertaking be thought to hold? Were
% --- p. 16 ---
\setcounter{footnote}{0}%
\opage{16}there no leaders, then there was no plan of building either;
everyone had to go his own way, and that could
then not possibly lead to unity. Were there leaders, chief overseers,
then these must either apply compulsion
or persuasion. Compulsion they could not apply, for
when the multitude is to be put under compulsion, it gives itself out
at once as the people, and then it is said straightway that
violence is done to the people. The leaders have therefore
had to limit themselves to applying persuasion. Now
nothing is surely easier than to persuade the multitude
to do something which at bottom it has a mind to,
something to which it believes it feels an inner irresistible
urge; but to persuade the multitude to do something
to which it has no mind, that is the most difficult of
all. It has therefore, one may imagine, when the work
began to grow hard and the tower crooked in the valley of
Shinar, come to a serious opposition against the chief overseers,
to debates so motley, broad and
long-drawn-out, that the languages were confounded, and the speakers at last
did not understand one another.

And what is it, then, that the spirit does? It hastens
the decision and pronounces the judgment. It is its will that
the discord hidden in the nature of self-will shall become
manifest; it makes it known of what kind such a
cohesion, such a unity is. ``Come,
let us go down and confound their tongue, so
that the one shall not understand the other's tongue.
--- And Jehovah scattered them abroad over the face of the whole earth,
and they left off to build the city.''

We learn from this that early there came a breach
between spirit and people. As a consequence of this breach
the many tribes with the many tongues
% --- p. 17 ---
\setcounter{footnote}{0}%
\opage{17}had then to wander each his own ways --- to victories or defeats, exaltation
or humiliation, honor or disgrace. The will of the spirit
has not been satisfied; there is wrath in the look
of the spirit, and yet the spirit is love: there is a glimmering of
spirit in the natural ground of the peoples. The gleam is different
in proportion to the particular endowment of the individual peoples
and their peculiar gifts; but something
divine stirs in them all. They yearn, they struggle, they
grope and stumble; but the movement goes on with unceasing
speed --- ever forward! On a thousand reefs,
of the power of the great, of the tyranny of order and the wild
multitude's hatred of order, the ship of the people has hitherto been stranded:
the shores of memory are strewn with wrecks. It
is the old offence, the old breach of the fundamental law, which
again and again has taken its revenge. Always the right of the stronger!
``But then righteousness itself must at last
become the stronger''; that is the inmost thought
in the mind of the peoples. It is this sight --- a rainbow
in the cloud --- that they yearn for; it is the goal of reality
toward which they strive. After the confusions, sufferings and defeats
of struggles a thousand years long, a turning-point has at last
come. The sun of enlightenment has in our
days risen so high in the heavens that we can see
where we are. We cannot yet see the issue of the struggle,
far from it! but we can nevertheless see the goal. If
the goal is to be reached, there must be will. It is not the will of the individual,
not the will of the multitude or majority will,
that is in question; no, it is the popular will. But the
true popular will is inseparable from the will of the spirit; if
the people is to conquer in the hard struggle and reach its goal, then
it must experience the power of consecration: it must be baptized with
spirit.

% --- p. 18 ---
\setcounter{footnote}{0}%
\lecture{Second Lecture,} {given in Fredericia.}
\opage{18}It is taught in Holy Scripture that the human race dwells on the whole circle of the earth of one blood. But it is also related how it came about that the one lineage became many lineages, the one family of peoples many tribes of peoples, the one original, common language many different tongues. The dispersion and the difference are now so great that the investigators of our scientific age must begin their inquiries afresh and put to themselves the question: is the whole human race really of one descent? Can all these white, black, olive-yellow and copper-red human beings be supposed to descend from one pair of parents? The inequality is more than outward, it is not confined to the skin and the hair, to the shape of the skull and of the face, the iris of the eye, and the like; it is just as great in the inward as in the outward. Some tribes of people are wonderfully gifted, others not; some distinguish themselves by a high degree of enlightenment and culture, others are still on the lowest step of rudeness. Even among the gifted the inequality is great, the Indians, the Egyptians, the Greeks, the Romans and the Northmen have each had
% --- p. 19 ---
\setcounter{footnote}{0}%
\opage{19}their own direction of spirit, each its own peculiar view of life. This is seen especially in their gods, for they are gods of the peoples; these gods of the peoples were, so to speak, the holders of the spiritual capital on which the peoples lived.

What a difference between two such circles of gods as the Greek and the Norse; with all that they have in common, there is, as regards depth and inwardness, a considerable distance between them: Odin does not understand Zeus, Aphrodite and Freja cannot live together under one roof. These gods are the original bearers of the peculiarity of a people; when they age, grow pale and sink, the strength of popular life withers away. If now what we call the spirit of the people were only a hazy fancy, a web of unclear stirrings and strong dreams: then we might indeed explain the gods to ourselves as the first bright figures in which the human being, the mind overfilled with thoughts, presentiments and half-understood longings, found a colorful, manifold reflection of its own moved interior; but then we should also at the same stroke have to reduce the spirits of the people to vanishing spirits of the times. Either we should then have to suppose that popular life in the present could thrive and flourish without spirit, or expect that the old faith of the people with the old gods would arise from the dead and displace the new enlightenment. The first is just as unreasonable as the last. We must therefore regard the spirits of the people from heathen days as spirit-like powers, scattered gleams and dim proclamations of the one eternal Spirit of truth, whose light has never ceased to shine in the darkness.

When light passes through layers of cloud, the rays are broken and the colors many; so too when the spiritual light is to make its way through
% --- p. 20 ---
\setcounter{footnote}{0}%
\opage{20}the unequal natural ground of the peoples. We may very well ascribe validity to the spirits of the people without making them persons or setting them up as independent beings. The Spirit of truth will not efface the peculiarity of the peoples, but it will remove what is foreign, what is hostile, and, as it is permitted to rule, ennoble the distinctive and glorify the peculiar. The Spirit of truth is personal in the true religion; its human revelation is Christ. In order to see how the Spirit of Christ is related to the spirit of the people, we shall in this lecture consider Christ and the people.

Christ and the people: this relation is not the same as: Christ and the congregation. For the faith of the congregation, the Only-begotten of the Father stands as the revealer of God in flesh and blood, the Reconciler, the Crucified One, who suffered and died for the sins of the world, rose from the grave and entered into glory; Christ does not stand in so inward a relation to the people as a people. The popular eye sees the divine in him from the human side, regards him as a heroic figure, an eminent, unique character in the history of the world, and dwells on his public life. In his public life Christ shows himself as the man of the people, its guide, teacher and leader. What he teaches the people, what he judges about the people, what he will do for the people and what he wills that the people shall do, thus becomes, from this point of view, the main thing.

Immediately at his first appearance Christ sends out a light which separates the multitude and the people; the multitude seeks him, and in the multitude he seeks the people. His Sermon on the Mount is a speech to the multitude and yet in the highest sense a speech of the people. Here is something that
% --- p. 21 ---
\setcounter{footnote}{0}%
\opage{21}distinguishes the speaker of the people in Galilee from the orators in Athens and Rome. The great orators in the capitals of the world so readily began by reminding the assembled of their natural rights and claims. ``Athenians! --- Quirites! Free citizens, freeborn men, you know indeed that you yourselves have power to make laws; you know and feel that it is unworthy of you to obey other laws than those you have yourselves made.'' By such speeches the passions were set in motion, and the multitude took courage. Of course, it was not directly the multitude that was spoken of; it was expressly, with all the ring and force of a fully resounding eloquence, to the people that the speech was addressed, about the people, solely and only about the people. None the less there was something questionable behind it. There was something political, not to say Polis-ish, in the whole procedure. It looked to the life in the eyes of the multitude as if freedom came with a leap, as if it sprang out of the head of the orator, just as the full-grown, fully armed Minerva once sprang out of the head of Jupiter. Thus the multitude broke loose, but the people was and remained a prisoner, it merely changed its leading-strings: it sounded like freedom, but ended in slavish dependence.

The speaker of the people in Galilee goes about it differently. In order to set the people free, he begins by binding the multitude. ``Think not that I am come to destroy the law or the prophets; I am not come to destroy, but to fulfil. For verily I say unto you: till heaven and earth pass, one jot or one tittle shall in no wise pass from the law, till all be fulfilled.'' In thus speaking from the mount, he speaks from the height of the spirit, of eternity --- not of civil laws, which can be transformed
% --- p. 22 ---
\setcounter{footnote}{0}%
\opage{22}and changed, but of the unchangeable law, the holy fundamental law that is written in the hearts of men, the law of truth, of righteousness and of love, which endures though heaven and earth pass away. And he sharpens this law and makes it still stricter. ``Ye have heard what was said to them of old; but I say unto you'' . . . . A sum of outward commandments is not enough; what matters essentially is the inward. In thus impressing upon the listening throngs what it is that is originally required, an inner change, a thoroughgoing transformation of the mind: not to will because there is compulsion, but sincerely to will because it is a need of the heart, the strictness itself becomes gentleness. Thus he transforms the law from a law of the letter into a law of the spirit, from a law of compulsion into a law of freedom, and helps the people to freedom.

One might fear that the multitude would murmur and feel offended at so bold a speech from a single defenseless man --- for that Christ at that time could not stand before the eyes of the many in the light in which he now stands for us goes without saying. He stood among them as one who wished to be judged by his words and deeds alone. But he knew how to lay hold of the consciences and transform the multitude into a people. The impression of the speech is --- significantly enough --- rendered in the words: ``The people were astonished exceedingly at his doctrine, for he taught them as one that had authority, and not as the scribes.'' As one that had authority, how strange! Had the scribes, then, no authority? Judging by externals, it should surely be they alone who had the authority; they were indeed the authorities of the land, the interpreters and upholders of the law, the appointed guardians of the people.
% --- p. 23 ---
\setcounter{footnote}{0}%
\opage{23}But their authority was only a semblance. The people had a feeling that these scribes, with all their word-splitting and legal learning without vitality, were without deeper sense for truth, without spiritual superiority. The whole existing order was an empty externality, something outlived, which clung to forms and conventions.

In order to restore the true authority, Christ must begin from the foundation by tearing down the false. He does it as the leader of the people, without any final regard, without circumlocution, without sparing. He tears the mask from hypocrisy; in the hearing of the whole people he cries woe over ``the whited sepulchres.'' His words could fall like the dew of heaven, cooling the house like mild airs that stir about the child of the mountain; but when insolence wounded him, and the holy wrath broke forth, then his speech was a thunder, and the lightnings crackled. ``The scribes and the Pharisees sit in Moses' seat. All therefore whatsoever they bid you observe, that observe and do; but do not ye after their works, for they say, and do not. For they bind heavy burdens, grievous to be borne, and lay them on men's shoulders, but they themselves will not move them with one finger.''

We can understand that the scribes and the Pharisees, the men of order, of sobriety and of the old godliness, must hate the new prophet with a thorough hatred. Seize him openly they dared not, catch him in words they could not. They therefore lay in wait for an opportunity, and meanwhile gave vent to their feelings in various judgments. Soon it was said, with the assumed superiority of a learned distinction: he is an idiot, an ignorant man, who
% --- p. 24 ---
\setcounter{footnote}{0}%
\opage{24}does not know the Scriptures; soon with brutal slander: he is a glutton and a drinker and consorts with the vulgar rabble; soon with undisguised wrath: he hath a devil, he is a past master of seducing the people, a devilish man. Others, who wished to pass for being more fair-minded, thought perhaps that what he said about the founding of his kingdom, his kingship of David and his throne was not so straightforward to understand according to the letter. But then one had to deplore deeply that he filled the multitude with delusions and spoke to the people of something that far exceeded its power of comprehension.

The words of hatred and the outbursts of blindness do not trouble us; but the question whether Christ as the teacher of the people could really be understood by the people weighs upon our mind. If the opposite were the case, not only would the work of Christ be in vain; but those who even in our days fear everything from a popular enlightenment and earnestly warn against speaking to the people of lofty things would indeed be perfectly right. We can console ourselves over this misunderstanding. The orators in Athens and Rome hardly always knew what they were doing when they spoke to the people and meant the multitude; but the speaker of the people in Galilee is clearly aware of what he does when he speaks to the multitude and means the people: ``he that hath ears to hear, let him hear.''

Christ knew full well that in the whole population there were not only noble but also ignoble elements; he knew the hardness of hearts, of those who see and yet do not see, who hear and yet do not hear. He compares the word of his kingdom to a sowing and speaks of the sower: some falls by the wayside, some on stony ground, some among thorns, but there is also
% --- p. 25 ---
\setcounter{footnote}{0}%
\opage{25}some that falls on good ground and bears, thirtyfold, sixtyfold, a hundredfold. All the people hung upon him and heard him, and he knew well that they could understand him when they sincerely wished to. When not only the twelve, but afterwards also the seventy disciples whom he had sent out returned and joyfully announced to him how the word had been received, ``he rejoiced in spirit and said: I thank thee, O Father, Lord of heaven and earth, that thou hast hid these things from the wise and prudent, and hast revealed them unto those under age.'' Here is, for the cause of popular life, a turning point of the times. There is something that surpasses the power of comprehension of the common people, and that is in learning and science: the course of the stars, the long history of the development of the earth, the natural descent of living generations from lower to higher in ascending links; there is in these enlightenments something glorious and great, something that popular society cannot do without. And yet the common people cannot understand it, a proof that it is indeed a higher thing, but not the highest. Thus: the higher the whole people cannot understand, but with inwardness of soul, in simplicity of heart, it can understand the highest. Wonderful wisdom! It is this thought of life that came into the world with Christ; it is he who vouches for an equal understanding of the great principles of equality and freedom, an understanding that is essentially compatible with the difference between the wise and the simple, the learned and the unlearned. By putting this thought of life into force Christ has founded popular enlightenment and taken the first decisive step toward the liberation of the peoples.

The popular enlightenment that Christ has founded has, however,
% --- p. 26 ---
\setcounter{footnote}{0}%
\opage{26}been misunderstood from two opposing points of view, one one-sidedly worldly and one just as one-sidedly godly.

From the worldly side it has indeed been willing to concede that there is knowledge of human nature, feeling and deep emotion in everything that the Nazarene says. His words about the lilies of the field and the birds under heaven are fitted to bring a beneficent relief to the mind depressed by cares of livelihood and anxieties. But, it is objected, moods are not enough. We live on the earth in earthly conditions. Enterprising people must attend to their trade; they must work and spin, they must sow and reap and gather into barns, they cannot, like the lilies or like the birds, be free of care for the morrow. It would thus --- so the worldly-wise think --- be a deficiency in Christ's teaching that he took no account of the practical, the nearest, the useful. Instead of letting his hearers consider the lilies of the field and the birds under heaven, he ought at once to have taught them to read, write and reckon, to have initiated them into the economy of the state, especially into politics; all of which must be regarded as useful and necessary. A great misunderstanding. What is here adduced as a deficiency is precisely perfection. What Christ imparts to the people contains precisely the presupposition and condition for all that is practical, all that is useful, in which there is blessing.

There is something unhealthy in the feverish zeal and ruthless greed with which a worldly mind can give itself up to gain and gain in agriculture, in handicraft, in trade and industry, in short, to everything that procures advantage, embellishment and enjoyment of the present
% --- p. 27 ---
\setcounter{footnote}{0}%
\opage{27}life. Especially, nothing is more dangerous for a people that is to learn to use its freedom than to throw itself violently and headlong into a political commotion, supposing that politics is the highest thing on earth. If the mind is not filled and moved by something higher, it will easily happen that self-righteousness, willfulness and unclean passions confuse the sight and drag even the better away to a sorry imitation of the artifices that stain the politics of the great, that web of injustices in which servitude sits.

Christ's word is simple and deep: ``Seek ye first the kingdom of God, and his righteousness, and all these things that ye have need of shall be added unto you.'' How this word is to be taken in the congregation shall not be treated here. We take it in this connection as a word to the people; we shall convince ourselves that it finds a great, decisive practical application. Far from paralyzing a people's powers of labor or making it indifferent to the temporal tasks, this word is on the contrary the only one that gives a full answer to a decisive practical question. What does a people need in order to become the most enterprising and industrious in the world, to become a people that has a sense for everything exalted and beautiful, a hard-working people that is able to bear even heavy burdens? The kingdom of God --- the kingdom of the spirit.

To be in the kingdom of the spirit and to feel oneself borne by an invisible power is, in other words, to feel oneself lifted above the temporal. But to feel oneself lifted above the temporal is a condition for being able to sink oneself into it and work with it without danger of going under. He who knows within himself that, even if misfortune breaks in, he has a greater capital in reserve, can venture the utmost, and he who need not be afraid of
% --- p. 28 ---
\setcounter{footnote}{0}%
\opage{28}the consequences can seize the moment. We take life as it is; let us consider those who enjoy. Who enjoys the goods of temporality best of all, the one who sets his only goal in enjoyment, or the one who in the midst of enjoyment is raised above enjoyment and has spirit and strength to break off whatever moment it may be? Surely the latter; for the first is a desperate soul, who laments when the possibility of enjoyment ceases, and dies of boredom if it is bored enough to hold out with him. Let us consider those who fight. Who is the bravest, the one to whom wounds and death are the greatest misfortune, or the one who knows himself raised above the outcome of the strife and can fall, if so it must be? The latter is in the midst of danger raised above danger, and danger itself is almost afraid of him; the first is a wretch, whom danger is always after and as it were takes its pleasure in cowing, because he so inwardly wishes to run from it. Let us then take the matter in earnest, think of Christ's word and consider the conditions of temporal life. Hindrances and entanglements, resistance and barriers everywhere! not only for the small, but also for the great; yet by the everyday rule especially for the small.

A poor man has many burdens to bear, how shall he hold out? He steels himself, he pulls himself together, he wishes in his strife with the temporal to lift himself above the temporal. If now in his mind there is nothing but temporal strength and temporal desire, what will the end be? The temporal is overcome by the temporal, loses courage and lets his hands sink. But if there is in his soul a spark of eternity, a spark of spirit, then he is nourished by the invisible
% --- p. 29 ---
\setcounter{footnote}{0}%
\opage{29}bread, then he holds out; his courage transforms itself into patience, and the courage of patience is invincible.

A small people must in the great relations of the world feel itself in many respects like the poor man. Even if it does not exactly suffer from cares of livelihood, indeed even if it abounds in well-being: there is nevertheless a pressure, an anxiety, that will not cease. The superior power may come with violence and plunder, how shall the small people find means to defend its right, to protect its independence and to keep up its courage? If now in such a people there is no spirit, no higher power of will, no thought of eternity; if all is only temporal, temporal calculation of temporal gain and temporal loss, then it will happen that the temporal is overcome by the temporal: the exertion flags, the powers slacken, despondency whispers: what can it help? % print: hiælpe (misprint)
But if the spirit is present, a thought of eternity, an uprising of the will, which seizes the mind of the people and lifts it above temporal struggles, then the conditions of existence change. Everything, says the cold understanding, goes according to laws. Right. But --- let us not forget it --- there are in the world two kinds of laws, laws of necessity and laws of freedom, laws of nature and laws of spirit. However inwardly they hang together, they are nevertheless different. To want to build on freedom and yet believe in blind necessity is a contradiction, is: to want to rule one's own life and yet submit oneself to fate. Christ's word will break the yoke of fate and place the strongly willing people under Providence. The governing Providence, which lets its sun rise on the evil and the good, transforms courage into patience; the fatherly Providence, without whose will not a sparrow falls to the ground,
% --- p. 30 ---
\setcounter{footnote}{0}%
\opage{30}transforms patience into courage: ``fear not, ye are of more value than many sparrows.''

The light-shunning godliness sees the matter from another side; it will not admit at all that the spirit in Christ's word has anything to do with popular life. It understands the admonition: ``Seek ye first the kingdom of God,'' as if the highest were attained when one simply cast the earthly away and concerned oneself with nothing else than weeping and wailing over the sin and folly of the world. We dare not forget the holy word: ``Blessed are they that mourn''; but the sorrow in which there is blessedness cannot possibly consist in extinguishing all the lights of gladness merely to look sourly at the world. If it were in itself a sin to use the temporal and to work with the worldly, then not only popular life in the special sense, but human life, from whatever side one may consider it, would be made impossible at a single stroke. The sickly piety has never been able to carry through its view of life, has never managed to give up everything temporal. It has remained standing at a half-way state without giving itself an account of the claim of the spiritual upon the natural; it has wished to take part and yet not to take part in the advancing society, it has out of sheer decorum not known whether it dared to live among real human beings. The hermits and monks of the Middle Ages proceeded with quite another consistency ; but history teaches us how it ended.

One of two: either Christ has wished to help popular life forward too, or he has wished unconditionally to halt it; he was --- let us speak humanly of him --- not the man to stop half-way. Had the latter lain in his plan, he would surely
% --- p. 31 ---
\setcounter{footnote}{0}%
\opage{31}not have rejected the fasting holiness; he would have broken the staff over the whole lawful order of society, and above all: he would have begun by forbidding marriage and banning the propagation of the race. When now he does exactly the opposite. When now he regards the human order, the given relations between princes and people, authorities and subjects, parents and children, masters and servants as something that follows naturally of itself, indeed as a divine order; when he again and again uses these relations as a ground for his parables of the kingdom of heaven; when he calmly and freely goes in and out of the houses of men, looks upon their goings and doings with mild eyes, takes part in their feasts, helps the suffering, comforts the sorrowful, blesses the little children and lets the wine of gladness be poured for bride and bridegroom: who then dares to say that Christ has broken the staff over the temporal, the human, the popular? To bless family life is to bless popular life, for what is popular life but a widely ramified family life embracing manifold families?

What the angels proclaimed to the shepherds, ``a great joy, which shall be to all people,'' that Christ as the teacher and leader of the peoples has in the highest measure wished to make truth. That he, who was the Savior of souls, also wished to be the Savior of his people, he has himself expressly shown. It lay heavy on his mind how it would go with the small people whose great hour was drawing near. It was a vigorous but refractory people, which lived on great memories. He saw it split into parties, Pharisees, Sadducees, Herodians, the Roman-minded, those who were zealous for the old and
% --- p. 32 ---
\setcounter{footnote}{0}%
\opage{32}held to the national, those who yearned for the new and admired the foreign; he saw its teachers without spirit, its authorities without authority, the common multitude like sheep scattered on the mountains without a shepherd; he saw the mighty, lurking enemy, the Roman eagle with the sharp claws, which was soon to seize its prey; he saw it all, and he sorrowed.

At his entry into Jerusalem, when the disciples rejoiced and the multitude cried Hosanna, the festival of the people was no festival of gladness for him. ``And when he was come near, he beheld the city, and wept over it, saying: If thou hadst known, even thou, at least in this thy day, the things which belong unto thy peace. But now they are hid from thine eyes. For the days shall come upon thee, that thine enemies shall cast a trench about thee, and compass thee round, and keep thee in on every side. And they shall lay thee even with the ground, and thy children within thee, and they shall not leave in thee one stone upon another, because thou knewest not the time of thy visitation.''

He \emph{would}, but they \emph{would not}, that was his lament. Even when he went to Golgotha and bore his cross, and the women wept over him, he forgot his own sufferings in the thought of what was in store for the unhappy people. ``Daughters of Jerusalem! weep not for me, but weep for yourselves and for your children.''

There are knots in the history of the peoples, turning points, of which it holds: now or never. If the opportune time is neglected, the good conditions misused, a wrong grasp taken, outrages committed: then generations must pay for it through centuries. No people on earth has fought with such a lion's courage, with such a magnificent strength of despair for its right and
% --- p. 33 ---
\setcounter{footnote}{0}%
\opage{33}independence as the Jewish people at the last. Too late, the Helper was gone: woe to the people that knows not the time of its visitation. More than fifteen hundred years, full of sufferings, disappointed hopes and bitter experiences, lie between that time and us. Now the time of visitation comes again; for the peoples are awakening.

The speaker of the people from Galilee will come forward once more. Believe me, he does not forget his altar and does not forsake his congregation because he shows himself once again, as he is humanly, on the stream of world history; he has strength to spare. Long enough have the self-appointed taskmasters of the peoples, spiritual as well as worldly, held him shut up as in a close air within official church walls. He wants out into the open, under the open sky; from the mount, among the lilies of the field and the birds out there, he will be seen by the multitude and heard by the people.

We cannot of ourselves begin the great work; our strength is too slight, our sight too dim, our heart too weak. But if He himself comes, the Mighty One in the power of the spirit, then his voice will be able to be heard by all. He knows what dwells innermost in man, he knows what sacrifice is, and what it is that is called courage. He has the word, he is indeed the Word himself, the Word of truth that goes to the heart, the Word of righteousness that strikes the consciences, the Word of brotherhood that asserts freedom and equality, the Word of love and reconciliation that smooths away all strife between small and great.

``Will ye now?'' that is his question to all the tribes of peoples; will ye? --- He pledges his crown of thorns that the work shall prosper. They can hear
% --- p. 34 ---
\setcounter{footnote}{0}%
\opage{34}him, each in his own tongue; they can all understand him --- for the people understands the highest; they can all follow him, each on his own way, in whatever land they may dwell, and whatever colors they may bear. The common element that he offers, the first thing needful, is precisely fitted to consume what is hostile and nourish what is peculiar.

Would that our people, our small people with the great memories under the heavy pressure, might receive him with open arms! We bear indeed his banner; the red cloth with the white cross, that is his banner of victory. Forward, onward, where the banner waves; forward, onward, for truth and right; forward, onward, where he himself leads, the King of the peoples, the Teacher and Leader of the peoples.

% --- p. 35 ---
\setcounter{footnote}{0}%
\lecture{Third Lecture,} {given in Aagaard.}
\opage{35}In order to know human life, its aptitudes and aim and the
way that leads to the aim, we must, as far as we are able,
investigate the laws of existence. We must make a
thorough acquaintance with the real world,
we must consult nature and history. It does not, however,
depend only on what it is we see, or how much
we get to see; it depends also, and chiefly,
on the eyes with which we see, the light
in which we see, the standpoint on which we place ourselves.

Here it is at once obvious that man
lives not merely in one world but in several. Outward
nature with its many creatures, formations and productions
lies spread out before our gaze; its summer
and autumn, its winter and spring speak to us in a language
of their own and call us to different work. But we
have also an inner nature, a world of feelings
and thoughts that we bear within ourselves, and this inward world,
how different it can be. We also make
the experience that it is chiefly the inward
world that spreads its glow over the outward.
When we are glad, then everything smiles in field and
% --- p. 36 ---
\setcounter{footnote}{0}%
\opage{36}forest; but when we are sorrowful, then the sun may shine as
much as it likes, it is as though all nature bore a
quiet grief. The inner eye may be very different,
and the explanation comes from within; the one sees
in nature a wondrous play of matter and forces,
a web of necessity woven of blind laws; the other sees
the traces of God's wisdom, a revelation of the Creator's
omnipotence.

Something similar is repeated when we consider
history. The course of the human race advances
independently of us, for each of us is only a vanishing
member of the race. But each single human being
has nevertheless his own course of life, his own little history; if
the many little histories were to be told at length and honestly,
how different they would be, for all that they have in common.
Here, then, it is also remarkable that world history
and our own history, the great and the
little history, so far as understanding is concerned, essentially correspond
to one another and show themselves to us in a common reflection of
a common light. It is not the events, the outward,
finite occurrences, of which we speak here; it is the interpretation,
the judgment, the apprehension of the whole, on which it
depends. He who cannot find God's providence in
his own life, how should he discover it in
world history?

We have many teachers of history, and we are grateful to them;
we have learned researchers who know how to distinguish
between true and false and to spread light over the connection
of events; we have thinkers who work
to find out the laws of history: we are very grateful to them.
Should it ever succeed in discovering the
fundamental law by which world history moves, just
% --- p. 37 ---
\setcounter{footnote}{0}%
\opage{37}as one has already discovered the fundamental law by which
the heavenly bodies move in the space of the heavens: then we should
surely at last come so far that we could calculate
what will happen in the future. Will it ever
succeed? Hitherto the judgments about the future, I
mean not merely about our own nearest future, but about the
historical future as a whole, differ widely. There are
those who expect a great age, a glorious and blessed
age, although they are indeed at variance among themselves as to
wherein the glory and the blessedness are to consist. Others
look upon the future with dark anxiety.: Whirl,
dissolution everywhere, breach of old customs, crumbling away
of the old pillars of society, how fearful,
how sad!

The difficulty in discovering the laws of history
lies in this, that they are not merely laws of necessity, as
in material nature, but also laws of will, laws of freedom.
Those who cannot catch sight of any laws other
than the so-called unchangeable laws of nature must
arrive at whatever explanation they please, their explanation
is always one-sided and false; and for the simple
reason that the laws of history are more than mere laws of nature:
they are laws of spirit.

If we are to think of a teacher of history who can
hold up the mirror and show us features of the face of the future,
it must be he who, in the fullness of spirit and of the Godhead,
is himself the turning-point of history; it cannot
be anyone other than Christ. But what is it
that Christ teaches us about history, especially about the historical
future? The visions appall us, we are seized with horror.
First the destruction of Jerusalem and then immediately after ---
the end of the world; ``Immediately after the tribulation of those days
% --- p. 38 ---
\setcounter{footnote}{0}%
\opage{38}shall the sun be darkened, and the moon shall not give her light,
and the stars shall fall from heaven, and the powers of heaven
shall be shaken. And then shall appear the sign of the Son of Man
in heaven; and then shall all the tribes of the earth wail.''
(Matt. 24:29--30).

Those who build their historical view on pure laws of nature
have an easy time freeing themselves from
the terror. When they wail, they wail only in honor
of ``the unchangeable laws,'' and for that reason lodge
the most definite objections against the whole prophecy.
Their proof is short and clear: ``It is foretold that
the world shall perish \emph{immediately after} the destruction of Jerusalem.
Nevertheless the world has since that time
now stood for eighteen hundred and five years;
therefore the prophecy is refuted. To this it must be added
that the whole view of heaven and earth is antiquated.
The earth is not the center about which everything
turns, the stars do not sit so loosely in the firmament
that they should fall down,'' etc.
So long as one follows exclusively the direction of natural science,
these objections will be irrefutable.
Woe to those who have taken upon themselves the superhuman
task of defending the word of revelation about the signs in
sun and moon with scientific reasons: their future
is dark, they are in a sorry plight.

Otherwise, when we place ourselves on the standpoint
of personality and ask about the laws of spirit. The earth is then
the center, in so far as man, the being in whom
spirit dwells, is the center. Christ, ``the light of the world,''
is the center of history. What he has accomplished on
earth would not be heightened by having been carried out
on a greater star, were it even the common center of
% --- p. 39 ---
\setcounter{footnote}{0}%
\opage{39}gravity about which the solar systems revolve. Everyone who
has spirit will understand how false it is to apply
a measure of outward greatness in space when
the question is of spiritual greatness. For Christ's prophetic
vision the whole historical future dwindles to a few
days; for he measures with the measure of eternity. What
is a day compared with a hundred thousand years, and what is
even a hundred thousand years against eternity? Therefore
Christ says: ``Immediately after those days.'' But how,
one asks, can the sun be darkened, and the stars fall
from heaven? --- It must surely happen when the powers of heaven
are shaken, and the sign of the Son of Man appears
in heaven; it must surely happen, as surely as it is spirit
that bears nature.

Natural science is good and its guidance excellent,
when we take it for what it is, an elucidation of the
real, an illumination of natural causes, a
pointing out of the given laws. But when it is employed
to prove that ``the given laws'' are not \emph{given}, because
there is no legislation and no lawgiver: then
this is a scientific outrage and an affront to
spirit, for spirit is the original. For spirit
itself, the almighty unity of wisdom and will, temporal
nature in all its shapes, in all its duration, with all its greatness
is a vanishing quantity.
The power of spirit in matter is eternally the same, whether
the mass be only a speck or a thousand times greater than
the sun, for the power of spirit is omnipotence.

Christ, the revealer of spirit, of the Almighty,
knows within himself that the innermost mover in him
is that which moves heaven and earth. From this
height his prophetic vision of the future must be apprehended.
% --- p. 40 ---
\setcounter{footnote}{0}%
\opage{40}The unclear mixture of true and false, of good and
evil, which confuses us so long as finitude with its
many conditions stands in the foreground, disappears when
everything is led back to the first source. History
shows itself here in a super-historical light. It is not enough that
nation rises against nation and kingdom against kingdom, that there
is hunger and pestilence and earthquakes here and there;
the innermost thing in the ``birth-pangs'' is that spirit rises against
spirit, that the cutting oppositions lay themselves bare and collide
without mediating middle terms: the powers of truth and
falsehood, of light and darkness, of life and death
move here in the last conflict.

However hard the conflict is, its outcome is
nevertheless not doubtful. The victory is at bottom already
won. Christ sees himself in the vision as the victor,
although his suffering and death still lie ahead. ``I
saw,'' he says ``at another place, Satan fall
from heaven, as lightning''; the power of evil is in the eternal sense
a powerlessness, however much power it may have in
the temporal. It is the same thought that in John's
Revelation (12:7--12) is presented under the following
image: ``And there was war in heaven; Michael and
his angels fought against the dragon, and the dragon fought, and
his angels. But they prevailed not; neither was
their place found any more in heaven\dots{}. Therefore
rejoice, ye heavens, and ye that dwell in them; Woe to them
that inhabit the earth and the sea; for the devil is come down
unto you! He hath great wrath, because he knoweth that
he hath but a short time.'' The historical struggle between
the good and the evil is thus to be apprehended as the
last scenes of a great battle in which the enemy has lost,
but still fights the fight of despair. The
% --- p. 41 ---
\setcounter{footnote}{0}%
\opage{41}continuing history is the long day of battle, which stretches
out in time, so that the liberation may be completed in reality,
so that as many as will may take part in the fight
and become partakers in the victory.

But --- here a new objection arises --- if
the whole historical future is only to denote the continued
destruction of the world, what then becomes of our expectation
of the age of the rights of man, the new
age of the people and the peoples, from which we promise ourselves so much?
If Christ is to be our teacher of history, then it is probably
best to let that sort of fancy go and to join
those who clearly see that everything is getting worse and worse.
If the good can only conquer by the world's going under,
then we do best to work along toward its destruction.


What, then, is the meaning of all this talk about the popular,
the new, the great, the hitherto unseen; unless
one is to understand the matter in this way, that the awakening of the people,
popular life and the spirits of the people are to be the most powerful
means of hastening the destruction of the world? --- The
cold mockery betrays a hidden despair. On the way of the future
there are many stations. ``Learn a parable
of the fig tree: when the sap is come into its branches,
and the leaves shoot forth, then ye know that summer is nigh.''
This parable applies not only to the very last
breakthrough, it applies also to the many intermediate ones.
Many summers may come, each with a
harvest of its time, before we reach the summer of eternity.

The law of development is now this: A confusion of the age
arises every time an existing order dissolves, and
something new is to come forth. The old falls into decay; the new,
which is in the process of forming itself, has not yet
% --- p. 42 ---
\setcounter{footnote}{0}%
\opage{42}won shape: out of fermentations, disturbances and wild
disorder there then clears itself, no doubt under violent
exertions, little by little a new and better order
of things. Society comes into equilibrium, distributes its
labors, sows and plants and prepares for a rich
harvest. But the more abundantly one enjoys all temporal gifts,
the more commonly one dreams of eternal peace and unbounded
well-being, the more keenly is the surprise felt
when the upheaval that a creeping corruption
has long prepared suddenly sets in. Out of
new fermentations and new struggles there must then form itself a new and
better condition, always on other and higher terms;
but every condition on earth, even the most perfect,
is transitory. How intimately this thought is connected
with the main Christian thought can be expressly
shown.

In John's Revelation (20:1 --- 3) we meet the following
vision. ``And I saw an angel come down from heaven;
he had the key of the abyss and a great chain
in his hand. And he laid hold on the dragon, that old serpent,
which is the devil and Satan, and bound him a thousand
years, and cast him into the abyss, and shut and sealed it
over him, that he should deceive the peoples no more,
till the thousand years should be fulfilled; and after that
he must be loosed a little season.'' This is the thousand-year
kingdom whose coming is announced; but what
does the thousand-year kingdom mean? We dare not confound
the super-historical with the historical and give ourselves
up to enthusiasms. How the expectation of
a thousand-year kingdom is to be understood in the congregation is a
matter apart: here we stand on popular ground and have
simply and plainly to do with a law of spirit. Every time a
% --- p. 43 ---
\setcounter{footnote}{0}%
\opage{43}new and better order of things triumphs, the spirit of darkness
is literally subdued; the evil will is chained, forced
down into the abyss, and must keep itself hidden until a time
when it again comes loose: that is the general thought.
That even the thousand-year kingdom, however glorious and mighty
it is said to be by reason of ``the first resurrection,''
is nevertheless beset by strong enemies and exposed
to great dangers, is depicted in vivid colors (Rev.
of John 20:7--8). ``And when the thousand years are expired,
Satan shall be loosed out of his prison. And shall go out
to deceive the peoples which are in the four quarters of the earth, Gog and
Magog, to gather them together to battle; the number of whom is as
the sand of the sea.''

That is an important lesson we can draw from this.
We dare not deify the popular or set up a kingdom of a people
in place of the kingdom of heaven, in which Christ's
revelation is perfected. Nevertheless the Millennium
may be to us an image of this, that Christ's spirit
will give itself an imprint also in earthly conditions: the Savior
of souls is also the Savior of the peoples. One may now place
the popular as far below the Millennial as
one will, yet this must stand firm: if truth
and right are to gain the upper hand over lies and wrong, then it
must happen by virtue of the spirit of righteousness and truth.
If the temporal is too slight and the finite too
wretched to furnish matter for the spirit striving upward,
what then has the future to offer? Slackness of will and
emptiness of mind, enfeebled nerves and burnt-out brains,
raw masses and unused powers, surfeit and yet
hunger, plans and yet perplexity, loathing of
everything and yet greedy longing for enjoyment, all this in a
jumble, stirred together into the strangest
% --- p. 44 ---
\setcounter{footnote}{0}%
\opage{44}mixtures and sediments, shall then not give human life
a higher worth, shall by no means weaken the
simple faith that it is Christ's spirit that must come
with the signs of spring, that summer is near only when
the fig tree gets sap in branches and leaves.

The historical driving force can be measured by the tension
between the eternal and the temporal. Much is ventured
when life is at stake, most of all, however, when there
is a strong feeling that it is eternal life that is at stake.
The spiritual cannot make an impression on human beings,
least of all on a whole people, unless in some
way or other there is something sensuous with it. Christ's
exhortations are clothed in parables; they are nothing but
exciting scenes, decisive moments, that glow in
spirit with sensuous colors. ``Watch therefore, for ye
know not what hour your Lord doth come. But
know this, that if the goodman of the house had known in
what watch the thief would come, he would have watched
and not have suffered his house to be broken up. Therefore be ye ready,
for the Son of man cometh at an hour when ye
think not.'' It is no idle waiting; all hands
are kept busy: here is something to do. It is as
when ``a man traveling into a far country called his servants,
and delivered unto them his goods; and unto one he gave five
talents, to another two, and to another one, to every man according to
his ability. Then he that had received the five
talents went and traded with the same, and made other
five talents. And likewise he that had received
the two talents, he also gained other two. But he
that had received the one went and digged in
the earth, and hid his lord's money.'' How hard and
severe is the judgment upon this last servant. He
% --- p. 45 ---
\setcounter{footnote}{0}%
\opage{45}steps boldly forward, as though he had a just
cause: ``Lord, I knew thee that thou art an hard man,
reaping where thou hast not sown, and gathering where
thou hast not strawed. And I was afraid, and went
and hid thy talent in the earth: lo, there thou hast that
is thine.'' Alas, it goes ill with him.

The temporal must be used, for the gifts of spirit must
be used. The spiritual does not exist for the sake of the temporal;
it is the other way around, the temporal exists for the sake
of the spiritual. The mark that the temporal
is used in the right way is that the powers of spirit
grow by the use. There are in all temporal conditions of life,
be they private or public, small or great, conditions
for a development of the life of spirit: hardy
men with weather-beaten faces and horny fists can
quite as well have spirit as those who move in the
perfumed atmosphere of fine cultivation. Especially must
it be emphasized, what the parable inculcates, that no one
may excuse himself by saying that his endowment is only slight.
The servant who got the one talent got the same obligation
to work with the one as he who got
the five to work with the five. What holds
of the single human being holds of a whole people:
a gifted people has a great obligation. Not to
use the gifts of spirit is to sin against spirit; but
sin against spirit is punished doubly. The first punishment
is that the gifts are taken away, that the person concerned becomes
spiritless. ``Take therefore the talent from him, and give
it unto him which hath ten talents. For unto every one that
hath shall be given, and he shall have abundance:
but from him that hath not shall be taken away even that
which he hath.'' The next punishment is that the spiritless man
% --- p. 46 ---
\setcounter{footnote}{0}%
\opage{46}is condemned as a criminal against spirit. ``Cast ye the unprofitable
servant into outer darkness: there shall be weeping
and gnashing of teeth.''

That the expectation of the end of the world can in
history coincide with magnificent transformations
of the conditions of the world and corresponding breakthroughs of
the life of spirit rests precisely on this, that it is linked
with the expectation of the return of ``the Son of Man.''
The spiritual law expressed in this shows, namely,
that as the temporal exists only for the sake of the eternal,
outward nature only for the sake of spirit, the world only
for the sake of the Son of Man, so it is also
both here and hereafter human life and its perfection
about which the whole turns. How the super-historical
in the progressive movements intervenes
in the historical can be seen by casting a glance
at the stations already traversed. Three points of time:
the age of the apostles, the age of the Crusades and the age of the Reformation
present themselves for consideration.

In the age of the apostles the opposition between the kingdom of spirit
and the world-empires, spirit-denying yet so mighty,
stood at its height. A few fishermen from Galilee,
defenseless men and simple women, that was the whole
equipment; those were the men of arms with which
the new spirit was to take up the struggle and complete the
work of him who had died as a malefactor on the cross. Dark
--- and yet bright prospects! The great Departed One had
whispered something to them in the chamber; that they were to preach
upon the housetops. He had not remained in the grave and
forgotten the Forsaken; he had visited them again,
he had greeted them as the conqueror of death, and promised
them that he would soon come --- well equipped. He
% --- p. 47 ---
\setcounter{footnote}{0}%
\opage{47}had not disappointed them with a vain hope, as
if it were at once a feast of victory. He knew the world, the
hatred of the mighty and the blindness of the multitude. ``They shall
deliver you up to be afflicted, and shall kill you;
and ye shall be hated of all nations for my name's sake.
And then shall many be offended, and shall betray one another
and shall hate one another. And many false prophets shall
rise, and shall deceive many. And because iniquity
shall abound, the love of many shall wax
cold. But he that shall endure unto
the end, the same shall be saved.'' The first prediction
has found abundant, literal fulfilment, a pledge
that the last also will constantly be fulfilled.

``Maran Atha, the Lord cometh, he standeth at
the door and knocketh!'' These omens, which were heard incessantly,
had an electrical effect on sinking courage and awakened the
Combatants, each time they began to fall asleep. It
was decided in heaven that ``the beast with the seven heads,''
the world-city on the seven hills, the rich
city of the Roman empire, ``Babylon, the great whore, a prison for every
unclean spirit and a cage of every unclean and hateful bird,''
should fall. The description of its fall is vivid
and gripping. ``A mighty angel from heaven took up
a stone like a great millstone and cast it into the
sea, saying: Thus shall Babylon, the great city, be thrown down
with haste, and shall be found no more. And the voice of harpers
and musicians and pipers and trumpeters shall
be heard no more in thee; and no craftsman of any
craft shall be found any more in thee; and the sound of a millstone shall not be heard in
thee any more. And the light of a candle shall not shine in thee
any more, and the voice of the bridegroom and of the bride shall not be heard in thee
% --- p. 48 ---
\setcounter{footnote}{0}%
\opage{48}any more; for thy merchants were the princes of the earth,
for by thy sorcery were all the peoples deceived.''

History has here furnished its most magnificent proof
of what the spirit of truth with spiritual weapons
can accomplish in temporal strife, when the powers of heaven are shaken,
and ``the sign of the Son of Man'' shall appear in heaven.''

In the Middle Ages the intermingling of the
heavenly and the earthly was so peculiar, and the separation
of the godly and the worldly at the same time so
outward, that spirit and spiritlessness almost flowed into
one another. The unrest and lawlessness that set in with
the dismemberment of Charlemagne's empire; the Mohammedans'
attacks, the Hungarians' inroads and especially the Normans'
plunderings, all this in union with exalted notions
of monastic holiness and monkish abstinence stamped
the whole earthly existence with a dark mark of
frailty. People believed themselves to be in ``the last
birth-pangs.'' Plague, famine, eclipses and comets came to be
regarded as signs of ``the loosing of the dragon'' and ``the coming of
Antichrist.'' All pious souls prepared themselves
for the end of the world, which was fixed for the year 1000.
When the year passed and the world remained standing, the date was moved
forward to 1030, and the signs became more manifold still.
A monk in Hungary saw a dragon, lifelike
and shaped like Leviathan, advancing with lightning and thunder;
a French bishop caused it to be proclaimed that he had
received a letter from the third heaven, in which it was commanded among
other things, ``not to partake of anything but water and bread
every Friday;'' etc.

But however wild and motley it looks in those
half-barbarous conditions, spirit is nevertheless present in the
mighty pressure; the eternal weighs upon the consciences
% --- p. 49 ---
\setcounter{footnote}{0}%
\opage{49}and portends a great upheaval and a day of judgment.
The heavy pressure conditions a corresponding upswing. The
enigmatic signs at last received a historical interpretation.
In the Pope's word about the holy sepulchre the aim was made clear.
The sepulchre of the resurrection! now the watchword was given, and a rushing
flood of released powers found an outlet. The main idea of the
Crusades is a popular idea of faith; it
contains all that the Middle Ages possess of the peculiar
and the great. When the many repellent features of blind
fanaticism, wild cruelty, cunning self-interest and raw
sensuality have aroused our displeasure, we are then drawn
again by that inwardness of piety, the mighty enthusiasm,
the mountain-moving faith which, despite the admixtures
of superstition, is yet drawn from Golgotha's spring and tastes
of the source. The sign of the Son of Man is stained by crimes,
but the human race makes giant strides forward.

The age of the Reformation denotes on the whole the dissolution
of the Middle Ages, and yet in a way its continuation.
The Reformers too had a strong
impression of the vanity of the world, and the possibility that
the last day might be near at hand stood
vividly before them; but it was chiefly Christ's
spiritual return for which they hoped. The Holy
Scripture was the mirror in which his figure could be seen
distinctly; so they drew forth the Scriptures, and the Bible became a
book of the people. How this has since been
misunderstood by many, as if the meaning were that one really
could read Christ into oneself by reading in the Bible, shall
not be treated here. What we have before our eyes are the conditions
under which, and the tension with which, the
half-extinguished life of spirit broke through anew. Christ
did not come down from the clouds, but the sign of the Son of Man
% --- p. 50 ---
\setcounter{footnote}{0}%
\opage{50}began to reveal itself in the inner being of men. It
was a living inwardness that was set against a dead outwardness.
The outwardness of the papal church was rejected and its worldliness
condemned. Rome was again ``Babylon,'' the great
whore, who by her sorcery deceived the peoples. It
is the struggle, a \emph{spiritual} one, yet through a multiplicity
of worldly relations a \emph{popular} struggle between
the old and the new Church, that gave the Reformation
its tension: Luther was on a grand scale a man of the people,
a popular man.

To an outward view the whole
historical movement now looks as if Christ's return
were to take place not from above, from heaven, but much rather
from below, from the spiritual depth in which life has its ground.
From this point of view those who regard the laws of history
as mere laws of nature can explain the whole to themselves
in their own way. ``Spirit,'' they say, ``is nothing other than the
human nature's own inner, universal being.
It was this pure human being that, in the apostles'
imagination, had taken flesh and blood in the man
Jesus Christ. Since now spirit as spirit is at the same time
the divine, those disciples, filled by the Master's life
and teaching, believed that Christ as the Son of God
was himself both God and man. The kingdom he
had founded by taking upon himself the form of a servant
he would then, after his exaltation, complete and lead to
victory by coming as the Lord in the clouds of heaven with
power and great glory. But he did not come; something
else happened. Instead of the expected Jerusalem,
the holy city, which was to descend from heaven,
there arose an earthly, and yet heavenly, Jerusalem upon
the ruins of Babylon, the old whore: the Catholic
% --- p. 51 ---
\setcounter{footnote}{0}%
\opage{51}Church was, in other words, the realization of the expected
return. Yet the spirit of man did not
yet understand itself. The earth trembled under the violent
shocks of a society torn asunder; people stared up
into the clouds, and believed firmly and fully, what all the signs pointed
to, that now heaven must open, and the Judge
come with the angels and all the saints. But he
did not come; something else happened. Weary of waiting
for the heavenly Jerusalem, people set about
conquering the earthly. It was a conversion of
imagination into reality. Christ's return transformed
itself into the rise of the papacy. Life with its stirring
powers came from within; enthusiasm kindled, it was
like a sea of flame of sensuous spirit and holy passions.
Only with the Reformation did the spirit of man begin
to surmise how it really stands with the
return. The Son of Man is man himself,
and man's self, the true man, righteous before God,
is in the inwardness of the heart, in faith. But the faith of the
Reformers too was medievally afflicted with the notion
of a kingdom of God that stood separate from worldly
life. Therefore the Reformation, with all the
awakening of the people and purification of morals that it unmistakably
called forth, could not carry the matter through. Its
official church organization went under the state; the state
made the holy a means, and the thrones by the grace of God
pressed the peoples down into masses. Only when
the masses rose and burst the absolute, self-deifying
state did the spirit come to consciousness of itself,
and the supernatural return received a natural
explanation. The rights of man: freedom, equality and
brotherhood, that is the sign of humanity. This sign
% --- p. 52 ---
\setcounter{footnote}{0}%
\opage{52}is not among angels in the airy clouds of heaven, it is in
clouds of dust in the midst of men down here on earth.
What the sign signifies, the future shall show, when the powers
stir and the fight is waged for the holy cause
of humanity with real, human weapons.''

This explanation has something appealing in the eyes of the present age,
because, though exalted, it yet hovers so low
and keeps so close to the real. There is something
true in it, but the truth is only half; when a
half-truth is given out for a whole, then we have the lie.
It is quite true that the way of spirit is the way of inwardness,
that the powerful inner being must burst forth and give itself
form in a real, present outwardness. But it is
not true --- it is in a lamentable manner untrue --- that
the great thoughts of freedom, equality and brotherhood,
these thoughts of eternity, which embrace all the temporal
tasks of humanity, should be only selfish
thoughts of nature, conceived in bitterness and born in rebellion;
the so-called new thoughts are old thoughts: they are
distorted thoughts of Christ. It was not for nothing that
Christ came forward as the friend of the poor, and took
up the cause of the lowly. If the new thoughts with the new demands
could really be carried through by a denial of
Christ's spirit, then the revolutionaries of 1789 and
1871 would be right: the bloodier, the better. But history
appears as a witness, and its testimony is: that it is not
merely tyranny, but also the freedom of the people, which under % defect: Folketriheden (a broken f)
the red banner of the rights of man can be smothered in blood.

With all our ``modern knowledge'' we are by no means
on the straight road of humanity, freedom and progress;
we are on the contrary about to turn
onto a byway. Our faith-denying condition is a
% --- p. 53 ---
\setcounter{footnote}{0}%
\opage{53}striking counterpart to the superstition of the Middle Ages. The Middle Ages
imagined with the super-historical, the present age
would help itself with the merely natural-historical. The truth
is that the super-historical must bear the historical, that
spirit from heaven must be present in all vigorous deeds
on earth. That is the law of movements; that is
the foundation of popular life; that is, in sum, what
Christ teaches us about the historical future.

% --- p. 54 ---
\setcounter{footnote}{0}%
\lecture{Fourth Lecture.} {(Fredericia).}
\opage{54}``Evil and good in every land, with thorns the roses mingled stand'', that is an old ditty; it applies
not only to the many lands but also to the many ages;
our own present age has roses, but also
thorns; it has strong sides of light, but also troubling
sides of shadow.

Our age is the age of enlightenment; undeniably! The night
itself bears witness to it. Our streets and houses are
so brilliantly lit that we cannot help
feeling a certain pride when we think how dark it
looked in ``King Hans's day''. And the many lights of the
understanding burn just as brightly. How much has
not already been done for improved instruction of the young;
to what a degree have the sense and taste for knowledge
and cultivation not increased in the last two generations.
The horizons widen; attention is
no longer confined and held captive by the nearest
everyday concerns; our thoughts wander far afield;
foreign regions, foreign languages and foreign customs
invite comparison with our own. We understand
that what goes on in distant lands can concern
% --- p. 55 ---
\setcounter{footnote}{0}%
\opage{55}us: we feel the pulse-beat of the great whole, and the events of the world
are outlined as clearly as
never before.

But the shadow follows the light. The manifold,
which distracts the mind, can so easily crowd out the first
simplicity; and yet it is precisely the simple that
gives concentration, inwardness of strength, and resolute will.
Instead of whole thoughts, in which the soul can rest, culture
hatches a multitude of briskly playing but
dissimilar notions, a swarm of fractions of thoughts,
for which one searches in vain for a common denominator. To
know a little about everything is not the same as to know
something thoroughly. The man who must be in on everything, so as
not to be one-sided, is seized by a restless striving; and the
inquisitive man who absolutely must have wind of everything and an
opinion about everything is seldom in a position to pass a
thorough judgment on anything. The worst of it is that it tells
upon the character. The man never acquires
character who is not in a position to form himself a
firm conviction. In an age in which a given
view of life prevails with authoritative uniformity equally
for all, it certainly looks as if every
individual in society had been born with conviction; but
under advancing enlightenment the uniformity begins
to dissolve. Many ways of regarding life,
many and conflicting points of view present themselves to the eye,
and what is the consequence? There is, one says, something in
the one; but there is also something in the other: evil
and good in all doctrines! All that one gets out of it
is that much can be said for and against every thing.
In the end everything becomes uncertain; one must go by impression
and mood. But impressions change and 
% --- p. 56 ---
\setcounter{footnote}{0}%
\opage{56}moods vary: everything depends on circumstances.
He who ten years ago was an arch-liberal is now gravely
gloomy and a die-hard conservative: it is circumstances.
The fir tree, which was once the Christmas tree, now stands in the loft,
and has been ``set to growing'': that too is circumstances.
How uncanny to meet one's first original view
on the way of development and then not to be able to
acknowledge it again; it is almost as
vexing as having to meet one's ``first love''
and to see that she has become engaged to another. To
doubt everything, to criticize everything and to be dissatisfied with
everything is often regarded as an infallible sign of
superiority; and yet this superiority is just as
deceptive as the flush on the cheeks of the consumptive.
Endless doubting is a wasting disease, and in the age of enlightenment
the disease is contagious.

Our age is an age of freedom; we acknowledge it
with gratitude and say it with pride. But
the age of freedom is also the age of dangers, for in freedom
there are many kinds of possibilities: to deal with freedom
is to deal with possibilities, but to deal
with possibilities is like dealing with fire and light in
a powder magazine.

In older times, indeed far into the eighteenth
century, what is now called freedom was a forbidden
fruit. Popular freedom was under lock and key,
its thoughts lay hidden in the book of secrets, a
book with seven seals. Woe to him who dared to break one
of these seals, his doom was sealed. Even with religion
and science, faith and reason, the power of the state kept strict
watch; it brooded over the jewels of the spirit like a
dragon over its treasures. When we Protestants cross
% --- p. 57 ---
\setcounter{footnote}{0}%
\opage{57}ourselves at hearing of all that the Catholic Church has
managed to accomplish by means of censorship, the Inquisition, witch trials
and heretic pyres, then we must above all things not
forget to include our own Church. In this respect it is
particularly instructive to become acquainted with
the history of England in the time before and after the Reformation.
After the persecutions with their bloody punishments had raged themselves out,
a drowsy order in the High Church and a mystical
enthusiasm among the sects became the natural condition. To
illustrate this the painter Hogarth has furnished
a picture in which the satire surpasses art. In the foreground
one sees the sectarian spirit going about its mischief. A
woman in garments full of folds lies flung down on the floor,
and out from the folds there swarms a multitude of small animals,
mice, as far as I remember. The groups around her
are full of life, some in quiet rapture, others in
such consternation that their hair stands like bristles
on their heads. It is a miracle that is taking place;
the favored woman gives birth to mice. So: a great
miracle of mice! The background is devoted to the High Church. The congregation
is assembled, the priest in the pulpit, here is
order and quiet. As commanded, the shepherd of souls reads his
sermon aloud from the paper. A strip of it informs
us what the text is whose content
he lays upon the hearts of the devout. ``You shall not
think that I am come to bring peace on earth;
I am not come to bring peace but a sword.'' The ``pure doctrine''
of Orthodoxy is, however, so reassuring
that not even the word of fire and sword can disturb its
peace: the priest looks sleepy, and the whole congregation
is already sweetly slumbering.

That whose essence is freedom cannot be held under
% --- p. 58 ---
\setcounter{footnote}{0}%
\opage{58}constraint; it is an unnature that takes its revenge. For when
freedom is spoken of, it is not enough to think
of freedom in the outward sense, such as freedom of trade,
freedom of the press, a free constitution and so forth. Outward
freedom is an inestimable good; but it
receives its worth only by furnishing conditions for the
inward. Ignorance and superstition are incompatible with
inward freedom. Science is by its essence free,
for science wills to know the truth, and the knowledge of truth
is in an inwardness. To will to prescribe to thought by outward
commands what it shall recognize
as true is the same as to will to command that two and
three shall not be five, but four. Religion is
by its essence free, for religion is in the highest sense
an inward concern. In order rightly to use
outward freedom, we must acquire inward freedom.
The mind that is seized either by covetousness or by
lust to rule or by any other disease
is unfree. There are many who have made a great noise about
liberality, freedom of mind, and yet, incredibly, have somewhat suffered from
lust to rule: they have then not granted themselves the good
that they so gladly wished to secure for others.

Our age is in a special sense an age of debates;
that goes with enlightenment and freedom. Much of
what in former days one scarcely dared think to oneself
or whisper into a confidant's ear now comes openly
forward in writing and speech and is discussed before the eyes of all the world.
If we could now trust that the many voices,
the many tones and the many opinions that let themselves
be heard in the continually prolonged conversation with which the
civilized world daily entertains itself really
stood in the service of truth, then we would have rich
% --- p. 59 ---
\setcounter{footnote}{0}%
\opage{59}means of getting the burning questions illumined
from all sides and of getting to the bottom of the real conditions.
But what is the procedure on the whole?

When one wishes to test a new machine, to appraise its
building material, its composition and working power,
one does not content oneself with listening to the eulogies with which
the inventor himself recommends it; just as little does one
let oneself be deterred by the objections of those who manifestly
have an interest in running down and censuring the machine. One
proceeds carefully, takes experience and the experts
into counsel, notes exactly what speaks for
and against, and lets oneself be convinced by reasons. It is
not moods or hasty, preconceived opinions
that are here allowed to decide the issue: no,
it is reasons, weighty reasons.

As regards the debates on the great, intricate
machinery of society, it seems on the contrary as
if the enlightened and free nations of the present preferred to strike
out on another road. If we apply the saying about
evil and good in every land to the now so rich newspaper
literature that fertilizes all fields, irrigates all meadows, and shines upon all roads,
then it would have to run: Evil and
good in every paper! But no! When we see the parties
group themselves to the Right and Left, extreme Right, extreme
Left, right center, left center and so on, then
it looks to the life as if the old saying
suddenly changed and meant: pure good in some
papers, pure evil in others. Those who have sworn themselves
to the Right assure us that there are no roses at all in the Left papers,
but only thorns; in return those who
have sworn themselves to the Left must naturally make known,
and that both with pictures and in bold type, that the 
% --- p. 60 ---
\setcounter{footnote}{0}%
\opage{60}Right papers are withered thorns, and that the genuine, true, unadulterated
fragrance of roses is to be sought in the Left papers alone.
So a conviction and an
opinion do indeed form; but it is of a negative kind. The conviction
is that one has an opinion, and the opinion is that one
will not let oneself be convinced.

The ship that carries the great cargo of the debates of the present
does not draw deep enough; it does not plough in
the open sea, but prefers to keep to where there are shallows.
We see it when we cast a glance at the debates
on religion and science. That no people can live
without positive religion, that thrones fall and states totter
when the holy pillars are taken away, on this statecraft
is so much at one with itself that persons who,
for their own part, have manifestly given up all faith
in the doctrines of the Church are often the most zealous defenders
of the existing ecclesiastical order. Religion is,
in short, publicly recognized for what it is, a spiritual
great power. The right and significance of science must,
especially in an age of enlightenment and freedom, be exalted
above every doubt: science too is then recognized
for what it is, a spiritual great power. If
now these two great powers could unite to
work together toward a common goal, then slumbering
forces would awake and a new day would break over peoples
and states; but it is precisely the torment of the present that between
religion and science a chasm has opened
as deep and wide as never before.

It is taught by the Bible and the Church that the world was created in
six days; but science proves that between
the coming into being of the first animals and the coming into being of the
first human beings there must lie millions and millions of years.
% --- p. 61 ---
\setcounter{footnote}{0}%
\opage{61}It is taught in the Bible and the Church that sun, moon and stars
were created for the sake of the earth and of human beings; but
science proves that the earth is a vanishing
point in the infinite space of heaven, and so forth. Were
it only deviations in the conception of the outward,
of relations in time and space, it might pass; but the strife
concerns the innermost kernel of existence, the divine
power and will. Religion rests upon the miracle;
science cannot acknowledge the miracle; how
should two spiritual powers so much at variance with one another
be able to work together: what the one builds up, the
other must pull down.

In the old days the strife was a secret; the Church
took care that no scientific deviation from the
orthodox doctrine came to the knowledge of the people, or that
at any rate, if a bold thinker spoke too loudly,
a deterrent example was made. % defect: aískrækkende (a broken f)
In our time, the age of enlightenment,
freedom and publicity, the secret
cannot be kept: every other person knows
that natural science does not conform to what stands in
the Bible. None the less the public debates move
along the surface. Commissions are appointed
to investigate the rights of the Church and of the state;
but it does not occur to the enlightened present to arrange
a serious inquiry into how things actually
stand with the relation between religion and science.
In the midst of an age of fermentation, of doubts, of criticism
and of restless thoughts, one settles down to
an instruction in which the biblical and the scientific
accuse one another in the most naive way.
One holds church for those who will believe, and
science for those who will know; but a smoothed-over
% --- p. 62 ---
\setcounter{footnote}{0}%
\opage{62}hodgepodge of faith and knowledge is not for the edification
of the people.

``Such settlements,'' one objects, ``cannot be
forced.'' --- There is here no question of force at all; what
is here demanded is an accounting, an account of how
it stands with the great moving powers of the spirit.
Should such an account be either dangerous or
impossible, then it must stand badly with the conscience of
society. The only objection that could with reason
be made would then be this, that such a debate
would become so learned and protracted that the
great multitude could not follow. We answer:
what essentially concerns all must be capable of being understood
by all; --- the people can understand the highest.
What is here aimed at is no declaration of authority,
either from the side of those who represent the doctrine of the Church, or
from that of those who represent science, but simply
and plainly a pointing out of the boundary by which religion
and science necessarily must be parted and meet.

It is chiefly three points on which the inquiry
turns: \emph{the origin of things}, \emph{the nobility
of man} and \emph{the verdict of conscience}.

Concerning the origin of things religion teaches that the world
is \emph{created,} that it is, in other words, a work of the Almighty's
will. If science could really
prove that the world is not created, then one would have to see, even
if by gentle means, little by little, to getting
religion abolished. But what does natural science now know
really about the origin of things? Nothing. The science of experience
takes things as they are given; by investigating
their forces, connection and laws it discovers that
the one state has preceded the other. The
% --- p. 63 ---
\setcounter{footnote}{0}%
\opage{63}first state, the outermost boundary to which it can
reach, is, as is well known, an enormous ``fog''. If the science of spirit
(philosophy) is to give an account of the origin of the fog,
it must necessarily presuppose that an active
reason, a divine wisdom, an almighty
shaping will has been present in the fog. But
from this standpoint it is impossible to prove that the
world is not created. Everything in nature develops from
the very beginning according to strict laws. Can the Creator
change his laws? The miracle stands as an indelible
boundary marker between religion and science; but the
alleged fundamental contradiction is removed. Religion is
right on its side when it asserts the miracle; for the
innermost original of everything is the freedom of the
lawgiving will. Science is right on its side when it
maintains that the laws cannot be broken. But as
surely as the laws of nature are the very own
revelations of lawgiving freedom, so surely are they not broken by being taken up into
the connection of the laws of spirit and made subject to the free
omnipotence.

The nobility and high vocation of man
religion expresses briefly and pithily in the saying: Man is
created in the image of God. What has science to
object to this? Natural history regards --- and from its
standpoint with full right --- man as a link in
the series of animals. Of course man occupies the
highest place, the place of honor above apes and bats.


Great preparations are made to prove that, just as
the higher forms of animals have little by little developed
out of the lower, so man has developed
out of a particular animal form, probably an ape-like primal form.
% --- p. 64 ---
\setcounter{footnote}{0}%
\opage{64}Darwin emphasizes that apes have a predisposition to
consumption, just like man; that they become angry when
one makes fun of them, exactly like man; that
they can learn to smoke tobacco, like man, and so on.

The partiality with which many natural historians
are at present working to efface the last traces of
\emph{original} difference between animal and man is so
warm that one can only cool it somewhat by directing
their attention to the relation between
science and man. Science has weighed
the globes, science has measured the speed of light,
science has discovered stars that sit so high
that the light they send down to us must be millennia
on its way, and that although light covers a distance of
over 40,000 miles in every second. These are high things.
When the natural historian now considers that science is
human, is a spiritual product of the human
brain, he may even pass over to the opposite
extreme and be carried away into such an enthusiasm that
he sets the human brain in the place of God and deifies
the human spirit. What a distance between
the animal and man! How many millions, billions, trillions of
years must not pass before the ape's
brain becomes so developed that it can weigh the globes,
before the bats manage to measure the speed of
light.

The whole alarm about man's animal descent
arises quite simply from this, that a famous
natural scientist, in a commendable zeal to demonstrate
the connection of the generations and the gradual transitions,
has forgotten that there is something in nature called
\emph{endowment.} It is by virtue of his divine endowment that
% --- p. 65 ---
\setcounter{footnote}{0}%
\opage{65}man appears as an image of God on
earth.

Concerning the verdict of conscience religion teaches that it
is an immediate and natural declaration of the
all-knowing, holy and righteous will of God in the inward
parts of man. Should science raise an objection to this
doctrine, it would first have to prove that an all-knowing,
holy and righteous God cannot possibly exist; but tasks
of this nature lie outside the domain of science
or, what says the same, outside
its reach. Conscience is not like
an animal instinct, which is at once finished from
the hand of nature and always remains on the same level.
Conscience is awakened and developed on the basis
of the religious, moral insights that society offers,
and is conditioned on the whole by the relation between religion
and science. Real life with its many labors
and duties must, so to speak, be illumined by two suns:
as to what concerns the inward and personal, by religion; as to what
concerns the laws of the world and the real conditions, by science.
Conscience can make a man lonely;
but it can also hallow and strengthen the bonds that tie
men to one another: there is a public conscience just as surely
as a private.

That conscience can hardly be counted among the
saints who have so far found special worship in
politics may perhaps be regarded as generally conceded.
It is remarkable to see how the states everywhere,
where their ``honor'' and ``interests'' demanded it, have
known how to preserve complete independence in face
of the moral obligations that State-Church
orthodoxy prescribes. From their mutual 
% --- p. 66 ---
\setcounter{footnote}{0}%
\opage{66}negotiations, wars, conquests and so on one would, if
one did not know it beforehand, hardly be able to guess
that they were --- what they after all really are ---
Christian states. Such coats of arms as lions,
leopards, eagles, double eagles and the like are on the other hand more
characteristic, for in them there is something of the beast of prey.

The age of enlightenment and freedom has already come
so far that public opinion has become a power:
politics, which is so independent of conscience, is
by no means independent of public opinion. When
the great beast of prey wants to advance, it scouts on all sides and
cautiously draws in its claws if it notices that
public opinion is becoming seriously angry. But public
opinion, a product of moods, considerations of the understanding,
interests and passions, is for all
its sound sense nevertheless no reliable guide; it is
the child of the moment, too much a child of the world. A
clear-eyed daughter of calculating politics, witty and shrewd,
but by her many suitors incessantly exposed to all sorts of
temptations, public opinion sometimes acquires a
disquieting likeness to the kind of ladies whom in Paris one
calls the demi-monde.

Conscience, the conscience of society, on the other hand,
is a reliable guide. But the conscience of society
forms only gradually, as a
popular mind and a popular will form. Public opinion
cannot on its own get the better of the great beast of
those who hold power. Conscience must come in: when
the movements of the beast betray that it hungers after prey
and thirsts after blood, the popular will must say No, and
the soul of the people through millions of eyes fix its gaze
on the beast. It works: the eagle's wing is lamed, the 
% --- p. 67 ---
\setcounter{footnote}{0}%
\opage{67}leopard forgets its spring, the beast of prey is
tamed.

The scientific view of the course
of history is not very comforting. It goes with the peoples
--- so it is said --- as with individuals, they have their
ages. First a youth with golden visions, mighty
striving and strong faith; next a manhood
with keen understanding, sure plans and a well-calculated
energy of action; at last comes old age with the wrinkles of criticism,
doubt of everything except the material, inclination
to rest and enjoyment, evil dreams, omens of dissolution
and death. The same laws by which the states of antiquity
sank hold also for the present. The nations of Europe
are already strongly on the decline; it goes
downhill with them all. Science goes forward,
but religion goes backward. Scientific
research must necessarily dissolve what is ecclesiastically handed down;
it cannot help it. It cannot, alas,
replace the loss. Science is not able
to give the people what it sorely needs, a new religion.
The common people and a part of the half-educated hold
to the old faith, as best they can; but it
dwindles imperceptibly from them, it evaporates. The end
will be that the peoples now in history must at last step
down, and new tribes with fresh roots and young shoots
fill the stage.

On this occasion we must again emphasize what
from a scientific standpoint is so easily overlooked,
that the laws of history are laws of spirit. The peoples do not age
by the same law as individuals; their old age
sets in only when the wellsprings of spirit dry up.
The nations of the present live in 'this respect under other
% --- p. 68 ---
\setcounter{footnote}{0}%
\opage{68}conditions than those of antiquity. At the turning-point of the times
a life has been sacrificed; a capital of life has been put in, which
with interest and compound interest is inexhaustible. Science
shall not burden itself with giving the peoples a
new religion; there is no need. Every people must
become answerable for how it has known how to appropriate
and apply the rich inheritance.

Instead of the dark picture of nothing but shadows,
weakness and old age we can with more reason sketch
a picture of rejuvenation, a picture of the future in
bright colors. We know indeed that nothing is perfect on
earth, that the real condition, whose features are to
answer to the picture, will be reached only approximately; but
for that there may yet be truth in the picture. It is
joyful to hope; yes, it is more than joyful, it
is a duty. It is the rejuvenation of the peoples, the growth of popular life,
for which we hope.

When it comes, then, the expected time, when
the people has become a full-grown man, who has fully
learned to use his powers, the spiritual as fully as
the material, then barriers will fall that now seem
insurmountable. If society, the great body
with the many limbs, is freed from all that is called
dead flesh, and gets its proper mobility, sound juices in
all tissues, fresh blood in all veins: then there will be
capability to take up the struggle of reality, and courage
to meet the possible dangers. If moral self-respect
dawns upon the mind and thought of all, if the conscience of the people is
so awake that willfulness must become shamefaced,
obstinacy creep into hiding, and short-sightedness
and presumptuous stupidity yield everywhere to better insight
and to clear reasons: then unity will be great and
% --- p. 69 ---
\setcounter{footnote}{0}%
\opage{69}the will strong; for it grows with understanding. If it comes
some day, the expected time, when the burden of labor
and the wages of labor shall be weighed against one another according to the
claim of equity on the golden scales of justice, then
'many hands will work to one end, and the many
interests mesh into one another like wheels and pinions in
a well-joined machine. If it is one soul, one
spirit, one stream of life that drives the works with the
many wheels, then the burden will not be too heavy for anyone,
and no one will be thrust outside: the spacious workshop has
room and a sphere of action for small and great. Yes, when it
comes, the expected time, it will be a time of summer:
light clouds and cool winds, bright nights and sunny
days. It will be a time of harvest; then the kernel swells
in
heavy ears; then ripe fruits are plucked from
the tree of freedom; then the apple of knowledge has no poison; the gnawing
is stifled, the adder in the tall grass is dead: then all
the roses of life and health send forth their fragrance.

% --- p. 70 ---
\setcounter{footnote}{0}%
\lecture{Fifth Lecture.} {(Kolding).}
\opage{70}A sound and vigorous popular life is impossible without a corresponding
spirit of the people, the spirit of the people is impossible without the faith of the people;
but the faith of the people is barren, impotent and meaningless
when it is unable to prove its power in deed.
What we here call the faith of the people is not superstition, the belief of the common people
in old legends of the underground folk, ghosts, wizards
and witches, a faith from the nursery; the faith of the people is
a faith that endures the light, a faith in the power of light and of life,
a faith in spirit. When one otherwise asks about
the faith of a people, one usually means its confession of faith.
One then easily comes to make the observation that
one dare not without more ado infer from faith to life, that
the confession of faith as a rule points so far
beyond life, that one must choose another line of sight
when one wants to have the real life in view.

The faith of the people is more than an exalted feeling or
a beautiful imagination; it is a faith in reality, a
living faith in deeds; as the faith is, so is life,
and as life is, so is the faith; here it holds
unconditionally: show me your faith by your works. To be sure,
the Divine in which one believes is something
% --- p. 71 ---
\setcounter{footnote}{0}%
\opage{71}superhuman; but this superhuman is, in its
immediate presence, nevertheless wholly and entirely human:
all the principal features of man's nature and being
are recognizable in its revelations.

In this respect it might then seem as though
the heathen gods, the old gods of the peoples, suited
far better the nature of man and the peculiarities
of the races than the eternal true God of the spirit, the
God who has revealed himself in Christ. The Greek
gods have Greek faces, Greek virtues and vices;
they are Greek human beings of a higher order. Something similar
holds on the whole of the gods of the Romans and of the Norsemen.
But precisely because they are only human gods,
gods with human faults and merits, they bear
the seed of death in their hearts and, for all their immortality, can
live only for a certain time.

It was the religion of the Aesir, the faith in the gods of lofty Valhalla,
that held up the spirit of the North and carried the popular life forward
on strong shoulders; but only for a certain time. In the 8th
and 9th centuries after the birth of Christ the relation is
reversed; then it is men who bear the gods.
The hero relies on his strong arm and his sharp sword:
he believes in Thor because he believes in his own
valor. What, after all, does the mysterious element
in the fate of the gods, which only the Norns know, the riddle
over which Baldur broods in melancholy, Mimir's well,
where Odin has set his eye in pledge, signify, other than the
dim, unexplored depth in the Northman's own inner being? With
all their courage and their sense of independence and freedom, the
strong men nevertheless feel themselves constantly subject to dark
powers; there is an inflexible necessity, something sorcerous
and wizardlike, that ensnares them in its net.
% --- p. 72 ---
\setcounter{footnote}{0}%
\opage{72}They are silent, shut up within themselves, sparing of words, when they hatch
plans and proceed to action. The mind is weighed down, it
is the yoke of necessity, the unfree, nature-bound being.
A Sigurd Fafnir's-bane can slay the dragon and lift
the treasure, but the serpent of fate, the great dragon of
enigmatic existence, not even the spirit of the people
itself is able to slay. Thus arises the presentiment of an All-father:
Harald Fairhair believes in one single, almighty God;
thus there is, in the midst of the din of battle, a quiet longing
for peace, a search for light, a sigh for deliverance.

In his incomparable romance, ``The Death of Hakon Jarl,'' Oehlenschläger
has presented this transition in a symbolic
manner. ``The nights brood so long and so black,
the Seven Stars glimmer so faintly'' . . that is a Nordic
winter night, the dark night of heathendom, whose cold breath
we sense. Around the sacrificial mound stand the gods, old % defect: tornemme (a broken f)
and gray, in cracked stones: ``Our time is past, soon
we shall sink''. The counter-image to this is a shadow-outline of the
Gothic church with the red walls. Through the leaded window
the moon's ray reaches the crucifix of the altar and
shines with a promise: ``You will surely conquer, you white
Christ, soon shall the wild northern zone bow to
your crown of thorns''. Then reality begins: ``Olaf Trygvesen lands in Norway''. Masses are sung, the banner of the cross
waves, ``The halberds wildly in the aether tumble,
they flash, they flash away''; it is a popular collision, a
bloody tryst between the old and the new faith.

Seen from an evangelical standpoint, flashing
halberds and clanging swords naturally cannot be used as
proofs of the truth of Christianity; but when the
discussion is of a faith of the people, then truth is a power,
% --- p. 73 ---
\setcounter{footnote}{0}%
\opage{73}vital power, heroic power: show me your faith by your works.
--- Human life cannot bear itself and still less
bear its own gods. The holy God of truth,
who bears all, is infinitely exalted above the
human; but precisely therefore also --- like love
itself --- the true God of men, a true human God.
The word of true life must be impressed upon
the hearts and win the souls.

More than a hundred and fifty years before Olaf Trygvesen's
appearance in Norway, Ansgar, an apostle of peace
from the monastic cell, had proclaimed the Word among us. Mild as
the Gospel, gentle, self-denying like one who
treads in the footsteps of the Savior, courageously patient like one
seized by the Spirit who yearns for the martyr's crown,
Ansgar had by the warmth of his word already begun to melt
the ice, thaw the cold minds and kindle altar lights around
those clad in mail. So it went both ways, that of the Word
and that of the Sword, forward with the transformation of the faith of the people.

But however much that is evangelical there may be in Ansgar's
teaching and personal life, we must nevertheless not forget
whence he comes. His Gospel is not the
originally pure one; a color has been put on it. It is
Catholicism, the Roman, not yet fully developed
Catholicism, that he brings with him for the conversion of the North.


How high Catholicism has lifted the faith of the people,
what change of mind it has effected, how many
chains it has loosed, how it has sheltered the
defenseless, ennobled woman, taken up the cause of the poor and of the
serfs and labored to bring unity, peace and
order into the wildly confused society, shall be remembered
with gratitude. Our Gothic churches: the
% --- p. 74 ---
\setcounter{footnote}{0}%
\opage{74}buttresses, the pointed arches, the high vaults and the slender
towers that point toward heaven, are more than monuments
of a peculiar art of building, they are monuments
of a peculiar stage of popular life, testimonies to
the power of faith in times gone by. From these silent walls
a petrified spirit speaks of vanished might and splendor.

But Catholicism has from the beginning committed
a sin against the Spirit; there is a fault, a single
fundamental fault, from which all its excesses, mistakes, indeed even
its crimes against the spirit have their origin; it lies
in the conception of the concept: \emph{Authority.}

After the pattern of Christ the holy Word is to be proclaimed
with authority; but after the pattern of Christ
this authority is then also expressly to derive
from spirit, i.e. from convincing truth, self-sacrificing
love, personal superiority. How has
Catholicism now conceived and applied the authority of Christ?
It has quite rightly taken its starting point
in the spirit. But however much it has been zealous to
get the personality of the Holy Spirit fixed as dogma,
it is nevertheless unmistakable that in the ordering
of the whole ecclesiastical system it has come to do what
it did not wish, in that it has in reality come
to make its Holy Spirit as impersonal as possible.

Upon the dogma of the immediate presence of the Spirit
in the infallible assemblies of bishops of the infallible Church
the Catholic doctrine of spiritual authority is founded.
The teachers and rulers of the Church thus do not have the
divine authority, like Christ himself
and his Apostles, from the personal element in their own inner being.
On the contrary! the religious-moral inner being can, as with
other mortals, find itself in much frailty;
% --- p. 75 ---
\setcounter{footnote}{0}%
\opage{75}their authority suffers no diminution because they are not
paragons of virtue. The clerical persons are not to be judged
by their own personal worth, but by
the position and rank they have in the Church. They are the Church's
organs, its servants and yet at the same time its bearers:
their authority is the authority of the Church system, personal
in what is fallible, impersonal in what is infallible. The Church system
is outward and the ecclesiastical persons likewise.
The Church cannot act except through its persons;
the whole Church system must then in reality comprise
both persons and Church. But the persons can indeed
err; in order that the Infallible may free itself from the Fallible,
a miracle must occur, a miracle of the Spirit,
every time the bishops assemble. That the bishops nevertheless not
so seldom have been assembled without the miracle having
occurred is notorious, since the one Catholic assembly of bishops
has indeed solemnly condemned the other. The assertion
of the miracle and the authority corresponding to it
thus becomes a purely outward assertion. The authority
becomes not less, but even more outward, when
a single bishop, a single, in and of himself fallible person,
is, by virtue of the miracle of the Spirit and the Church system,
to be transformed into an infallible one.

On a groundless miracle the spiritual authority was
founded and brought to completion in the Middle Ages.
How outward the authority was is seen from the
ecclesiastical separation between clergy and laity. All
the clergy were of age; all laity, i.e. non-clergy,
were --- in the spirit --- under age; with that the
matter was settled. That such an authority was by no means,
as one from the Protestant side has sometimes
wished to pass it off, an arbitrary invention of clerical
% --- p. 76 ---
\setcounter{footnote}{0}%
\opage{76}lust for power, is clear from the whole ecclesiastical way of thinking.
It shall be willingly conceded that outward authority was precisely
suited to semi-barbarous conditions in a semi-barbarous
age. It shall be conceded that a sensuously tangible authority
was needed in order to have an educative effect
upon the rude masses. But if we consider the
educative element in this authority more closely, it also
shows itself that precisely in its outwardness it has its
limit.

If the point is to produce sudden and
strong religious impressions, subdue coarseness of mind, awaken
deeper longings, chastise doubt and take obstinacy
by force, then surely there is no church system that can
measure itself with Catholicism. The same holds
when the task is to nourish a mystical inwardness and
to bring the outward life with its countless particulars
under strict general rules. But further one does not
get; in inwardness it remains with the
mystical: outward reality must make do with the
sheen of piety which a colored light through the windows
of the Church spreads over ``the good works.'' Masses
and images of saints, choral songs and incense, ringing of bells
and processions, all such things are suited to fill
a sensuous imagination and make the spiritual
outward. To kneel before the Mother of God, say one's rosary,
lie with bare knees on peas in order to atone for a
church penalty can be disciplining enough for those under age;
but it is a sin against the spirit that understanding
is lacking, that those under age never get so far that they
could ever become of age.

At the beginning of the iconoclastic controversy we have a word from
Gregory the Great (590 --- 604); he says: ``Let them
% --- p. 77 ---
\setcounter{footnote}{0}%
\opage{77}read upon the walls what they do not read in the books.''
Against this there would be nothing to object, if we
dared understand the words so that the provisional
reading upon the walls were as soon as possible to be completed
by a thoughtful reading in the books; but to
educate a people to religious thinking for themselves would, what
experience has sufficiently taught, be the same as
to work against outward authority.

Fiat and outward authority are most closely
united; outward authority wants, what
Christ and his Apostles never wanted, to rule by
the use of outward means and outward weapons. Hence
the long embittered struggle between papal power and
princely power; the former was in the eyes of the Church like the sun, the latter
like the moon, which must borrow its light from the sun. Among
the many outward means none is so effective as
money. The Church needed money, always money, more
money. When now not only money, but also remission of sins
is an outward thing, then there was, one would
think, no contradiction either in exchanging the one
outward thing for the other, thus no contradiction at all
in selling remission of sins for money. But --- strangely
enough --- this clear logic the peoples could not
understand. That the Pope took their souls, their conscience,
their inner freedom and spiritual independence, they put
up with patiently as good Christians; but
that he took their money, that was revolting. Against the
revolting one revolts; at last the measure was full:
the Reformation broke through.

At the Imperial Diet at Worms (1521) the world saw an
example that showed what spiritual authority can do,
even if it meets in outward lowliness, and presents itself
% --- p. 78 ---
\setcounter{footnote}{0}%
\opage{78}in an unprepossessing form. Imagine the
glittering assembly; the Emperor, the princes, the papal
legate, the Cardinal, archbishops, bishops,
abbots, what representatives of worldly greatness,
of ecclesiastical power and authority! And then the humble
monk, the accused, whose defense they did not wish to
hear, but whose disobedience they would sternly punish,
unless he humbled himself under the infallibility of the Church and
recanted all that he had spoken and
written. It was a great moment when the official
of Trier broke in upon Luther and demanded a plain and
straight answer whether he would recant or not; for
Luther answered: ``Unless I am convinced by the testimony of
the Holy Scripture or by public, clear and evident
reasons and caught in my conscience
by the Word of God, then I can and will recant nothing,
for it is neither safe nor advisable to do anything
against one's conscience.'' Here was a newborn conscience
with a new authority; the inner broke forth like
lightning from the dark cloud, and the lightning kindled.

What awakenings and hopes, what streams
of spiritual life the Reformation brought with it is generally
acknowledged. Luther made the cause of the Reformation a
cause of the people. ``Every true Christian is of the clerical
estate; the cup in the Supper belongs to the people; it shall
stand free to everyone to believe the newly invented doctrine of transubstantiation
or any real communion whatsoever
between the consecrated bread and the divine body.
Against the whole significance of the outward Church is set
the sole saving omnipotence of faith. Therefore the
Christian freedom must not be bound by any proposition and
arbitrary vow'' etc. But the new Church, which
% --- p. 79 ---
\setcounter{footnote}{0}%
\opage{79}gradually formed itself, dared not stand weaponless
as Luther at the Diet. Through the Holy
Scripture one went back to the original: instead of
the authority of the Pope only the authority of Christ held, and
the congregation of the Apostles became the pattern of the evangelical Church.
If this thought of life, original in Luther,
were to have been carried through, then the
young Lutheran congregations must also have put up with
remaining at the conditions of life of Christ and the Apostles;
but voluntarily to choose suffering and martyrdom by such
a measure, for that there was, with all the inwardness of faith
and the heroism of sacrifice of which that
age has so many and beautiful examples to show,
nevertheless not strength of spirit enough. So one let the conditions
decisive for spiritual
authority remain outside,
lowered the claims of following the pattern and held to the
doctrine, ``the pure doctrine.'' Despite its almost unbridled
zeal for the pure doctrine the Church came to terms with
the historical circumstances and submitted to the
state. The inward authority gave up and fell into the
arms of the outward: thus our Lutheran
Church became a State Church.

If the people in this country could either overlook
or fail to appreciate how much that is noble and good our Lutheran
State Church has in the course of time furthered among us, how
much has been accomplished, as well by God-fearing kings as by orthodox
clergy, for the comfort of souls,
for the purification of morals, for quiet progress in cultivation of the spirit
and general enlightenment: then our people would
sadly lack --- what happily it by no means
lacks --- historical sense. Not to speak of
the fact that it was the State Church which with one blow struck the
% --- p. 80 ---
\setcounter{footnote}{0}%
\opage{80}hierarchy to the ground, so almost every step
forward, every measure by which the present
state of freedom has been prepared, can be traced back immediately or
mediately to the religious and civil influence of the State Church.
When a people is on the point of forming and working out something new, then it is a
duty of gratitude to take note of the merits
which the old, that is now to be shelved, has really
had. By fulfilling this duty one best secures oneself
against the rash error of discarding precisely that
which necessarily must be preserved.

In one point the Lutheran State Church has sinned;
it has, like Catholicism, set outward
authority above inward. Luther has with great
force emphasized that every sacrament that is administered
in the Church must expressly have Christ's command. Since
this holds only of Baptism and the Supper, he rejected
the Catholics' other five sacraments as self-made.
The old State Church now concluded thus:
if the sacraments have Christ's command, then it must above
all be the duty of the Christian authorities that Christ's
command be punctually fulfilled. To make sure of this,
a command of the authorities was then added to Christ's command:
outward authority stepped in to help
the inward. It was a monstrous mistake, not
to use the harsh word: a crime against the spirit. Where
Christ commands, he tolerates no interference from other
commanders; his authority is violated when unconcerned
helpers step in and want to support it with an
authority of quite another kind. Christ's command is to
faith, and faith is free: free faith, free obedience, free
appropriation are the first, indispensable
% --- p. 81 ---
\setcounter{footnote}{0}%
\opage{81}fundamental conditions of the religious life. Had Christ wanted to put through his commands
by the use of compulsion: then
he had indeed power, and could have used power. But that
he will not permit any outward authority whatsoever
under any pretext whatsoever to meddle in
that which by its nature belongs under spiritual authority,
he has proved in deed by walking the way of suffering
and letting himself be crucified. Luther understood Christ,
when he declared that it should stand free to everyone
to believe --- imagine --- ``the newly invented doctrine of transubstantiation''! But, alas, the Lutheran State Church
did not understand Luther.

Just as in Catholicism, so also in the
Protestant State Church outward authority has dug
its own grave; but, it goes without saying, the grave is
another, for here is another gravedigger.

The eighteenth century has in a special sense
received the name of the age of enlightenment, in particular of the enlightenment of the understanding.
With advancing science thinking was freed;
the investigating understanding threw itself upon everything; it wanted to comprehend everything, and rejected everything that
it could not comprehend. With blazing torch
Enlightenment pursued goblins, trolls, ghosts and
underground folk, who fled as fast as they could; at last
it took on the Devil himself. According to the orthodox
principle, one was to take reason captive under
the obedience of faith; but the obedience of faith was indeed the
same as obedience to the doctrine of faith prescribed by the State Church, a blind submission under an
outward authority. Such a thralldom of spirit
reason would no longer put up with. With the help of the understanding science soon brought it so far that it
% --- p. 82 ---
\setcounter{footnote}{0}%
\opage{82}itself became the supreme court, and religion had now to
submit to the authority of reason.

Holberg, who must be regarded as the first champion of Rationalism
among us, expresses himself on this among
other things thus: ``I have shown that Orthodoxy itself,
which grounds itself alone on the authority of others, cannot
be pleasing to God. For just as a great artist
finds no pleasure in having his work praised by
a blind or unappreciative man, who has no eyes;
and a musician cannot rejoice that his music is
praised by one who has no ears, so God can find no
pleasure in those who judge of his glorious
attributes, of which they have not troubled themselves to get the
least idea. --- If a hymn of praise or worship is to please,
then it must be made by one who knows why
he worships; and if a faith is to be called faith, then it must
ground itself on knowledge, so that the phrase: \emph{I
believe}, must be the same as: \emph{I know.}

The enlightenment of the understanding did not limit itself to
shaking the existing orthodoxy from a standpoint
outside the Church, it crept in through the church door.
Under the guise of an enlightened, rational orthodoxy
Rationalism began by criticizing the Bible and ended
by explaining away the miracle and diluting Christianity.


At last a turning point set in, where the strife
between science and the orthodox, officially recognized
doctrine of the Church became so violent that it brought the authority of the state
itself into embarrassment. The state had undeniably
committed itself to uphold the old Lutheran
orthodoxy; it had had all its clergy swear
an oath on the symbolic books. But on the other hand
% --- p. 83 ---
\setcounter{footnote}{0}%
\opage{83}the state had indeed also to hold to enlightenment and recognize
the right of science; here was a crossroads. It
showed itself, to mention an example, when the royal
Danish Chancery in the year 1834 came to a clash with
``the strong Jutlanders.'' The old-fashioned believing common people
namely could not come to terms with the new textbooks,
but sent the children to school with the old ones.
It was regarded as school neglect, and the Ordinance % defect: Skolefcrsømmelse (a broken o)
of 1814 was applied. Fining and
distraint were to convince the simple people of
the excellence of the new books; but it bore no fruit.
The poor peasants put up with distraint and sale of
their furnishings, their sheep, cows etc.; but in the use
of Balle's textbook they would not acquiesce. They did not limit
themselves even to simple refusal, they adduced
reasons, and gave a list of the passages that
stood in conflict with the old faith and had aroused their
offense; but the high authorities could naturally % print: Øvrigked (misprint)
not enter into religious grounds of proof.

Here stood right against right. The peasants really had the
right to demand the old strict Lutheranism upheld;
it was indeed precisely this upholding to which
the supreme authority of the state had solemnly committed
itself. But the Chancery also had right; it maintained the
authorized textbooks, because they were \emph{authorized.}
To yield to insubordination against the
Authorized was to weaken the authority of the state. In order to
uphold its authority the state broke its obligation:
the State Church authority had fallen into
discord with itself.

To cover this discord with a suitable veil,
one has appealed to something called the
% --- p. 84 ---
\setcounter{footnote}{0}%
\opage{84}Protestant principle, and which is supposed in particular to consist in this. that
the Lutheran doctrine in the name of the spirit and of freedom
lies open to critical examination and lets itself be reshaped,
according as the interpretations of Scripture with correspondingly adapted
elaborations of the dogmas vary, whereas the Catholic
doctrine of the Church, petrified and crystallized in old forms,
must defend itself against free science by withdrawing
from its criticism and applying fiat instead of
reasons. The essence of Protestantism is development: it
seeks the truth and keeps pace with the time. The essence
of Catholicism is standing still, it maintains that it has the truth
and wages a desperate struggle with the spirits of the new
time. Herein is something true, but also something misleading.


In comparison with the loftiness, character and
coherence which distinguish Roman Catholicism,
the Protestant State Church system of the present looks
wretched. Catholicism knows what it wants, on
what foundation it stands and what weapons it is to
use; it still has in its way authority, spiritual authority.
That is more than can be said of the
sadly crippled State Churches. On what
foundation would they stand? Some appeal to the
symbolic books, juridically, juridically! But when the juridical
authority of the symbolic books is literally, strictly
to be upheld by the power of the state, what then becomes of it
with the Protestant principle? Others appeal
to science: Criticism, criticism and more criticism!
But if it is really scientific criticism,
which in the highest instance is to decide what religion is,
and how many scraps of the Bible still may
% --- p. 85 ---
\setcounter{footnote}{0}%
\opage{85}have leave to be genuine: what then has the state authority
to signify, and what becomes of the Church?

Here is a task to solve. How its solution
is possible, especially for us under our conditions, is a
question whose full answer will be able to emerge only from
very detailed discussions. What I on
this occasion have to offer is a brief and definite
statement of the points of view which in my conviction
are and must be the decisive ones. I shall
endeavor to elucidate these points of view in a following
lecture.\footnote{See the Seventh Lecture.}
% --- p. 86 ---
\setcounter{footnote}{0}%
\lecture{Sixth Lecture.} {(Aagaard.)}
\opage{86}``Man dies, but never dies man's memory'', that is
an old Norse saying, and it bears witness to a high-minded
view of life. Those who live only for self-interest
had best not mention this saying, for it threatens them
with a judgment: base self-interest, crude egoism,
may achieve much in the world, but there is something it
has never attained and never attains: an honorable posthumous reputation.
What holds of the individual man holds
in a wider sense of an entire generation. We all know
how tempting it is to look to the nearest
advantage. When those who cling to what exists themselves
must admit that the fall of the old building is unavoidable,
and yet, for the sake of advantage, the nearest advantage,
cling to the rotten pillars and set against
every thrust of the future the poor comfort: ``it will hold
well enough as long as we live'': then such conduct gives no honorable
posthumous reputation. When a party of progress makes use of what it
regards as a favorable moment in so reckless
and passionate a manner that it must necessarily confuse
and harm the nearest future, when the watchword becomes:
``Now is our time, we must use our advantage, let the
% --- p. 87 ---
\setcounter{footnote}{0}%
\opage{87}more distant consequences be as they will!'' then such conduct gives no
honorable posthumous reputation. The verdict of posterity becomes a
sentence of punishment, a sentence of condemnation; ``Man dies, but
never dies man's memory.''

The verdict of the future is more than a human verdict,
it is a judgment of God. Through history there runs a
connected series of judgments of God; it is only a question
of how far we human beings are able to
discover them, and in what depth we are able to recognize % print: bvilken (misprint)
them.

The Middle Ages took the matter directly after the pattern of the Old
Testament; David felled Goliath, that was
a judgment of God. When a lawsuit or a matter of honor of
great importance was to be decided, then it did not essentially depend
on the witnesses and the legal hearings, no, the decision
lay in something else, for example in a duel. If
insulting accusations had been raised against the character of a noble
lady, her knight seized his sword; and
if he defeated his opponent, the latter's statement
was at once overthrown as untrue. It was a faith of the people
that God would by a miracle let the righteous man conquer.
If we consider the judgments of God properly so called, ordeal by iron,
trial by the sacrament, trial by the cross, and so on, then the
characteristic feature shows itself. People believed that the bread on the altar,
when a certain formula was recited at the Mass, would
turn over if the suspect was guilty. Or
they wrapped some pieces of wood, one of which was
marked with a cross and another not, in wool and laid
them on the altar: if the suspect now took the
piece with the cross on it, he was declared innocent,
otherwise guilty. Drought, plague, famine
were regarded as immediate arrangements of
% --- p. 88 ---
\setcounter{footnote}{0}%
\opage{88}God for the punishment of sinners and the advancement of righteousness.
Even a mighty enemy's natural death was a hint
from above. When the Emperor Frederick the Second, the Hohenstaufen,
was dead (1250), Pope Innocent the Fourth burst out in an
address to the clergy of Sicily: ``Heaven rejoices
and the earth exults, for the lightning and storm with which
God the Almighty has so long threatened your heads
have by this man's death been transformed into
a refreshing zephyr and a fertilizing dew.'' Even
some hundred years ago the Spanish
clergy wished to forbid the performance of pieces consisting wholly of
song (operas), ``because this ungodly entertainment
was to blame for the prevailing lack of rain.'' It has
indeed since been urged that such religious interpretations
were well suited to awaken consciences
and make the idea of God's omnipresent
justice vivid; but it is nevertheless obvious that
if the Almighty by interposing miracles himself
were willing to undertake the administration of justice, then it would be
impossible for so much injustice to thrive in
the world.

The modern age has reached the opposite extreme.
The divine freedom has caught itself
in the laws; the unchangeable laws are bonds which it
cannot loosen: in the connection of natural causes,
conditions and intermediate causes an inflexible
necessity rules. It is declared outright that one will not
hear of ``an irresponsible God''; and with the responsibility of men,
before the tribunal of history,
it likewise has little to signify. If the acting
persons are considered one by one, then a distinction must surely be made
between good and bad, noble and ignoble characters:
% --- p. 89 ---
\setcounter{footnote}{0}%
\opage{89}praise and blame cannot be wholly excluded. But if we grasp
the whole with its countless particulars in connection,
then we are convinced that the great men of history have always acted as they,
according to the circumstances,
necessarily had to act. A moral
standard for good and evil, just and unjust,
cannot be applied to world events and the
mighty spirits, the shopkeeper's scales of civic morality are
too small here. The great will be neither praised nor blamed,
neither loved nor hated; it will, in one word, be admired.
A Tiberius and a Nero have just as much
attraction for the historian as either a Paul
or a Luther. World history is the world's judgment:
the judge is like Fate itself, and the judgment without
conscience.

Should anyone venture to meet this grandiose
view of the world with religious-moral weapons in the simplicity
of faith --- as David met Goliath, then
one would certainly with a superior smile remind him
that David is dead, and his faith a long since abandoned
standpoint. ``A judgment of God'' one does not fear;
one is sure that the miracle will fail to appear, and so on. ---
How remarkable! In the Middle Ages Omnipotence was free,
but the whole age unfree; in our days the age is to
be free, but Omnipotence itself unfree. There must be
something behind this.

Every time a great goal with a great task
is set into history, it is Omnipotence itself, the eternally
free spirit, that opens the depth of possibilities and sets
men the task. The task is not pronounced in plain
words from heaven; it is through existing facts
and inner impulses that it announces itself.
% --- p. 90 ---
\setcounter{footnote}{0}%
\opage{90}A historical task does not come naked into the world;
it is clothed in the many-folded garment of natural conditions, events, circumstances:
it has a veil over its
face. The human will awakens as the Veiled One
shows itself. It is freedom that awakens freedom, will
that speaks to will. As the task is more closely
determined and its demands step out into reality,
it sets off a multiplicity of separate,
subordinate tasks: here is work enough in all
directions. The human will, which chooses, acts
and works, has in all parts and from all sides, whether
it knows it or not, to do with the divine will.
Our freedom is limited, our will in
countless ways conditioned by circumstances. But
circumstances are colored by the light of knowledge, and
the obstacles themselves give the will its elasticity. Different
circumstances may lead the human will
into highly different tracks; but there is no circumstance
that can prevent our will from willing.

The manner in which millions of wills work together
and in relation to the common task form a common will,
a ruling general will, characterizes the \emph{spirit}
\emph{of the time.} On the spirit of the time the children of the time all without
exception, the mightiest as well as the humblest,
are dependent, and yet in relation to the task all without
exception, the humblest as well as the mightiest,
as regards possibility, independent. Everyone can
will and choose in his own way, the one the good, the
other the bad; but Providence makes them all responsible.
When the set task has then been worked at
so long that its content is exhausted and the forces worn out,
then its age goes into its grave and gives
% --- p. 91 ---
\setcounter{footnote}{0}%
\opage{91}place to a new time. In so far as at every
such transition, whether it takes place quietly and slowly or
suddenly as in a spiritual earthquake, we give heed to the spiritual will's
infinite superiority over the finite, yet each
severally responsible, human wills, we recognize a judgment of God
in history.

From this point of view a peculiar light falls
on the world-historical struggle between right and wrong,
good and evil, which has aroused doubt in so many
thinkers. The ruling spiritual will does not award us
possession of this or that earthly good; it
will, in the eternal sense, help men forward in willing and acting, but it does not take sides directly with
the good. The good which it historically wills, it lays into
the great task. Men may seize this
task from whatever side they will; those who fight under
its banner have on the whole Providence with them. On the task's account
crimes can be committed; indeed
the crimes may even be the greater, the
deeper the task lies, the more historical right it has.

Boccaccio, an Italian poet of the fourteenth century,
has in his ``Decamerone'' a series of tales,
among them the following. Two merchants, a
Christian and a Jew, were very good friends. The
Christian merchant had had so many proofs of his
Jewish friend's honesty, faithfulness, readiness to serve and
upright way of thinking that he, being himself a good
Catholic, felt an inward distress at the thought that
so honest a soul should be lost. What
did he do? He naturally sought to convert him
to Christianity. But as he soon noticed
that his own arguments of persuasion would not do,
% --- p. 92 ---
\setcounter{footnote}{0}%
\opage{92}for the Jew was sure of himself when he defended his religion,
he thought to himself: ``I must get him to the priest.''
But that did not work. ``To the bishop, then.'' But that
did not work either; and so the good man at last did not
know what he should do. One day his
friend steps in to him and says: I want to make a journey to
Rome; I have business there, and wish to look
about me in the famous city. --- Good Lord, sighed the Catholic
to himself, if Abraham comes to Rome and sees
that Sodom and Gomorrah, then I shall never get him
converted; now my last hope is gone. --- The Jew set out
and was away for a longer time. When he came back there
was talk of commercial affairs and all manner of worldly
things; but that point of religion the Christian
friend dared not touch, he dared not even mention either the Pope
or the cardinals. How astonished he was when
his friend suddenly fell silent, regarded him with a peculiar
mild look and then said: I will confide something to you;
I have resolved to become a Christian. --- Is it possible?
How could you, in the midst of the pool of insolence, blasphemies
and debaucheries, get good thoughts of Christianity? It is clear, answered the Jew, the God who
can bear such things must be more merciful than mine.

What is here said as a biting mockery can be taken
in full earnest. The height of crimes to which
the medieval Church at last carried things is
at the same time a measure of the religious depth and moral
greatness which the task originally contained. By
this standard we must recognize the judgment of God under which
it sank.

Let us cast a glance at the period that directly
borders on our own, and yet lies closed
% --- p. 93 ---
\setcounter{footnote}{0}%
\opage{93}for us, namely the age of Napoleon the First, and there we
see an example of what a judgment of God is. The task
that was here set was not directly religious; and yet
the spiritual will was with it, it was in truth a great task, significant
for the cause of humanity (humaneness).
The Revolution broke brutally with the historical
past. In its terrible wrath it wanted to break
the old yoke, as one breaks a withered branch
on an old tree; with a bloody axe it wanted to hew
the tree down, and did not notice how deep the roots were.
At no turning point in history have the wildness of the masses
and the impotence of the multitude, where a guiding Providence is concerned,
shown themselves in a clearer light than just then when ``the Mountain'' sank
and ``the Valley'' began again to rise. Despite astonishing
upsurges, immense exertions of calculating
shrewdness, fanatical zeal and a robber's courage that staked
everything, it nevertheless looked as though the ``Rights of Man''
were to be stifled in the inhuman, and the defenders
of the old servitude to return with tenfold
right. When Groa, says the myth, sang magic songs
over Thor, the god became so glad that he told her
how he had carried her husband, Ørvandel,
across the Elivagar. Superhuman powers were required
to carry the new society, convulsed to its innermost core,
across the Elivagar of the Revolution. But the spiritual will
was with it; there rose out of the depth of possibilities a
heroic figure, a holder of power of a higher order, a
world-historical genius, Napoleon Bonaparte. He had
a will.

That such a will cannot possibly be a simple,
naturally necessary product of circumstances, shines
forth most clearly from this, that it everywhere, to the
% --- p. 94 ---
\setcounter{footnote}{0}%
\opage{94}world's astonishment, makes itself the master of circumstances; it brings good fortune to go with it and teaches the
chances it meets to bow deeply before the superior man.
So the new world had received a new will,
a central will. Millions and millions of wills, strong
and weak, small and great, are set in motion. Surprise,
envy, admiration, friendship, enmity,
hope, fear, hatred, enthusiasm, freedom, compulsion, all
one! all thoughts, all moods, all wills must
join in: they must all revolve as % print: sóm (accented for stress)
planets about the one
sun, all move about the one central will. Here
is freedom in the midst of this necessity; here is a world
of wills.

If we look at the course it took in the solution of the task,
then the hero of the age unfortunately forgot on the way
that it was the ``Rights of Man'' from which he had
set out. Instead of being the servant of the spiritual will, he became
an instrument of Fate: the imperial eagle swooped down on
the whole like a bird of prey. When one considers how
one step led to the next, how one extension of
power demanded another, how the extraordinary
was heaped upon the extraordinary, Ossa upon
Pelion, as in the days of the wars of the gods, until the colossal
edifice at last crashed down, and the myth of
the Titan, Prometheus chained to the rock, was
made reality on St. Helena: then there is undeniably
matter enough for those who immerse themselves in the
concatenation of events and admire necessity. None the
less this necessity is from first to
last conditioned: outwardly by the clash of chances, inwardly
by the choice of the will. In the interaction of chance
and choice lies conditioned freedom. It is the
% --- p. 95 ---
\setcounter{footnote}{0}%
\opage{95}world-ruler who, from selfish impulse, continually transforms the spiritual will
into the will of fate. The great drama, the tangled
knot, the far-radiating points of splendor, the
shattering scenes and the tragic ending, all is so
interwoven with the threads of freedom and necessity
that Providence and Fate both enter upon the
right of the stronger and challenge each other to a world-shaking
trial of strength. Such a trial was never seen
on the stage of the old world. An Alexander,
Hannibal, Scipio, Caesar are examples of what genius
can do by virtue of invincible will of impulse; but Fate
rules, freedom is overshadowed by necessity.
The heroes of classical antiquity could not possibly have had
such a consciousness of the essence and conditions of freedom
as the son of the Revolution from the eighteenth century.
So the Mighty One becomes responsible, so the spiritual will
asserts its superior power; so the judgment becomes
more than a world's judgment, it becomes a judgment of God.

When it can at times offend us that the right of the stronger
is to hold as fully in history as in nature,
then we dare not forget that it has, on the conditions
of freedom, a peculiar connection with what we call
the stronger. It is said in the language of religion: better
to suffer wrong than to do wrong; but the
meaning of the holders of power is: better to do wrong than to suffer wrong.

There may at times come upon crude holders of power such
an involuntary lust for destruction that they visit
countries with fire and sword and burn down peaceful towns,
having no other purpose than that of wanting to show
the world what a wretched thing it is that one calls law and
right. Nero is not the only one who takes up a distinguished
place among the crowned arsonists. Such an
% --- p. 96 ---
\setcounter{footnote}{0}%
\opage{96}abuse of power is vanishing in all civilized
countries. But with that the ``Rights of Man''
are as yet but little helped. The positions of states toward one another
are, and by the nature of the case must be,
positions of power. To maintain right there is required
power; the proposition, true in and for itself: \emph{Right is Might},
is then turned around and comes to mean: \emph{Might is Right.}
Thereby arises a tacitly acknowledged rule of power
which, more than anything else in the world, undermines
the moral consciousness of peoples and states. It goes
without saying that the holder of power never fails to give
his claims, even if they contain the most
crying injustice, a semblance of right; there is pretext
enough for war when one wants to overwhelm a weaker neighbor.
So wrong is done on the grand scale. And when the wrong
is done, what then? Then one judges by the outcome.
What has happened has happened; the wronged party must
resign himself to his fate. Let the wrong be manifest
to all, so manifest that the simplest man
can feel it with his hands; let the overmighty one
himself see that it is really wrong, and that it only
costs him a word to make the wrong good again: he
cannot --- so weak is he, the strong one, on this
point --- he dares not do it. And why
not? Because it touches his honor too closely; it would
look like weakness. The holder of power who voluntarily
gives up a single piece, however small, of his
power has given himself up and betrayed the whole rule of power.
To make a wrong good is, after all, to concede
its existence. ``Might is Right''; how should a
holder of power humble himself to concede that he ever
either had done or could do wrong. To
% --- p. 97 ---
\setcounter{footnote}{0}%
\opage{97}make a single wrong good? No, no! If one is to enter
upon the particular, then the debt becomes too great;
there stands an army of creditors outside the door. Better
to do wrong than to suffer wrong.

And now public opinion? Politics is not
sentimental, but practical. Public opinion is, from precisely this point of view,
to be disciplined to understand
that questions of right and wrong have their home only in
the internal affairs of the administration of the state, in civil and criminal procedure,
but in no way concern the position of power.
That is no contempt for right and justice. On the contrary!
The higher we esteem right and justice, with the
greater esteem and awe must we look up
to the holder of power who is placed so high that he, without
diffidence, without a blush, can show the world that he
has right under his foot, and teach us what wrong
means. --- That was, then, an explanation of the right of the stronger;
we now pass on to another.

Better to suffer wrong than to do wrong; that is
the poor man's comfort in the eyes of the mighty, and alas:
by being misunderstood it can only too easily become poor
enough. \emph{To suffer} wrong is usually understood
in such a way that the emphasis falls on the bare suffering,
the deedless endurance, which in dullness consoles itself with the thought
that here there is nothing to be done. That is an outrageous notion,
offensive especially when it is to go under
the name of Christianity, for it is an open denial
of the thought of Christ.

There has been One in history, One alone, who
in spirit and truth knew what it was to suffer wrong.
Millions have suffered and still suffer to this very day greater
or lesser wrong; there is hardly anyone among us,
% --- p. 98 ---
\setcounter{footnote}{0}%
\opage{98}who in this respect has not something or other
to complain of. We may be in the right in much, but
we are yet always in the wrong in something. No one dare say that he
only suffered wrong, but never did wrong. Only the one
Sufferer suffered wrong, not \emph{although} he was in the right in
much, but \emph{because} he was in the right in all, because he had never
done wrong. The many sufferers must admit
that what chiefly brings suffering upon them is the lack of
the goods of which they have wrongfully been deprived. Who dares
say that for justice's sake he has voluntarily
chosen to suffer? We do not love right itself so much
as we hate those who deny us our right. Could we
come to our right by abating a little of justice,
we should gladly enough compromise: it is the intractability
of our opponents that forces us to suffer.
Only the one Sufferer, he who voluntarily suffered for the sake of justice,
would not compromise, and
therefore he suffered to the uttermost. Had his suffering
been a feeble and deedless endurance, would he perhaps have accomplished nothing
great in the world? He too applied the
right of the stronger: he was in truth the Strong One.

If to suffer wrong were only a sign of
weakness and impotence, Christianity would never
have come into being. Christ with his witnesses vouches
that he who suffers wrong may be far stronger
than he who does wrong. Transferring this to popular life,
it is obvious that a people which, in
having to suffer wrong, understands how to transform the
pressure of necessity into a task of freedom, may, even if it is a
small people, attain a great significance in history.

To love right and to hate wrong is the first condition
for attaining a better state of society. But
% --- p. 99 ---
\setcounter{footnote}{0}%
\opage{99}in which people is one now to expect to find the greatest
love of right and the strongest hatred of wrong?
Surely in that people whom the dispensations of Providence have taught
what it means to suffer wrong. Here there is no question of
a feeble and deedless suffering. If it goes so far with
idle endurance and acquiescing in everything that at last one
becomes reconciled with the wrong and thinks in one's quiet
mind: the stronger is after all in the right; we should
have done the same as he, if we had had the
power, then the spark of independence is extinguished, and
the people must die. There is nothing that so rouses
the contempt of the overmighty as fawning, cowardly submissiveness
in those whose right he has trodden under
foot. We learn it from the Romans, we learn it from the history of the
French Caesar.

The people that is brought up to love right by
suffering wrong receives a good education. In all
states with a free constitution there are debates enough about laws
and rights, --- an immense labor of legislation! The debates
sometimes take such a turn that
they are close to driving civilized society to
despair. The machinery of justice labors, groans and puffs;
screws and wheels crack and creak. But the more
people crowd about right and the right, the more crooked and
warped it comes to look; at last it becomes
impossible to bring forth a single work cast whole. Calculated
accommodation, clever agreements, strategic
tricks and tactical maneuvers will not do:
what is inflexibly demanded is an all-pervading sense of
justice, love of right and hatred of wrong. A strong and living sense of right is a thousand
times more worth than ten thousand learned folios of law.
% --- p. 100 ---
\setcounter{footnote}{0}%
\opage{100}Where the first is lacking, the last gives only a
wretched substitute; but when the first is given,
the second will follow as of itself.

If the business of legislation were to depend wholly
on juridical insight into all particulars,
then the whole people through its elected could not possibly secure
for itself any real and living participation in legislation
itself. But when it is not juridical hair-splitting in
small matters, but an ever-wakeful sense of justice, a
sound and open eye for legal relations on the large scale,
on which it essentially depends: then the people can
gradually educate itself to an ever more effective
participation and deeper intervention in all branches
of legislation. Self-education then advances on the following
conditions.

1) The electors must take care that the elected man
can on the whole be said to represent the
mood of the constituency; for the people has an indisputable right to
put its mood into the scales.

2) The elected man is, in all debates in the
legislative assembly, bound to remain true to his character,
the character in which he appeared
before his electors; but he is by no means bound
to vote with a closed party in one and all or
to cling fast to a preconceived, petrified opinion;
he is on the contrary bound, bound on honor and conscience in the holy
name of justice, to seek
enlightenment, to listen to the statements of experts, to take note of the
arguments of opponents, to form an independent conviction
and to vote according to his conviction.

3) The elected man is accountable to the electors.
If these electors have once for all got it
% --- p. 101 ---
\setcounter{footnote}{0}%
\opage{101}into their heads that only that is right which pleases them, and
all that is unpleasant to them wrong: then they constitute
only a crowd, not an organic part of the people.
But if they are receptive to arguments and sincerely wish to know
what right is, he can inform them of the connection of the matter
and thereby put into their hands the means of
passing a free and independent judgment. If they then find that the
elected man has not hit upon what, in their view,
is the useful and right thing, then they have full right
to elect another. In any case a step forward has been taken
on the way of better knowledge.

These propositions, with the truths lying in them,
contain nothing new; but the old thing which constantly
must become new to us is this: how does it stand
with our sense of justice? We, who outwardly must suffer
wrong, have we inwardly learned to love right? Has the injustice
of majority power become so repugnant to us
that in relation to one another we would rather,
if it must be so, ourselves suffer wrong than wrong
our opponents? Such questions are not political;
they are questions of conscience. But politics dare
not exclude conscience; as surely as there is a
popular life, so there is also, and must be, a conscience of the people.


Who then is, in the higher sense, the stronger, he
who proves his power by doing wrong, or he who
has learned to develop his independence by suffering
wrong? The first has placed himself under Fate,
the last stands under Providence. The necessity of Fate
spreads, when wrong must constantly make good
wrong, and what is acquired by violence is possessed by
compulsion. On a broad foundation of the serpent-knots of wrong
% --- p. 102 ---
\setcounter{footnote}{0}%
\opage{102}great state-edifices can indeed be raised, and it
is this semblance of greatness that incessantly dazzles and tempts
the holders of power. But as little as one can pluck
grapes of thorns or figs of thistles, so little can
one reap fruits of right on the field of wrong. Heaps
of old wrong deposit cancerous growths in the body of the state.
The strong constitution endures much, and the organism of society
has many healing outlets; but the substances
of wrong must be removed. One imagines that
wrong can be forgotten, that time can gradually
stop its consequences and put it out of action.
That is a delusion of the holders of power. The elements of wrong
are like a poison, slowly consuming, unceasingly working.
The rule of power imagines that the art of politics
will with time be able to transform the deadly substance into
a wholesome nutriment. Impossible! If there does not come a
healing hand which cuts out the cancerous sore of wrong, then
the disease becomes incurable: the evil passes into the blood.
It is a necessity of Fate that the great work of wrong
must fall, for it bears the seed of its destruction in
itself. Wrong is the destroying, right the sustaining
principle in history. A great body of state, grown together out of violence and wrong,
may by its very growing together at last
become one single great, tangled cancerous knot. The worshippers of Fate
admire the great cancerous knot and extol necessity.
But necessity stands under freedom;
necessity is that the spiritual will must be fulfilled;
every time we see its fulfillment in history, we recognize
a judgment of God.

% --- p. 103 ---
\setcounter{footnote}{0}%
\lecture{Seventh Lecture.} {(Kolding).}
\opage{103}The old Lutheran State Church has, in the sincerity of its heart,
no doubt from the beginning with the very best
intention, fused together two components of the life of faith that belong together
but are nevertheless unlike in kind, namely: faith in
\emph{salvation and faith in morality.} Thereby the Church has
been placed in an ambiguous light.

Faith in salvation is the faith of the congregation, and the Church signifies
the congregation; faith in morality is the faith of the people, and
the Church signifies the People's Church. Faith in both forms
has essentially the same source; for spirit is, seen in its origin,
the same: only the spirit of the congregation is able to
furnish the eternal, fully valid foundation for a true
spirit of the people.

But spirit works in human life in different
ways, with different force in different directions.
``Our citizenship is in heaven''; that is the faith which
spirit awakens, confirms and nourishes in the congregation.
``Our citizenship is on earth''; that is the faith which
spirit awakens, confirms and nourishes in the people, a faith
in the significance of a moral life full of energy for action, a popular
faith, which shows itself in deeds. The faith of the congregation
% --- p. 104 ---
\setcounter{footnote}{0}%
\opage{104}is the original, the faith of the people the derived. The difference
is obvious: we can immerse ourselves in the
original with a clouded, uncertain eye for the derived;
and human beings by the thousand can make use of, employ and work upon
the derived without therefore, in any deeper sense,
troubling themselves about the original.

Spirit is the source of life, all life is from spirit;
therefore the different fundamental chords of existence
are to be brought into consonance in one great, full-toned harmony
of the many souls, the many self-sounding instruments of spirit.
But how different these fundamental chords are!
When the banner waves, the red cloth with
the white cross, and the people sing in a thousand-voiced
chorus: ``King Christian stood by the lofty mast,'' then
something stirs wondrously in the breasts of many; it is the faith
that we are one people: it is like a rustling of the wings of
the spirit of the people that we hear. But the spirit of the congregation
puts the banner of the cross in another light. The faith of the congregation
sees in the Crucified One, the One crowned with thorns, who has
bowed his head, a portent of another struggle, the last
struggle; then it sings another song, to deep organ tones
a death-song: ``Who knows how near my end
is to me?'' These are different fundamental chords; these are
opposite poles of life; is the one of them, then, really
to be destined to blot out the other?

The fundamental law of life is two-sided: to die to the world
and yet to live in the world; we are to live vigorously in order to
die, and to die in order to live vigorously. The two-sidedness
rests on this, that the divine Word points in two directions:
toward the blessed and toward the moral. ``Whosoever
smiteth thee on thy right cheek, turn to him the other
also,'' so spoke Christ when he taught from
% --- p. 105 ---
\setcounter{footnote}{0}%
\opage{105}the Mount; but when he stood before the High Priest, and one
of the servants gave him a blow on the cheek, then he did
not turn the other cheek, but said: ``If I have
spoken ill, prove it, but if I have spoken rightly, why
smitest thou me.'' --- ``Love your enemies, bless
them that curse you, do good to them that
hate you, and pray for them which do you harm and
persecute you;'' that is a condition of salvation, the mind
that can understand this learns that there is Something
which is called blessed. But what people will be able to
understand it as if the law of morality forbade
bearing arms against the enemies of the state and of the fatherland?
The moral can be taken up into the blessed; but no one
will assert that human beings cannot be
moral without being blessed. Here, in spirit and faith,
there is a difference between the blessed and the moral. This
difference must in earthly life, under the conditions
of finitude, develop into a real opposition ---
not contradiction --- between the faith of the congregation and the faith of the people.


To the faith of the people corresponds the People's Church. There stands in
our Constitution in plain words that the Evangelical Lutheran
Church is the People's Church. By this it is by no means meant that
the Evangelical Lutheran faith of the congregation is to be forced upon
the people by law. Had this been the meaning, another
word would have stood there: there would have stood, not \emph{People's Church},
\emph{but State Church.} It was precisely, as recently
adduced, the characteristic of the old State Church
that it fused in an unclear manner the faith of the congregation
and the faith of the people. The government held unshakably fast
to the principle that no one could be a good
citizen of the state, a faithful and obedient subject, unless
% --- p. 106 ---
\setcounter{footnote}{0}%
\opage{106}he were a true, believing Christian, and that no one could
be a true, believing Christian who did not confess
the faith commanded by the state, in accordance with the teaching of Scripture
and of the orthodox Church. It
was not only for the sake of morality, of civil concord
and uniformity in outward matters, that
it applied such severe measures against those of other faith,
no, it was expressly also for the sake of salvation,
for the sake of the souls under age themselves.
The ruler of the state really supposed that he, so far as
the ecclesiastical was concerned, held the conscience of the subjects in his
hand and on the Day of Judgment would have to render account to God
as supreme-commanding shepherd of souls.

How entirely and utterly this State Church point of view
has been abandoned in the Constitution is clearly shown by
§ 76. ``Citizens have the right to unite in congregations
in order to worship God in the manner that agrees with
their conviction, provided that nothing is taught or done
that conflicts with morality and public order.''

Every jot and tittle of this paragraph is a
protest against the State Church. It is not a minority
of those of other faith, some few religious parties,
who simply will not swear by the symbolic books
of the Lutheran Church, and who on that account receive
permission to worship God in their own way, of whom
it here speaks. No, it speaks here not of permission, but
of right, not of the members of the congregation, but of
the citizens, all the citizens without exception. But when
all citizens have the right to worship God in the manner
that agrees with their conviction, then undeniably the
whole people has the right to do the same.
On the other hand, citizens have by no means the right to teach or
% --- p. 107 ---
\setcounter{footnote}{0}%
\opage{107}do anything that conflicts with morality. The separation
between faith in salvation, the particular faith of the congregation,
and faith in morality, the general faith of the people,
is thus expressly recognized. The 3rd paragraph
of the Constitution, which says: ``The Evangelical Lutheran
Church is the Danish People's Church and as
such is supported by the state'' stands in good agreement with this,
for here it speaks not of something that \emph{shall}
be, but of something that is. If one would seek a remnant
of the State Church in this, that the Lutheran Church is supported
by the state: then we must ask why it
alone is to be supported. It cannot be because
the majority of the Danish people are for the time being Lutherans,
for then there would have had to stand \emph{Majority Church};
but now there stands expressly People's Church. It is a
true blessing for land and realm that there stands People's Church;
had there stood Majority Church, then not only the faith of the majority, but also
the will of the majority, the right of the majority and the might of the majority
would at the same time have been solemnly
recognized.

That the Evangelical Lutheran Church is the Danish
People's Church means: the Evangelical Lutheran spirit of the congregation
and the Danish spirit of the people are essentially akin;
the faith of the people and the education corresponding to it
ought, in accordance with the peculiarity of our national nature,
to go predominantly in a Lutheran direction and on average
to take on an Evangelical Lutheran stamp. How this
is to be understood more closely will show itself as we consider
what follows: ``The constitution of the People's Church is regulated
by law.''

Weighty voices are raised against letting the law
by which the constitution of the People's Church is to be regulated proceed
% --- p. 108 ---
\setcounter{footnote}{0}%
\opage{108}from the popular legislature. ``The Rigsdag, which has no confession,''
it is said, ``cannot make laws for that which especially
concerns the Evangelical Lutheran confession of faith.''
We cannot subscribe to this point of view, because we
cannot approve its presuppositions. It is here presupposed
that it is the Evangelical Lutheran affairs of the congregation,
the doctrine, the arrangement of divine service, the administration
of the sacraments, the calling of the shepherds of souls, etc.,
the regulation of which the Constitution speaks. But that the Constitution's
thoughts go in another direction is seen from § 76.
It is here presupposed that the laws are to be laws of compulsion,
not, as in the days of the Apostles, spiritual laws of freedom.
But to want to govern the inner affairs of the cure of souls by
the application of laws of compulsion is surely simply to tread in the
footsteps of the old State Church. Here is finally presupposed
a legislative and governing authority of so purely
ecclesiastical a kind that it takes up an independent position
over against the secular. But if it is already unreasonable
to want to introduce two entirely coordinate governing authorities,
an ecclesiastical one right beside a secular one, in
one and the same state, then it is more than unreasonable, it
is manifest contradiction to want to withdraw from the people the right of legislation,
when that which is to be regulated by law is precisely
the concern of the People's Church, the Church of the people.

There are Others who think that the time has now come
when one ought at last to make a serious matter of separating the Church
from the State. Insofar as it is again the old notion
of a State Church that haunts their brains, we
have nothing to object to this. Religion and the state can
both be served by the divorce: religion is freed
from an oppressive burden, and the state likewise. But
the question here is not of a State Church, but of the
% --- p. 109 ---
\setcounter{footnote}{0}%
\opage{109}People's Church and its peculiar position in relation to the congregation.
What matters is to make the relation between the Church of the congregation
and the People's Church as clear as possible. The Church of the congregation
with the sacraments and the word of salvation has
its definite confession of faith; the People's Church with the word of morality
has no sacraments and professes only
the general principles of the religious-moral life of the people.
It is in good agreement with this
that the light, the life and the power of inwardness with which
the great principles of morality are filled and appear
can give us a reflection of the spirit that stirs
in the congregation. Since it is precisely the People's Church
that contains within itself the means of spirit and the conditions
that form the foundation of the people's instruction
and education as a whole, how will one then be able
to justify separating the People's Church from the people's state.

The constitution of the People's Church can and ought to be regulated by
law; here, on this point, there can be talk of laws of compulsion,
civic laws. To want to apply compulsion to that
which by its nature does not let itself be compelled is a blunder;
but to rage against all compulsion where laws of compulsion
essentially belong is also a blunder.

From old dread of State Church authority, several
have spoken up for an almost boundless freedom
in the People's Church. It is, in their opinion, supposed to stand
in conflict with the principle of freedom that the teachers in the People's Church
are appointed by the state, for thereby they become, one says,
imposed upon the members. It is supposed to be an encroachment
on the rights of the family that children without the consent of their parents
are obliged to attend school and go through a
course of instruction ordered by law, which among other things also
includes religion and so forth. Would not what is
% --- p. 110 ---
\setcounter{footnote}{0}%
\opage{110}unwarranted in such objections become clear if
one begins by asking oneself what sort
of state, what sort of church and what sort
of educational system is here in question? The
power of the state from which the laws proceed does not stand either
outside the people or above the people; on the contrary, according to
the Constitution it is a popular power; its will is a popular will expressed
in laws. It is the people itself that, through
the popular state, by way of legislation, orders its
church constitution.

By the order fixed by law the people secures itself
against the Church being given up as a prey, whether to the changing whims of the multitude
or of individuals. It is not a
self-willed state that sets itself over the people, but it
is the people that by means of the state sets itself over the individuals;
it is the people that, by its legislative authority,
procures for itself security that all individuals in
the whole population are made partakers of the instruction
and education that the faith of the people requires.

In the different classes of society the particular
education and cultivation can naturally be very different
and be attained by very different ways; but there
is Something that must be common to All, and that is the
popular. All children, girls and boys, whatever they otherwise belong to,
families of whatever rank and station,
of whatever confession of faith, are not only to be
brought up and instructed; but they are, so far as the popular
is concerned, expressly to be brought up and instructed in
the way that the popular will prescribes: they are to be brought up
and instructed in the faith of the people. The humane people
may, on the general conditions of morality, grant
the citizens the right ``to unite in congregations and
% --- p. 111 ---
\setcounter{footnote}{0}%
\opage{111}worship God in the manner that agrees with their
conviction;'' but if there is a society whose
confession of faith does not permit such a moral education
as the faith of the people demands, then it is surely
evident that the members of such a society
cannot possibly be counted as belonging to the people. Those who assert that
parents have the right to decide whether their
children shall enjoy any instruction and education
or not, assert at the same time that parents have
the right to decide whether their children shall have
title to enjoy the rights of citizens of the state or
not. Those who grant parents full power to
decide whether their children shall be educated in the faith of the people
or not, must at the same time grant them the right to decide
whether the children shall be counted as belonging to the people or not.
But when one grants the individuals of the population such
a right over against the popular will: where then is
the authority of the people, what then is the people?

That the People's Church, as it exists at present,
is an ambiguous being can be seen most plainly
in the in many respects difficult position in which
the servants of the Word find themselves. Christ says: ``No man
can serve two masters.'' But as long as the constitution of the Church
is still as it is, the servants of the Word must
admit to themselves that in reality what is demanded of them
is the unreasonable, that they should serve two masters.
They are to be Christ's servants, with authority from
Christ's authority, spiritual authority; and they are at the
same time to be the state's servants, for they are indeed the state's
officials and have in this capacity all their authority
from the state. How difficult such a double service
is, we learn on every page of the history of the old
% --- p. 112 ---
\setcounter{footnote}{0}%
\opage{112}State Church. Only when the opposition between
the Church of the congregation and the People's Church comes clearly forward
will the two-sided relation of service be freed from all
ambiguity.

In the congregation the servants of the Word are unconditionally Christ's
servants; for this service they stand as shepherds of souls
accountable to Christ and the congregation: here there applies
no interference of authority in the law of the state. In the People's Church
the servants of the Word are, in a religious-moral spirit, servants of
the popular, and must in this capacity be subject to
the popular authority of the state. In both services
they are workers, and a worker is worth his wages;
but the work of the shepherd of souls cannot be paid in the same
way as that of the teacher of the people.

``When you go out,'' says Christ to his disciples,
``then preach and say that the kingdom of heaven is come
nigh. Heal the sick, cleanse the lepers, raise
the dead, cast out devils. \emph{You have received it}
\emph{for nothing, give it for nothing}. You shall not have
gold, nor silver, nor brass in your girdles; nor a bag
for your journey, neither two coats, neither shoes, neither
a staff; for a worker is worth his food.'' Here
the talk is not of general popular instruction; here
it speaks expressly of the proclamation of the word of salvation,
of the soul-saving activity. Everything is given for nothing,
for it is given out of love, and love is always given
for nothing: it is a commodity that can be neither sold nor
bought! On the other hand it can and ought to be presupposed that where
the word of love finds entrance into the hearts of human beings,
those who receive the spiritual goods will be willing
to relieve the temporal need of those who for
nothing impart to them these goods, in that they acknowledge that
% --- p. 113 ---
\setcounter{footnote}{0}%
\opage{113}``a worker is worth his food.'' Such a relation
cannot be ``regulated by law''; the work of the shepherd of souls is
no piecework: he is obliged to give All
for nothing, but then he is also entitled to receive
voluntary gifts.

Otherwise with the position of the teachers in the People's Church.
The teacher of the people performs work commanded by law in return for remuneration
in a salary fixed for that purpose. He is obliged to
perform the work and he is entitled to demand his
salary. Just as the authority of the people as a whole will know
how to damp the arbitrariness of individuals, so will the
popular legislation also know how to secure for the servants of the popular
word the salary that such a service is
worth, and to place the religious-moral educators so freely
that they do not in an unseemly manner become dependent on
individuals who ought most of all to feel themselves in free dependence
on the educating word.

How little one has as yet thought of transforming
the State Church into the People's Church can be seen from the provisions
that have been made on the part of legislation concerning the establishment
of free congregations. What now forms the opposite
to the free congregations? The unfree congregations,
the congregations in the People's Church supported out of the means of the state.
The unfree congregations thus belong to the rule,
the free ones to the exception; an unhappy exception!
If there is any society in the world that has an original
claim to freedom, then it is surely the congregation of Christ;
its freedom is not usurped, it is dearly paid for: he
bought it with his blood. That the congregations of the People's Church
are unfree proves in fact that they still find
themselves in the old State Church condition.

The new arrangement with the separation between
% --- p. 114 ---
\setcounter{footnote}{0}%
\opage{114}free and unfree congregations is especially characteristic in
the manner in which it is defended. That a human
authority has taken charge of the divine
things and arranged the constitution of the Church so that everything
happens by virtue of a supreme command, that the word of salvation
is commanded by law, the sacraments commanded by law, the general
prayers commanded by law, etc., all this, for the sake of order,
Commanded by law --- one says --- makes absolutely no
encroachment on Christian freedom. And why not?
Because every human being who is not content with the
order commanded by law has surely the right to step out and
seek, if he will, his edification in a free congregation
or another society not commanded by law. So one comforts
oneself with the thought that faith is free, that we in reality
have a completely carried-out freedom of religion. There
is in this way of regarding things not so little
of a self-deception, after all.

Insofar as the talk here is of the individuals' freedom of choice,
the freedom is certainly so extensive that no one has
anything to complain of. But let one take careful note
that it is not the freedom of individuals that is in question,
but the freedom of the congregation, the Church of the congregation's own
freedom, how does it stand with that? Or is it supposed that
what we call Christian freedom has nothing at all to do
with the word of salvation itself and its proclamation, not with the sacraments
and their administration? Can the spirit of the congregation
perhaps work best of all as the spirit of freedom,
when the means through which it works are all unfree?
Can the congregation hear freely, when it knows that the shepherd of souls
dare not speak freely? Can any spiritual
discipline (church discipline) be exercised by shepherd of souls and congregation
together, when the co-operation of both can at any moment be
% --- p. 115 ---
\setcounter{footnote}{0}%
\opage{115}broken through by an authority intervening from outside?
These questions have not been answered at all,
because, so far as we can see, one has not yet put
them to oneself with clear consciousness.

The free congregations, which were to remedy the
total lack of spiritual freedom from which the State Church
suffers, are, as it seems, likewise not in a position to
solve the task. Free congregations may well form
which bind themselves beforehand juridically to
pay their shepherds of souls; but there is hardly any
free congregation in which the shepherd of souls can essentially be said
to fulfil Christ's command: ``You have received
it for nothing, give it in turn for nothing.'' Those who in
One and All aim at fulfilling the apostolic
freedom must also in One and All submit to the
apostolic conditions of life.

In order that that which cannot be subjected to compulsion, and
that which by its nature cannot do without laws of compulsion,
may come into the right relation to one another,
the mutual interaction of the Church of the congregation and the People's Church, of the faith of the congregation
and the faith of the people must of necessity be
made clear. That the Church of the congregation can subsist
by itself without a People's Church is evident in itself: the Apostolic Church
with all its popular stamp was no People's Church,
it was in the deepest sense of the word a Church of the congregation.
The People's Church can for its part work by itself
without immediate connection with the congregation; for
although, in order to be what it is especially to be
among us, it must be permeated by the Evangelical Lutheran spirit,
and must have a general stamp of an Evangelical Lutheran
life of the congregation: yet it is without definite confession of faith and without sacraments. The essential
% --- p. 116 ---
\setcounter{footnote}{0}%
\opage{116}separation is a condition for their being able to work together
in good understanding and peaceful unity. The congregation
with the word of salvation and the sacraments contains what is
particular and original; the great popular society of education
with the word of morality, the instruction of youth
etc. makes do with what is general and derived.
If now the union and interaction of both is conceived in the light of the faith of the people
from the popular side, then we have what
the Constitution understands by the Evangelical Lutheran
People's Church.

In order to be able to fulfil their office as servants of the Word in
such a People's Church, the young men who
feel an inner calling to it must go through a preparatory school that
corresponds to the preparation in both directions. As prospective
shepherds of souls they must make themselves familiar with the
spirit and doctrine of the Evangelical Lutheran congregation; as prospective
teachers of the people they must make themselves familiar with the
popular conditions, with the historical-moral way of thinking of the people,
in short: with the conditions of life under which the spirit of the people
works. As teachers of the people they are installed by the state authority;
as shepherds of souls they are chosen by the congregation
and receive the apostolic consecration. The strife
that could arise on occasion of the double election
and the double installation is in principle settled
by this, that the religious-moral teacher of the people surely may very
well work without being a shepherd of souls, and that the election
of a shepherd of souls is so free that the congregation, if
in a quite peculiar case it should feel the need of it,
can, without offending the appointed teacher of the people or
giving offence to the state authority, choose another
as its shepherd of souls.

As regards oversight, it is a matter of course that it
% --- p. 117 ---
\setcounter{footnote}{0}%
\opage{117}also must become two-sided. The teacher of the people stands,
like every other official, under the oversight of
the governmental authority and his capacity for work is assessed
in proportion to the spiritual power with which he
works in a moral-popular direction upon the population of the parish.
The shepherd of souls stands under the oversight of such persons
as have insight and spiritual authority to clarify the
religious tasks of life, develop the peculiarities of
the Evangelical Lutheran doctrine and point out the deviations
from it. How a governmental authority chooses its
overseers, and what the effects of the oversight can be,
is a well-known matter. How the congregations in the
apostolically free time got their overseers (bishops
and elders) chosen, and what effects spiritual authority
exercised, is likewise well known.

There prevails with Many a fear that a spiritual
oversight without legal ordering and juridical means of compulsion
will become indefinite and unfruitful, that confusion
in doctrine and a corresponding disturbance of the life of faith
cannot be avoided. It is an untimely fear. To be sure
it must be conceded that many of the deviations which
the State Church in its time stamped and punished as heresies
may come forward anew. But if these deviations are
dangerous, they can surely be brought to light; and when they are brought to light,
every congregation surely has the right to call its
shepherd of souls to account. Such storms of heresy
as shook the old State Churches are not to be
feared in more recent times; for the conception of \emph{dogma}
has on the whole undergone an essential change.

Dogma is a proposition of faith expressed in definite form.
In older days, when orthodoxy was juridical, one held
so exactly, so pettily to the form and literal expression of each single
% --- p. 118 ---
\setcounter{footnote}{0}%
\opage{118}tenet, that anyone who dared to speak freely and livingly out of a religiously moved
mind exposed himself to being accused of heresies
of which he had scarcely dreamed. Here there was in the
rattling off of heresies Enough to inflame the
passions, goad subtlety and increase the harassments.
To free itself from the dead formalism
of Orthodoxy, Pietism wanted to appropriate everything with feeling.
But as feeling, so to speak, softened the dogmas
to a mush, the doctrine of faith became jelly-like, and the whole flowed away into
foggy indefiniteness. Both extremes are avoided
when we consider that each single form of proposition is indeed
a definite but at the same time one-sided expression of the content
of faith, that the content is always richer than the limited
form, that the forms of expression must change when the richness of the content
is to be fully unfolded. It is not
dogmatic indifferentism, but a more thorough conception
of the nature of the truths of faith, on which we, by a rejection
of that untimely fear, can base our assumption.


What has here been set forth is not to be regarded
as a practical proposal for a transformation of the existing
church constitution; it is no proposal at all, but
an elucidation which, to the exclusion of every consideration of
practical difficulties, is aimed directly at clarifying
the principles according to which the transformation in the course of time
must of necessity take place. That in practical respects
there are difficulties to be overcome cannot count as
an objection to the view here set forth.

In the judging of this view only the principles themselves
are to be taken into consideration. If it can be demonstrated that the standpoints set up
% --- p. 119 ---
\setcounter{footnote}{0}%
\opage{119}are either wholly or in part mistaken, then naturally the
principles are false, and the whole collapses
of itself. But should it on closer
examination be found that the fundamental relation between the particular
faith of the congregation and the general faith of the people has been
conceived and determined in such a manner that the principal propositions in question follow from it and correspond to the nature
of the matter, then we have, even if in fleeting outline, a firm
foundation. On a firm foundation, in accord with truth,
the church constitution must be ordered. For surely there is
no one who will assert that in a matter
which with right can be called the highest and holiest matter of truth
of human life, and
thus also of the life of the people, there can be practical considerations that make it
permissible to rest content with the half-true or
even with the untrue.

% --- p. 120 ---
\setcounter{footnote}{0}%
\lecture{Eighth Lecture.} {(Kolding).}
\opage{120}We have spoken of the spirit of the people, the faith of the people and the teachers of the people
in general, but how does all this look
when it is carried over into reality: what, then, is
the task of the teacher of the people, properly speaking? What is it that makes
instruction popular?

For the enlightenment of a people many kinds of teachers are needed;
just as we, indeed, also have many kinds of schools: schools for the common people,
burgher schools, folk high schools, modern schools, agricultural colleges,
military schools, veterinary schools, navigation schools,
a polytechnic institution, Latin schools, the University,
besides other schools such as: commercial academies,
trade schools, drawing schools, singing schools and, then, also
--- dancing schools! Now, are all these schools popular,
or is it only some of them that are popular,
while others are not popular, indeed a few perhaps
even unpopular?

That the population as a whole has need of many
different schools is surely seen already from the fact that such
schools are in operation; to found a school and keep it
going costs money and labor: such a thing one does not do
for nothing. In order now to consider which and how many of
% --- p. 121 ---
\setcounter{footnote}{0}%
\opage{121}these schools can be said to be for the good of
the people as a people, we dare not confine ourselves to narrow
regard for the barest necessities, or even for what is
serviceable and useful in the everyday sense;
we must at the same time look
to the nobler, the spiritual; indeed, when
our view is free, we can only most unwillingly persuade ourselves to exclude
the pleasurable. From a strictly economic
standpoint it might well seem as if the singing schools in particular,
and especially the dancing schools, must belong to the articles of luxury.
But no! What the singing schools in recent
times have accomplished, everyone who has a popular
sense will certainly appreciate with gratitude. The songs of the poets delight
us most of all, after all, when they are sung, that is to say:
when they are sung purely and in the right tone; for when
the tone is false, it grates upon our ears. A naturally
good voice gains much by going through a
singing school. The same can be said of our natural
movements, especially when we hop, leap and dance.
The bare naturalness is not enough; a little
schooling is needed. If it becomes an enlightened people to sing
in an enlightened way, then it also suits it well to dance in an enlightened way.

But, one might object, one thing is instruction
imparted in differing measure and with differing aim
to the different individuals in the population, another thing is
popular instruction. To this it must be answered that that
which in every branch of instruction, and in particular in every
school, essentially grounds the distinction between the popular
and the not-popular depends not so much on the subject matter
as on the manner in which, and the spirit
in which, the instruction is conducted. In my judgment both
the popular and the unpopular, in so far as
there can be any question in instruction of forming
% --- p. 122 ---
\setcounter{footnote}{0}%
\opage{122}and developing a definite way of thinking, can thrive in all schools,
in some more immediately, in others more mediately.

According to an opinion which preeminently calls itself
popular, and maintains that it has the right view of what
the people and the popular are, the matter is nevertheless said
to stand otherwise. There are then said to be in this country
two great schools, of which the one is straightforwardly
popular, the other, to put it mildly, not-popular, indeed on the
verge of becoming unpopular, namely: \emph{the Folk High School
and the University}. That the University cannot
be absorbed straightforwardly into the popular, as the
Folk High School is, shall willingly be granted; but why the
scientific high school should absolutely be unpopular,
we do not see. It is maintained that it has a shadow,
a false shadow, which no longer corresponds to the side
from which the light comes, an uncanny black shadow, which
blinds the spirit and darkens the people; this sinister, corpse-like
shadow is --- Latin!

In order not to hasten into defending something
that cannot be defended, we shall proceed soberly and
thoroughly. We shall not begin by excusing
the University and the Latin schools, for that might easily acquire
a taint of one-sided partiality. No, we shall begin
the other way round; we shall begin by collecting, bringing forward
and pointing out all the evil that has partly been said and
can still be said about Latin. We shall not let ourselves be
frightened by a shadow, but boldly attack the Latin
``bogey.'' Much can be said against Latin,
a great deal.

In the first place Latin is a dead language. To
love a dead language more than a living one, especially the mother
% --- p. 123 ---
\setcounter{footnote}{0}%
\opage{123}tongue, to love it so highly that one not only reads
it, but writes exercises both when one speaks it, and when
one writes it, is to love the dead. But to love
the dead is possible only for the stillborn or rather
for ``death-biters'': ``all Latinists are death-biters.'' How
uncanny, when life is bright and the heart glad, to meet
on one's way a stillborn, sullen Latinist with a face
like a grammar; one is close to losing one's
childhood faith. The only Danish Latinist who really has
a claim to indulgence and absolution is Saxo, Absalon's
clerk; therefore not only his history but also
his name has been transplanted into Danish: he is called
no longer Saxo Grammaticus, but Saxe Sagemand;
which is very popular.

In the second place it is a well-known matter that Latin
corrupts the blood of the people. The Latin blood is originally
mixed with wolf's blood. Romulus and Remus, these
two first great Latinists of the world, did, according to the express
testimony of history, suckle at a she-wolf.
The marrow and juices of the noble stocks, of the Goths and Normans
in particular, would have kept themselves far purer,
and the peculiarity of the tongues would have stood far
clearer to us, if the old Latinism had been choked in
the migration of the peoples; but weeds do not perish so easily.
The Latin blood of culture made culture wolfish at
the same time as it deposited dead flesh in the hearts.

With us in the Middle Ages the whole literature, with
few exceptions, was Latin, the psalms in Latin, the whole
service in Latin. Most of what people learned
was Latin, and therefore it was at bottom a good thing that they learned
nothing. The ballads of the heroes, the rhymed chronicle and the Euphemia ballads
fortunately kept themselves free of Latin: the good
% --- p. 124 ---
\setcounter{footnote}{0}%
\opage{124}nature of the people resisted the poisoning. The clergy were indeed
Latinists, to be sure, but the infection did not spread widely;
they kept, so to speak, the blood to themselves: they were
not allowed, after all, to marry. Added to this, the contagion
was not so potent, either, as one might have feared,
for it was bad Latin, it was monk's Latin. More dangerous
it looked at the Reformation, since Latinism
too was reformed. Then the corrupted monkish muddle
was cast away, and the genuine, old, vigorous Latin
from the golden days of the Roman empire was zealously fetched up
from the classical cellar. Zwingli and Calvin took deep
drafts, and Melanchthon himself was not free from getting
a little tipsy. Luther kept sober; but then
his marriage, however highly we otherwise must praise it,
nearly occasioned a great misfortune. Let one consider well
what it means that the priests, who now, owing to
the improved Latin, got their blood in earnest corrupted
in the black school, married and thereby made
the corruption hereditary.

The famous naturalist Darwin has, as regards
the hereditary, discovered something he calls atavism.
There is born in a family a child with, for example, a peculiar
wart on its nose. People wonder and ask:
after whom can the little one have that wart? In
neither of the parents is there found the slightest hint of such
a protuberance; the grandparents, on the father's
as well as on the mother's side, are examined and are likewise found
to have normal, wart-free noses. And yet the wart is
inherited. There is, when one goes further back, a
great-grandfather or a great-grandmother who has had the wart;
heredity has skipped some intermediate links: it
is atavism.

% --- p. 125 ---
\setcounter{footnote}{0}%
\opage{125}This now finds a fine and rich and most fruitful
application to Latinism. There is born in a family, let
us suppose it, besides the healthy children, a pair of individuals
with weak nerves, headache, scrofula, dryness of
the heart and an extremely dull eye for the popular: whence
does that derive? From Latin! That the mother cannot
Latin and the father cannot either, makes no difference to the matter.
If there is not a grandfather, then there is a great-grandfather
or a great-great or a great-great-great-grandfather who has read
himself stupid in the black school. By going back far
enough one is sure in the end to run into Latin.

Where the Latin infusion has taken place in large
portions, it produces autocratic palpitation, fermentation
in the blood and a violent politics. In France, Italy,
Spain, in general among the Romance peoples,
the blood easily comes to the boil, and the parties clash violently
with one another. The parties! Yes, here is indeed a
peculiarity which at first glance must seem to
speak against our assumption. If it were really the
Latin blood that drove so strongly in politics, then
our men of the Left, who are known not to be Latinists,
would have to stand far behind the men of the Right, for the men of the Right
are, alas, stiff Latinists. But that the matter stands
the other way round cannot be denied; it is not the men of the Left,
but on the contrary the men of the Right, who breathe
with difficulty and are near to sailing astern. None the
less the thing has its correctness, when we look more closely,
and do not forget that it depends on the portions
and on atavism. In older days Latin was taken
in large portions; at the present time the
portions grow ever smaller, indeed often so small that they
approach a minuteness that recalls the
% --- p. 126 ---
\setcounter{footnote}{0}%
\opage{126}homeopaths' invisible globules. Now one sees at once how
the whole thing hangs together. The men of the Right, who
have received the infusions in small portions, are distinguished
by a certain coolness, circumspect caution, sober
reserve, and a peculiar squinting of the eyes, which is an
unmistakable sign that they often have difficulty
in catching sight of the popular. The men of the Left, who
have received no portions, are as a rule fresher, with
growing courage and a large eye upon the people. That
something violent sometimes comes over them, that they too,
as Latin puts it, now and then seem to suffer from
something human --- humanum quid patiuntur --- is
not their fault; it is the Latinism in their blood that
sweeps them along. It is not their own Latin, no, it
is that of their forefathers, which according to the law of atavism breaks
out in them, and that now and then with such strong
jolts that even at this sight we must sigh and say:
``all our misfortune comes from Latin.''

It is a great confirmation of the validity
of a hypothesis that what at first glance seems to
speak against it is found on closer consideration to
speak for it. We can then confidently go further and
put to ourselves the serious question whether the
great confusion in the monetary system, the false speculations
and dreadful bankruptcies of which one reads in
the newspapers, are not, at least for a great part, to be found
to have their first source and last root in
Latin. Horace says: \emph{virtus post nummos}, ``Money
first, then virtue may come after''; that is
classical Latin. Be that as it may with the bankruptcies;
so much is certain, that it is going backward with betrothals
% --- p. 127 ---
\setcounter{footnote}{0}%
\opage{127}and that this is the fault of Latin can be proved as
clearly as that two and two are four.

It is a notorious matter, a fact confirmed by many and reliable
experiences, that dead
Latin makes the heart cold; but when the heart is cold,
how then shall one be able to love? Hagbarth and
Signe, Axel and Valborg, Other Ebbesen and Sigrid with
the veil, such betrothals are no longer found in the North ---
unless it be as exceptions, by virtue of a quite
peculiar Old Norse atavism --- instead there are found parties of reason,
parties of money and, what can almost
make the hair stand on one's head, newspaper courtships
in abundance. Whence this derives is easy to discover;
it derives especially from the wretched, more and more
prevailing habit of wooing in writing; but this
habit, or rather bad habit, derives in turn from Latin.
What is it that Latin despises? The living word.
What is it that Latin esteems highly? Writing exercises. But
to woo in writing is precisely to write an exercise. By putting the
two things together one will be able to understand that a Nordic
human being with a Latin heart must let that moment
in life, in which the living word and, in a double sense,
the warm oral communication is and remains the highest,
pass him by, cold and indifferent.

A cold-minded, unloving heart cannot love
the people, a Latin heart can never possibly either
be or become a genuine Danish popular heart; that is
in sum all the evil that I at the moment know
to say about Latin.

So, back to the question: is the University
unpopular because the thoughts of antiquity and the language of antiquity
constitute an essential component of the foundation upon
% --- p. 128 ---
\setcounter{footnote}{0}%
\opage{128}which it is built? Can comprehensive investigations
and penetrating expertise be supposed to weaken
sound human understanding and cloud the natural
sense? Is it perhaps a matter of course that sustained
practice in connected thinking, many-sided apprehension
of the conditions of existence, thorough understanding of
the connection of the present with the past, careful distinguishing
between what is known and what is not known, etc.
should be a hindrance to a correct judgment of
what goes on before the eyes of all and shows itself to the
immediate sight? The teacher of the people who either directly
or by roundabout ways spreads such false notions
among the people leads the people astray. That there can unfortunately
be devotees of science, or rather learned busybodies,
who busy themselves with the past to such a degree that they at last
forget the present, is undeniable; but
it is the duty of the teacher of the people to point out
that the fault lies with the persons, not with science.
Those who maintain that strict science has nothing
to do with real life are distinguished
precisely not by clear sight of the real. It is
indeed precisely strict science that has taught us to
master nature through knowledge of nature's forces
and laws. Not only the steamships, the railways and
the telegraph wires, but even our implements of agriculture and
labor, the mixing of substances for the benefit of
the soil, the treatment of plants and animals, etc.,
remind us every day and every hour of the day at
every step we take, of science's wonderful
intervention in our natural existence. If anyone wishes to express
himself with the unprejudiced opinion that it is only the
useful in science, not science itself, that a people
% --- p. 129 ---
\setcounter{footnote}{0}%
\opage{129}must prize, then it is the duty of the teacher of the people to
enlighten him as to his error. It is surely just
as impossible to draw benefit from science when one
withdraws from its cultivators the necessary means, as it
is impossible to have benefit from a domestic animal when one
neglects its keeping and care. Parceling out, quibbling and
cavilling with a separation between the theoretical and
the practical in science have no home anywhere:
every science is, by its nature, theoretical.

On the other hand there is a difference, an essential and
thoroughgoing difference between true and false, if one
likes, lost sciences. Just as one has in real
life various things that are called false
--- false curls, false teeth, false hips
and the like --- so there is still to be found in the
circle of the sciences something that must be called false. To
consult the dead (spiritism), to perform miracles
or, what is about the same, to explain
miracles, to fathom the divine mysteries and comprehend
the Holy Trinity, such things belong under the
false sciences. It is the duty of the teacher of the people to
enlighten the people that there is something false which
calls itself theology, a humbug science, which weakens
and undermines the original sense for what is original in religion
and thereby harms the popular.

The inward connection between the scientific
and the popular shows itself especially when
the question is of our own people's share in the advancement of civilization
and science. In this respect we have
not only kept pace with the times; but we have a relatively considerable number of men
who stand on the topmost step and are honored by all enlightened
% --- p. 130 ---
\setcounter{footnote}{0}%
\opage{130}nations as science's foremost representatives:
a Niels Steensen, a Tycho Brahe, an Ole
Rømer, a Hans Christian Ørsted, these are names that
mean something in science. There hardly appears a larger work on the history of the sciences, be
it in Europe or in America, without these names
being found in it and shining like stars of the first magnitude;
they are names that cast luster over people and fatherland.
Such things ought to be emphasized, not to give national
vanity nourishment, but to awaken
popular self-esteem and to impress a corresponding
obligation.

In the age of enlightenment it is important that
a people keep its horizon open on all sides, not
shut itself up with its own notions and
become infatuated with what is immediately home-made; but it is
surely also important that a people, especially a small
people, not lose itself in empty imitations of what is
foreign, that it have confidence in itself, learn to
rely upon its spiritual leaders and appropriate what is
foreign in such a way that it thereby obtains conditions
for a richer, surer and more many-sided development of its
own. It is its own outstanding representatives,
poets, artists, thinkers, orators, natural scientists,
philologists, historians, etc., by whom a people
is to know its endowment. It is the duty of the teacher of the people
to point out that these gifted ones with their
intelligence do not stand outside the people, but belong in the very strictest
sense to the people, are the people's own organs.
Their knowledge is the people's knowledge, their intelligence
is the people's intelligence; it is with their eyes that the people
sees what it could not otherwise see; it is in
% --- p. 131 ---
\setcounter{footnote}{0}%
\opage{131}their thoughts that it recognizes its own thoughts and
can fully work along in the world of thoughts. Through
them science is the people's science, its clear
eye the people's eye. Therefore one shall teach the common
man that those who disparage intelligence scorn
the intelligent people. It shall be said loudly and freely that everyone
who attacks the highest scientific institution of learning
in the country is, if he otherwise knows what he
is doing, engaged in putting out an eye of the people and
stealing a jewel from the crown of popular life.

And one thing more. There is among the common people a
feeling that science is worldly, that it nourishes
unbelief and leads to contempt for religion. Some
would see in this a good sign, an omen of deliverance from
old oppressive prejudices, a promise that
the human spirit is at last becoming of age and coming
independently to itself. But most are agreed
that the growing, ever more widely spreading ``free-thinking''
is dangerous, that one cannot for that reason do otherwise
than harbor a certain distrust toward science. It is the duty of the teacher of the people
to throw light on this confusion and to do so
in such a way that the common people can follow
and form an independent judgment. Certainly it
would be impractical and unhealthy to burden the people with
scientific disputations and entangle it in learned
controversies, but that is not needed either. There is
something that can be made directly evident to
sound sense, and that is sufficient.

Religion and science each have their own aim and
their particular tasks. Religion addresses itself to
the conscience of man; it will awaken us to \emph{self-knowledge},
develop and clarify our \emph{knowledge of God}.
% --- p. 132 ---
\setcounter{footnote}{0}%
\opage{132}Science addresses itself to man's consciousness
of the visible things, of their essence and connection;
it will awaken us to \emph{general thinking for ourselves},
develop and clarify our \emph{knowledge of the world}. If one
now tries in the name of religion to anticipate the knowledge of the world
and to lord it over science, then one commits
an encroachment; if one tries in the name of science to
anticipate the knowledge of God and to lord it over religion,
then one likewise commits an encroachment. Here a distinction must be made
between religion itself and a false application
of the religious, and likewise between science itself
and a false application of the scientific. What
matters here chiefly is a sufficient
separation between religious truth and
the religious mode of representation, between scientific
truth and the scientific mode of representation.

According to the teaching of the Bible the world was created in six days.
The religious man who does not know science holds
in all simplicity to the biblical mode of representation,
and does well therein; for every feature of the grand
presentation makes vivid the creating omnipotence,
the creating wisdom, the aim of the creating will.
The scientific investigator takes real nature
as it is; he looks away from the biblical mode of representation,
and does well therein; for he must indeed look
at things with an unbiased gaze. So he comes
to the result that everything develops naturally and therefore
slowly; that for the cooling of the earth there must
have passed millions of years before living
beings could appear, etc. If one now confuses
the mode of representation with the truth, the confusion is
boundless. The religious hold fast to the notion
% --- p. 133 ---
\setcounter{footnote}{0}%
\opage{133}of the six days, because they hold fast to the notion
of the creation of the world. The scientific reject
the notion of the creation of the world, because they reject
the notion of the six days. From both sides
the truth is forgotten. The scientific truth, that a
development of nature may require millions of years, does not stand
in conflict with the religious truth, that real,
developing nature has its origin from the creating
spirit.

There is a series of boundary determinations concerning
religion and science which, for the sake of the faith of the people, popular life
and the spirit of the people, must necessarily stand
clear to the consciousness of the whole people. It is not the intricate
investigations, but the distinctly prominent
main points, on which it depends. Also from
this side it will be seen what significance an
institution of learning such as the University can have for popular life
as a whole. The enlightenment to which the common people
have a claim, and which they sorely need, if they
are really to come along and make their voice count
in the highest affairs of society, is of a quite
peculiar kind. Simple catechism instruction in
connection with some meager knowledge and a
number of practical rules of prudence does not suffice; the eye of the people
must little by little accustom itself to see the
real conditions of the world in the light of the present; the people must,
in other words, come to know the conditions of culture. In order that
this may happen, it is not enough that a teacher of the people
enter into the people's way of thinking and speak intelligibly
and comprehensibly. Let his speech be ever so plain, intelligible
and simple: it must nevertheless, as a whole and in
detail, have the color and gleam of the light of truth
% --- p. 134 ---
\setcounter{footnote}{0}%
\opage{134}which proceeds from the realm of clear thoughts; the teacher of the people
must be a man educated on all sides in science, he
must himself be a living link between the popular
and the scientific.

If we now turn to the folk high schools, I mean the
Grundtvigian ones, we have here a grand, indeed unique
example of how original, clear and deeply penetrating
a spiritual view of the people's needs and of the popular
itself can be. The ``enlightenment of life'' for which
Grundtvig fought and which he at last got put into effect
is in all main features, so far as we can judge, so aptly
planned and so peculiarly ordered that with it it can in truth
be said that a new beginning has been made. The ear is opened,
the mind awakened; mythology, poetry, the saga and history
unlock their hiding places and offer their rich treasures,
not in order to burden the memory and one-sidedly exercise the understanding,
but to open the inner sense, illumine
human life and furnish a measure of what is
small and what is great in existence. It is here
not the task either to extol or to defend the Grundtvigian
high schools; they defend themselves. However much
groping and unclarity may cling to such a
first, new beginning, it is nevertheless a fact
that these schools in the different parts of the country --- not
least in this region --- have exercised so awakening,
so educating and enlivening an influence that no other
institutions of instruction can show anything corresponding.

But just because we place so great a hope in the
steady progress and growing influence of the folk high schools,
just because we regard it as irreparable harm if
these schools should miss their task and disappoint
expectations, just for that reason it is our duty to call
% --- p. 135 ---
\setcounter{footnote}{0}%
\opage{135}attention to the fact that even these schools, with all their
popularity, can lead the people astray and thus become
unpopular.

To fall in love with our heathen antiquity and our Catholic-Romantic
Middle Ages in such a way that one
thereby neglects to make clear to oneself an understanding of the conditions
of the present, is unhealthy. A people of action cannot
live on mere memories, be they ever so glorious,
ever so great; still less can an enlightened people wish
to make an attempt to screw the present back to antiquity.
For the common people there is something appealing in looking
back upon the great age of the peasants. Thorgny the Lawman's
speech to Olaf Skötkonung is sometimes drawn
forth as an example of how a free people is to
conduct itself toward a free king. This picture,
however characteristic it may be of that time, is dazzling
and misleading when it is set up as a model
for the present time. Let one note well
that Thorgny's bold words are precisely indicative
of the challenging self-confidence with which the peasantry
thought it could rely on the old forces and the
old rights, without a notion that the two powers of culture,
the Church and the State, which were destined to
carry history further, were also destined to subdue
the defiant common people, indeed to press them down into a bondage of almost
eight hundred years. Thorgny the Lawman does not see
what two wondrously dark thunderclouds
are already rising and gathering over his
head.

We praise the time of the Valdemars; we look up to Bishop
Absalon and honor him as a man of the realm and of the people.
But it is well known that the time of the Valdemars was
% --- p. 136 ---
\setcounter{footnote}{0}%
\opage{136}no age of the peasants: Bishop Absalon was not the peasants'
friend; and yet he was, according to the conditions of the time, in the
true sense of the word, the people's friend. How little ``the common
man'' in Absalon's time was fitted to have a say
in war and peace, may be seen among other things from
the scruples and evasions that prevented King
Valdemar from carrying out his first intention: at once to attack
the Wends. Absalon's answer on this occasion is
known; in his quiet mind he was embittered that
the King had been so unwise as to take counsel
with the common man about such a matter. Saxo is entirely of
Absalon's opinion; ``for,'' he says, ``war shall be decided
with the counsel of a few wise men and carried out with the common
man's help and strength.''

The history of the Valdemars teaches us that the rule of the nobility,
no more than the hierarchy, was a product of
chance or of the arbitrary lust for power of the great families.
The oppression and steady sinking of the common people
was deeply grounded in the nature of the conditions, a sorrowful
but necessary consequence of the fact that the unenlightened and therefore
minor common people were not able to grasp the high
powers of culture, to secure their independence by using
the right means and to follow the course of civilization.

To give the appearance that the common man's view and
the common man's understanding, left to themselves, is capable
of mastering the conditions of culture, of organizing society
after its own head and of standing alone at the
helm, is to deceive the common man and to prepare for him
bitter disappointments. If the common people are to have a voice in counsel and
give their vote weight, then they must gradually learn
to see through the resistance which they have hitherto had
in their unacquaintance with the powers of the spirit whose
% --- p. 137 ---
\setcounter{footnote}{0}%
\opage{137}influence it has all too easily thought it could overlook.
We wish the great age of the peasants back, but, mark well,
not on the conditions of antiquity, no, on the conditions of the present.

There is a piece of the history of the fatherland that does not at all
look popular; it runs from the introduction of absolute power
to the loosening of the adscription to the soil. One cannot tell
it as one tells an old saga or
reproduces a scene from Saxo or from Snorre; and yet
it contains a section whose right illumination is of the highest
importance when the common people are to have their eyes opened to
the genesis of the social conditions in which we live.
Moods, imagery and lively descriptions are of
inestimable value for popular awakening; but in
reality the bare awakening is not enough. The conditions of the present
are intricate; insight and
expert knowledge are needed to unravel the web and untie the knots.
To go into particulars and attempt a real solution
cannot of course be the task of the teacher of the people.
But it can and must be his task to make the
common people aware that the knots are once and for all
there, and that they can no more be untied in
imagery than they can be cut through with a blunt axe.

If it is to succeed in our country, what is striven for everywhere
but nowhere yet attained, to carry through
a real education of the people, popular enlightenment and
popular cultivation, then there must be a field where
the University and the Folk High School faithfully work together.
Such a cooperation will not efface the difference
of the character of the two schools; but it will --- in
the popular sense --- efface the old uncanny
boundary line between learned and unlearned, common and
not-common, intelligent and not-intelligent, which wounds
% --- p. 138 ---
\setcounter{footnote}{0}%
\opage{138}and divides, it will bring light from above and forces
from below into such an interaction that the spiritual
sap of nourishment comes into circulation, that we, with all the
tones of color of an all-sided enlightenment, with all the shadings in
all transitions from the obscure to the clear, from the
simple to the manifold in knowledge and insight,
constitute an enlightened people united in the recognition of truth
and reality.

That the interaction will be fruitful not only
for the Folk High School but certainly also for the University,
is a conviction which in a striking manner has
confirmed itself to me through the spiritual alertness and
unbiased receptiveness that I have found in this circle.
I thank the honored assembly, ladies and gentlemen,
for kind attention. I shall not
forget that if what has been said here has, as I hope,
been understood as it has been heard, then it is
owing for a good part to the folk high schools which for a
series of years have worked in this region. The thanks which I
bring you are well founded; I am grateful for
the impression that I have here received of an open ear, of
an awakened sense and above all of a fresh, blossoming
popular life.

% --- p. 139 ---
\setcounter{footnote}{0}%
\lecture{Ninth Lecture.} {(Fredericia).}
\opage{139}We are a small people. ``Yes, alas,'' one answers
with a sigh, ``we are too small to assert our right.
Had we an army of 800,000 men and a fleet
to match, how it would brighten up in the old and new papers
of the law of nations. The moral effect alone
upon the diplomatic corps would be edifying
to behold. The fairness and obligingness,
the honest will and sincerity of heart with which
one would then, from all sides --- for the sake of peace and concord
and good understanding --- offer
one's bona officia, would surely have to make a touching impression
upon every friend of man; a restless zeal to get
torn treaties healed, sworn promises fulfilled,
forgotten indemnities paid, etc., would surely
give the now doubting Europe a proof as striking as it
would be surprising that there is still honesty,
fidelity to one's word and uprightness in the political world.
But, alas, we are a small people! That we should now just
lack the quantum of iron and lead which is so
necessary for setting justice in motion,
is surely as though a curse rested upon us.''
% --- p. 140 ---
\setcounter{footnote}{0}%
\opage{140}This way of looking at things is not only oppressive, it is
also one-sided and false. If it is already a questionable thing
when an otherwise well-endowed person has become so angry
with fate that he no longer cares to use the
powers he has, then the matter becomes still more questionable
if a well-endowed people were to fall into
a similar despair. Heavy burdens
may rest upon the small peoples; but it must not therefore
be overlooked that those which rest upon the great
may be, in proportion, just as heavy. Let us in this
connection regard our conditions from the bright side
and consider how much blessing may go with
being a small people.

A small people, however courageous and brave it may be,
cannot, especially in our days, expect to reap such
laurels in war as a great people; but can it
therefore do nothing to defend its independence?
The French people is, in the nobler sense of the word, assuredly
a great people, a people that has shone and in the future
will surely shine again on the stage of war; but
has it nevertheless not had to learn what defeat
and humiliation are, quite as well as we? There is for
the selfish element in human nature undeniably something satisfying
in setting oneself above the prejudices of morality
and making great conquests; there is, when the passions
are inflamed, something still more satisfying in
letting the sword of retribution flash and taking vengeance for
defeats suffered. Such satisfactions are reserved
for the great peoples alone. Should we now have to admit to ourselves
that the monstrous human slaughters with which the lust for dominion
and vengeance feed one another in turn belong to
the noblest and best things for which it is worth living,
% --- p. 141 ---
\setcounter{footnote}{0}%
\opage{141}then we should undeniably have cause for complaint; but if
the cause of humanity can also be advanced in a peaceful
way, then a small people has nothing in this respect
to complain of. A great people can sacrifice human lives
by the thousand without thereby putting its existence
at stake; it has masses enough to serve
as cannon fodder; but a small people must stake its life when
its independence is threatened. For a small people, then,
a war of defense acquires a quite different moral, personal
significance from what a war of conquest or a war of vengeance
acquires for a great one. In this moral, this personal element
there is a blessing.

A brave but otherwise peaceable knight does not
ask how great the probability is of winning
when his honor and person are attacked; he risks his life,
for honor and life are inseparable in him. He does not skimp
on a shilling when the question is of a
sword and a reliable suit of armor, for sword and armor
make up an essential part of his real personality.
He does not regard the time and diligence he has
had to spend on his exercises in arms as wasted;
he does not suppose that anything is thereby lost to his
peaceful life and pursuits, for he knows that a life in peace
and quiet is worthless when independence is lacking.

A small people must be like a peace-loving but
brave and well-equipped knight. It does not ask
how great the probability is of winning when its
honor and personality are attacked. It knows within itself
that true independence stands independent of victory
and defeat. So long as the majority of the population only
felt themselves to be blindly obedient subjects under a
self-willed state, it might seem natural that the citizens
% --- p. 142 ---
\setcounter{footnote}{0}%
\opage{142}should leave defense to the state, regard war as a
sheer misfortune and peace at any price as an unconditional
good. If the state could not defend itself, what was
that but that one got a new government; one
could always be sure of keeping the old burdens
and vexations. That was the political view
in the fable of the horse and the ass, repeating itself with a
slight change. Otherwise, when the state is popular, is
the state of the people; such a state is an expression of the people's
real personality, and personality must be defended.


A great and mighty people can be at one time warlike,
at another peaceably disposed; but it has, on account of its
massiveness, difficulty in making the opposite moods
go together into unity. In a small
people, which can wage war only for the sake of peace, the spirit
of the people must at all times be in equal measure both
warlike and peaceable. The independence of the small people
is not, like that of the great, now and then exposed to
danger and then again secured; no, the small people is,
essentially viewed, always in danger, just because it is small. It
must always be watchful, always in tension. But
precisely in this watchfulness, this tension there is
a blessing. The works of peace and the preparations
for defense of life and death ought to go hand in hand
with each other.

It is not overstrain, not a screwed-up chivalry,
that is hereby recommended. The screwed-up chivalry
is on the contrary on the side of those who, without a sense
of the real conditions of self-defense, sing of the berserkers of old
and the heroes of yore, supposing that it does not properly
depend on arms and exercises in arms, but
% --- p. 143 ---
\setcounter{footnote}{0}%
\opage{143}solely on ``the word and the spirit.'' That is lack of sense, that % defect: Utorstand (a broken f)
is hollow enthusiasm. To send unpracticed or only
half-practiced soldiers into the field is, to put it plainly, to
send the sons of the fatherland blindly to death. No, better then
to join those who take the matter materially.
There is after all always a sense of reality in working for
the nourishment of the population; one can after all in the end --- out of respect
for the real --- bid the independence of the people good day
and set the porridge pot higher than honor.

For an all-sided carrying through of the principles of freedom and
equality, the historically inherited
order of society undeniably puts great obstacles in the way. But
to charge the existing disproportions exclusively
to the account of the institutions of the state and to think one can get at
inequality by raking over the existing order,
is an enormous error. Inequality is more than arbitrary;
it is deeply rooted in the natural conditions of man.
Were it to be wiped out, it would have to happen by
so unnatural a compulsion and so merciless a despotism of society
that Spartan severity would be
child's play in comparison. How sharply
the classes of society, with the highly unequal distribution of
goods and rights, can stand against one another,
even at stages of cultivation that belong more to nature than
to culture, history shows in manifold examples. We
shall not here expatiate on the caste system in India
or in ancient Egypt, but hold to our own
past.

The old Norse poem \emph{Rigsmaal} describes for us
the origin of the different estates. ``The god Heimdal
had once, under the name of Rig, set out on a wandering.
He came in to a pair of old poor people,
% --- p. 144 ---
\setcounter{footnote}{0}%
\opage{144}Ae and Edda; it was wretched with them, the bread was baked
in the ashes and full of husks. Some time afterwards the
wife bore a son, who was sprinkled with water and called
\emph{Træl}. He was dark of color, the skin on his
hands was coarse, the knuckles thick, the fingers clumsy,
his complexion sallow, his back crooked and his heels
long. He grew up and learned to use his strength.
Once a woman came to the house with a heavy gait,
sunburnt arms and blisters under her feet. She was called
\emph{Thy}, that is, bondwoman, and married Træl. Their
sons, who were called Herdsman, Bungler, Clumsy, Sluggard,
Dawdler, Crookback and Drudge, busied themselves with
setting up fences, manuring fields, tending swine and goats and
cutting turf. The daughters were called Black, Lump, Crooked-nose,
Hunchback and Longshanks. From these the race of the thralls has sprung.
Rig wandered on and came to the farmer's
homestead, to a pair of happier people, Ave and Amma. A
fire was kindled in the middle of the floor. The man sat and
whittled wood for a loom; his beard was combed and
his hair cut across the forehead, his shirt was close-fitting and open
at the breast. The housewife twisted the thread and prepared the
yarn for the loom; her hair was twisted into a topknot,
her jacket stood open at the breast, and she had a collar % print: Brystet. og (misprint)
about her neck and buttons on her shoulders. The wife got a
son, who was wrapped in linen, sprinkled with water and
called \emph{Karl}, that is, Freeborn. He was a blooming
boy, light of hair and with sparkling eyes. He grew
up, learned to build, to forge, and to drive the plough. A girl
whose name was \emph{Flink} came driving to the farm in a goatskin
gown with the bunch of keys at her side. They married
and lived contentedly together. Their sons were
Man, Lad, Landowner, Peasant, Householder, Longbeard
% --- p. 145 ---
\setcounter{footnote}{0}%
\opage{145}and Brave-man. The daughters' names were Wordy, Bride,
Snow-white, Bashful, Finery-lover, Of Age and Woman.
Rig then came to a building whose door faced
south; in the door, which was shut, there was a ring, and straw
was strewn on the floor. The man and the housewife sat
looking at each other and playing with their fingers. The man
busied himself with weapons, and the wife wore a ring ornament
on her breast; her brow was brighter and her neck
whiter than the newly fallen snow. Nine months after
Rig's visit she bore a son, who was wrapped in silk,
sprinkled with water and called \emph{Jarl}; he had light hair,
and his eyes sparkled like a serpent's. As he
grew big he learned to shoot with the bow,
to ride, to swing the sword and to sport in the water. And Rig
came to him, taught him to know runes and bade him
win lands. So \emph{Jarl} went forth, moved about in war
and won eighteen farmsteads as his own possession. \emph{Erna} (the Brisk),
daughter of Herse, who was slender of waist and whiter and
prouder than other maidens, became his wife, and glorious
sons adorned their marriage.''

A society with inequalities so ingrown, fixed and bound to nature
is, as far as the example of the thralls is concerned,
not far from the inhuman: might not something similar,
even if under other conditions, be
said with respect to the proletariat that in our days is constantly
growing? In every society, be it ever so
democratic, there will always be an inequality which either
already is or is at least on the point of becoming revolting:
against the excesses of inequality the people must unceasingly
be at its post. But in the nature of inequality there is
a secret.

Equality becomes inhuman when one wants to
% --- p. 146 ---
\setcounter{footnote}{0}%
\opage{146}abolish inequality; inequality becomes inhuman when
one will not carry through equality: \emph{it is on the
foundation of inequality that equality must be carried through}. The torments of society arise from the fact that this simple
proposition is misunderstood from two opposite sides.

Those who have an interest in maintaining inequality
hold one-sidedly to what is necessary by nature, with all the
arbitrariness flowing from it. That the ladder of society has
many rungs, that on account of the division of labor there must
of necessity be different classes of labor, higher and lower,
that the inviolability of the right of property, the heritability
of fortune, different ability, different luck, etc.,
must produce inequalities which no human
hand is able to wipe out, of all this the defenders
of the existing order are excellently informed. But
when the matter takes on the appearance that inequality without
more ado is the solid, the rightly reasonable, whereon one
may safely rest; when it is again and again forgotten
that inequality, held fast for itself, is inhuman,
that inequality exists only in order that equality may be realized
in an ever richer measure: then the outlook becomes
unreasonable, and the party that fights under its banner,
intractable.

Those who suffer under the pressure of inequality and for that reason
are interested in asserting equality have
weighty grounds to adduce. ``Can it be called an age of freedom
and equality when one assures the possessors
of letting the untiring capital work for them, grow,
unwearied amid the progress of labor, while thousands
and thousands of real workers must wear themselves out
without attaining the bare necessaries, because the economic question
is not about the worker, but about the work, because
% --- p. 147 ---
\setcounter{footnote}{0}%
\opage{147}capital, this impersonal thing that sucks the blood of all, has become
labor, and labor capital? A horse one may value
and treat according to its work, for the work is
not for the sake of the horse, but the horse for the sake of the work.
Is public morality perhaps to be glorified by the
profit of political economy: to set a man who works
on a level with a horse? Can it be called
respect for the personality of man when barren abstractions,
general products of the exertion of real men
in the sweat of their brow, --- for nothing else
are capital and labor --- are to have a right
to use the working persons as living machines?
Is it perhaps the prerogative that the law grants to the living
over the lifeless, that capital is bound to
repair the dead machines as they wear out,
while the living may be thrown away when they are no
longer able to repair themselves?'' Here are
questions which, when put in practical form
with an eye to the particular case, may well embarrass political economists
and make civilized
society with its ``rights of man'' blush.
We could give the champions of equality our undivided
sympathy, if only it were possible to get the
error rooted out which weakens and hampers the whole
endeavor. The fundamental error lies in this, that one,
instead of recognizing the significance of inequality as the foundation
of equality, wants to have inequality itself pulled up
by the root. A sad error! The equality that
is not able to take the element of inequality up into itself is, and
remains, utopian.

Among the animals that live in society, the Lord of Nature
has seen to it that the individuals of different classes,
% --- p. 148 ---
\setcounter{footnote}{0}%
\opage{148}where such are found, exactly perform the work
corresponding to each class and in accordance with this receive
their share of the whole product of society. Ants,
termites, bees, etc., furnish examples of how
inequality can be reconciled in a natural way with
equality. The condition of the union is that each class of
individuals has by nature so peculiar a stamp
that it is driven from within to perform its peculiar
work and for that reason finds itself in such a dependence
on the whole society that it cannot possibly tear itself
loose from it. It is such a dependence on particular
natural dispositions to which the description in the cited
``Rigsmaal'' points. The thrall, the farmer, the Jarl all
have their origin from the visit of the god; but their
natural endowment is so different that one can even
see directly from their exterior to which class each
among them originally belongs. This explanation
is thoroughly false; inequality here gets so terrible a
right of the firstborn that equality cannot possibly come forth.
To be sure there is a lower and a higher natural endowment, but it
applies to individuals, not to classes; history has shown
that there can be clever thralls just as well as stupid
Jarls. From another, more exalted position inequality has been
defended by appealing to the express
will of God. It is then, according to the official
religion, to be quite settled that God himself has willed
such a separation of the estates, that one estate is set
in eminence, another in lowliness, that some men are
by Providence destined to toil and poverty, others to
enjoyment and riches. The fruits of this official interpretation
have shown themselves in various lands. The neglected
proletariat, whose ignorance of the spiritual
% --- p. 149 ---
\setcounter{footnote}{0}%
\opage{149}is in many places boundless, has, understandably enough,
confounded the outward ecclesiastical system with religion
itself and let itself imagine that faith in a living
personal God was the same thing as faith in a superhuman
despot who took the side of the great and
helped them to subdue the small. This imagination
must be rooted out from the ground up, if it is to
succeed in any degree in carrying through equality.

What is peculiar to the religion of Christ is that
it lays weight both on inequality and on equality;
it not only lets the existing opposition between
high and low, mighty and humble, masters and servants,
rich and poor stand, it goes further and refers
all this inequality to a divine order. But it
must above all not be overlooked that it regards
the presupposed inequality precisely as a condition for carrying through
equality. It speaks to master and servant, to
rich and poor, and what does it say? ``Remember that
ye are equal before God, that ye are brethren.'' On this condition
the servant has just as much obligation toward the master
as the master toward the servant, the poor just as much
obligation toward the rich as the rich toward the
poor. But the spirit of religion knows the world and knows
how hard existence is; just therefore it lays two
times, three times weight on the obligation of the rich
toward the poor, two times, three times weight on the obligation of the master
toward the servant, on the obligation of the employer
toward the worker. If thou knewest, poor worker,
who He is who has spoken and still speaks
thy cause before God and men, and with such
authority that thy personal worth, thy inalienable right
shall stand immovably firm, that thou in spirit and being canst
% --- p. 150 ---
\setcounter{footnote}{0}%
\opage{150}be the equal of the rich and mighty of the earth: in truth,
thou wouldst never forget him!

With the religious-moral recognition of the worth of man
the political-civil leveling out of the disproportions
which inequality introduces must be unceasingly striven for.
It ought to be admitted from both sides, as well from theirs who
feel oppressed as from theirs who are favored,
that the knots of inequality cannot be untied
as one unties a bow. There must be on the one side
a willingness to make sacrifices and on the other
prudence enough to give time and to be content with
something, until more can be attained. The main thing is
that a plan be laid for the endeavors as a whole
which both parties must acknowledge as appropriate;
that work be done continuously step by
step, so that what goes before prepares what follows;
that something really be accomplished by way of civil legislation.
Here, if anywhere, it can be seen how intimately disposition and action
belong together. If disposition is lacking, then all practical
proposals are in vain; if action is lacking, then one
may reason and moralize as much as one pleases: it
leads to nothing.

If we now look around for the most favorable conditions
for a real solution of the tasks of equality,
might it not turn out that a small people has
essential advantages over a great one? In a great people
the different classes of society will not only keep
apart; but each class will even, on account of
the greatness of the masses, the multitude of individuals, constitute
a world of its own. Two classes of the same population
can in manner of living, colloquial speech, customs and habits stand
% --- p. 151 ---
\setcounter{footnote}{0}%
\opage{151}farther from one another than two persons of different nationality,
when they find themselves at the same stage of cultivation.
Admittedly, mixtures, transitions and fusions
will take place in many ways; but the elements of inequality
with what is foreign, what separates
and repels, crop up in so many
forms and hold to so hard an immediacy
that the ethical-practical working through becomes, if
not impossible, yet exceedingly difficult. ``There are,''
says a witty French author, ``two peoples in
France: the provinces and Paris; the one eats, sleeps,
yawns, listens; the other thinks, dares, watches and
talks. The first is drawn along by the second like a
snail by a butterfly, being by turns amused and
disquieted by the caprice and daring of its leader.
It is the leader, then, to whom we must look.'' The last
years' events have shown in a sorrowful way that the
moral, political and social discord between the particular
classes of society in the great country by no means confines
itself to the opposition here indicated between
province and capital; it branches so far and
cuts so deep into the roots of society that it will surely
be a long time before equality in the fatherland of the rights
of man is reconciled with inequality and the contending
parties are made one.

With us the conditions are differently favorable, that is to say,
when the popular will seriously wants to use them. It is a
fact that the smaller a company is, the more weight
can be laid on its single members.
Precisely the circumstance that the classes of society with
us are too small to constitute worlds of their own, closed off,
must favor an approach, a so to speak
% --- p. 152 ---
\setcounter{footnote}{0}%
\opage{152}personal understanding between them. One individual of a hundred
is, in the capacity of individual, ten times as much
worth as one individual of a thousand, and ten thousand times
as much worth as one individual of a million, according to
the common calculation of masses. But what is a mass?
The word mass is precisely characteristic of matter, of the
dead, the inorganic, the lifeless. Matter consists of
indivisible parts, whose collected quantity makes up the mass.
The more parts, the greater the mass, the greater the mass,
the greater the force of attraction. The globes in space
act in proportion to their masses, in that they affect one another
according to the law of attraction: in this region the
great masses are great powers. If this view is transferred
to human societies, then human individuals become
parts of mass, and the law repeats itself: the more
individuals, the greater the mass, the greater the mass, the
greater the force of attraction. But who does not see the
danger, the corruption for popular life, in the dead
law of attraction being applied to personal beings. The
less feeling of independence there is in the single
individuals, the stronger the inclination they have to
flock together in a crowd and lump themselves together in a mass.
On these conditions the great people will always exert
a weighing, oppressive influence upon the small. Not
only in war, also in peace. A silly fashion, a thoughtless,
but lewdly glittering, nerve-titillating novel,
e.g., goes out in Paris, and what happens? The new
phenomenon is not judged according to its inner worth;
far from it! Though criticism also does its duty, it does
no good. It is voracious masses that hold the scales,
mass-individuals that pass the judgment. The novel comes
out in seven editions; the great nation, it is said, has approved
% --- p. 153 ---
\setcounter{footnote}{0}%
\opage{153}its significance for the history of culture. What follows from this?
``That a translation is naturally necessary in order to
enrich our literature; we dare not stay behind'': the
small people is quite impressed.

A small people cannot fight its way through by
means of masses; it has by Providence just been placed
under such conditions that sound sense must make it
attentive to the difference there is between
people and mass. The mass as mass is inorganic and
so far always crude; in the people as people everything is organized,
living, personal: the more the mass spreads
out, the more the people is narrowed; the more the people
expands, the more the mass is narrowed; or:
the less mass, the more people, the more mass,
the less people. But to be placed under such
conditions that for the will of spirit there is precisely the very strongest
summons to get the individuals developed
to personal independence, to get the mass organized
and transformed into a living people, to get the whole permeated, % print: gjennnemtrængt (misprint)
inspirited and moved by one thought, one will,
one mind, in this there is then no curse; no, on the contrary,
the greatest blessing.

Our fatherland, which in earlier days has at times
been near to perishing, but nevertheless in a wonderful
way has been saved, has, in this century in particular,
suffered rough treatment. Through all this we have yet
retained a good which is beyond price. There are in history
two great laws which, because they are laws of spirit and of freedom,
in the clashes of finitude so easily come
into conflict with one another, and yet must constantly be brought
into agreement, namely: the Law of \emph{Continuation},
a law for even development, and the Law of
% --- p. 154 ---
\setcounter{footnote}{0}%
\opage{154}\emph{Progress}, a law for the breakthrough of new ideas.
At every decisive turning point there shows itself
a tension between these two laws. In France, in
``the great nation,'' the Law of Breakthrough has gained so
decisive a preponderance that the foundation of legitimacy is
undermined. Every idea of society that is able to
inflame the multitude and gather the masses around it
has, because the historical Law of Continuation is broken,
fundamentally the same right as what exists for a moment.
The ship of state is holed in the keel and must constantly be kept
at the pumps. How can unity of the people come
to strength in a country where three dynasties and two republics,
at the same time in full activity, fight one another
to the death, in that each of these powers ---
as circumstances may be --- considers itself fully entitled
to seize the reins of government. We shall
not lord it over ``the great nation''; we shall
with admiration acknowledge that it, even in its defeats,
has displayed astonishing powers; we shall --- for the sake of
``the rights of man'' --- believe and hope that
the people which amid terrible crises and excesses
has been able to hold the banner of freedom aloft will continue
its course and reach its goal, to its own honor
and for the benefit of the civilized world. But
when we consider what it means to venture such
breaches of the Law of Continuation (Continuity) as have
here occurred, then we must count ourselves fortunate.

The remarkable thing about our constitution of 1848 is precisely
this twofold character: that the change is so thoroughgoing,
and its introduction nevertheless so lawful. It casts
a historical luster on Frederick the Seventh and his
men that continuation and breakthrough here show themselves
% --- p. 155 ---
\setcounter{footnote}{0}%
\opage{155}as two sides of the same thing. It is a foothold for
those who hold to what exists; the lawful way,
the even, quiet transition is fixed as a principle.
But it is also a support for those who boldly and
confidently work for progress; for from the principle it is seen
that one can --- understand, in sober earnestness --- take
a matter from the root up, without bringing about
a violent, lawless upheaval.

So we have, then, what is a blessing upon
our small people and makes ``the heavy unity'' light, with
the assistance of Providence all conditions for an energetic, yet
calm development of a sound and vigorous popular life in
Denmark.

Theoretical surveys must not be mixed up with the
practical tasks; but as surely as disposition and action
unceasingly condition one another, so necessary is it that the
people be sufficiently enlightened about its own conditions by
being enlightened about the conditions of human life as a whole. It
is teachers of the people, well-prepared, well-equipped teachers of the people,
on whom in this respect first and last everything
depends. We close by recalling the words of the great
teacher of the people himself: \emph{The harvest is great}, but the \emph{laborers
are few; pray} the Lord of the harvest that he
\emph{send forth laborers into his harvest} (Luke 10:2).

% --- p. 156 ---
\setcounter{footnote}{0}%
\lecture{In Memory of the Battle of Isted.} {(in Aagaard Wood).}
\opage{156}Today twenty-five years ago there was a solemn day here in
the country; a festival was held then, --- a bloody festival
on the brown heath of Isted; it was a sacrificial feast in the
old Norse style. We all know, of course, what it was about. It
was about a joust against the aggression of defiant arrogance, of
intrigue, that breaks faith and laws; it was about defending
our boundary-mark, our old right, our young freedom, our
mother country's independence and honor. To trample the serpent of rebellion underfoot and raise the Dannebrog
high was the goal of the struggle; and the goal had to be attained, for
it was the will of the King, the will of the people, and the will of the soldier.
So it went forward; despite resistance, rain of bullets, and
smoke and blood, yes, even despite the whims of fortune, our
brave men went boldly forward to the goal. It was a day of honor
for Danish hearts and for Danish lions. Schleppegrell
and Læssø! these names Saga has carved upon her shield;
there they shine with undimmed splendor in remembrance
among a grateful people. Yet it was more than
renown that was sought. The hero who
fought the great fight did not look too high, did not take aim
at the stars of fame; he was a sturdy, modest
% --- p. 157 ---
\setcounter{footnote}{0}%
\opage{157}hero: the honor of the colors, the girl's smile and the King's thanks were
all his reward; it was the brave Landsoldat who
brought his festival sacrifice here.

A glorious victory without bluster or boast. In the jubilation of
joy there was no arrogance, ah no, but a quiet sigh, a
melancholy blinded with tears. The news of victory struck like
lightning; but then there followed with it --- the dear reckoning,
the long, the terrible list. Everyone who has
lived through it knows well what it means: in anguish of soul,
with beating heart and trembling hand to run through
the list of wounded and fallen; everything grows black
before the eyes at the thought of finding there --- ah, precious
names. Of joy's play of colors and sorrow's shadows
Time, with a gentle hand, has little by little brought an equilibrium
about and painted for us a picture in memory.

Upon the victory of Isted there has, alas, followed a defeat, a
loss, that hurts; but honor is unstained. Superior force
broke through and tore our land asunder; but the Dannebrog
is whole still, in this cloth there are no rents. To
engage one, to stand against two, to give way a little
with the foot before three and without shame to flee before % defect: tor (a broken f)
four, that was
King Frode's law of manly courage, and this law the
Landsoldat has faithfully fulfilled. If we consider the superior numbers
against which we had to fight in the last battle, then
our defense is to our honor; for here there were more than
two and three against one, and our enemies, by the
laws of chivalry, have no particular cause to boast of the victory.

The struggle is still going on; there is still a
victory to be won for our independence, though with other
weapons and under other conditions. Independence
is, after all, of different kinds, partly something outward, partly something
inward. A great country with a great people can hold
% --- p. 158 ---
\setcounter{footnote}{0}%
\opage{158}itself independent outwardly, so long as its military power
is able to subdue and defeat the outward enemies; but
that is not to say that it is also inwardly an
independent people; far from it! It has often been seen in history
that such a people can be in much discord and
dissension with itself, that precisely for the sake of greatness and
outward independence it must bear heavy chains.
Also a small country with a small people can preserve its
independence outwardly, if Nature, with rare
favor, by steep coasts, high mountains, inaccessible regions,
has made it safe against attack; but, even surrounded by
a Chinese wall, such a securely protected small people can yet
inwardly be highly unindependent; it can split into
irreconcilable parties, it can consume its strength in an empty,
unfruitful bustle, it can wear itself down in petty
squabbling. For outward independence no
human power dares to vouch; even the strongest can
stumble: great kingdoms, great thrones fall. Over the
fate of the peoples in the course of the ages rules a guiding Providence, and,
let it be our consolation, Providence's eye looks as mildly
upon the small as upon the great.

What matters, then, is to win inner independence;
but to inner independence belongs concord,
mutual concord in the spirit of truth, justice, freedom
and love. Popular life is not a life in
Paradise; the popular soil must be cleansed, plowed
and cultivated, it costs labor. Around the tree of freedom
there are roses and love tokens, yes! but also thorns and
thistles. It is impossible that independence and freedom
can develop in a vigorously growing people without
inward tension and strife. So view stands against
view, party against party, and opposed endeavors clash
% --- p. 159 ---
\setcounter{footnote}{0}%
\opage{159}with one another. It is necessary that amid such
clashes there be a strong foundation of unity and concord.
If the foundation of concord breaks, then the independence of the people
is in danger.

Who among us has not experienced how difficult
it is to bring two persons, even if each, for his own part,
sincerely wishes what is right, to come to agreement,
when there is suspicion and secret grudge between them.
In public affairs, when the parties contend, it may
often seem to an outsider as if the clear
cause, the right view stood mediating in the midst of the strife;
but all at once it has vanished again, a But, a side glance,
a single word has chased understanding away: it is not
dishonest intent, not conscious duplicity, it is
the serpent of grudge, it is the prompting of suspicion. What
is to cleanse out suspicion and stifle the serpent of grudge? A
deep, strong and living feeling of what the people is,
and what the fatherland requires, when great interests
are at stake. Also in this respect the memory
of the Battle of Isted can be a blessing to us. There was concord,
yes there was unity, one spirit, one soul, one
watchword. ``With God, for King and Fatherland!'' that was
a word which gave the soldier courage; it dwelt in his silent
mind, it gave him certainty of his calling, it kindled
the flame in his breast, it sounded with hurrahs in forest and
vale, it roared across the heath. Could we see
these sacrifices, these dead, these dying, see the blood trickle
out from the springs of the heart, without forgetting our everyday squabbling
and feeling lifted above petty willfulness,
suspicion and grudge, then we would be an ungrateful people;
the memory of the warriors, the spirit from the graves of the heroes would have to
% --- p. 160 ---
\setcounter{footnote}{0}%
\opage{160}bear witness against us that we gave up our best possession, our inner
independence, and were not worthy of freedom.

In the great concerns of the people and the fatherland
concord is, on the whole, more than a matter of feeling and more
than a calculating consideration of the understanding; it is a matter of will;
there shall be, there must be concord on what is most important.
The man who cannot become of one mind with himself
never becomes independent, and the kingdom that is
at odds with itself must fall. To concord belongs
will; to a concordant and independent people a clear and
mature popular will. Measures for the fatherland's
defense, the first conditions for the many-sided development
of popular life in spiritual as in material respects, belong to
the things that tolerate neither evasion nor delay; here
is something that cannot wait until all byways have been
tried, and all means of resistance on the part of the contending
parties exhausted.

If the work of independence is not to come to a standstill, a
watchful eye must be kept on a twofold danger. The one
is: to give way too much and let oneself be blinded by empty
promises and meaningless agreements. The other is:
to press too hard and, cost what it may,
to carry through, as the saying goes, one's own war. The latter is
just as ruinous as the former. By using
means too strong the leaders of the people, with all their
zeal and good will, may easily come to behave like
those quacks who, in order to cure a single
injury in haste, ruin the whole body and kill the patient.
Also in this respect the men of peace have much
to learn from the warriors. --- Look to the right! Look to the
left! --- Straight ahead! that belongs to the young soldiers'
first drills. If the way of truth lay open for
% --- p. 161 ---
\setcounter{footnote}{0}%
\opage{161}us, like a ditched highway, we should certainly need
to look neither to the right nor to the left; but since we
must constantly fight our way forward; since the road may wind
through the impassable ground of errors, and the truth lie
now to the right, now to the left, it matters
to have an open eye for both possibilities: a plain trotting
``Straight ahead!'' could easily deceive us through illusions
of our own infallibility. And one thing more. In
the tumult of battle a double word of command sounds, now
one: Forward, on! now again a Halt, stand! The
latter is for the soldier fired up for combat often
harder than the former. To go on, when everything is staked,
and passion gives speed, comes naturally; but
to stand still in the midst of a rain of bullets, see comrades storm
and fall, without daring to hasten to their aid, costs
self-conquest. None the less it is, for the sake of victory,
just as important to obey the last order as
the first. To true courage in battle belongs self-control;
should anyone doubt it, let him consult
the heroes of Isted.

When a people is in earnest to fight for its independence,
it is not enough that a leader sets himself
a goal, gathers a crowd around him, a majority, and
says: we are the people. It is indeed possible that such a
majority will may be a first, beginning expression
of the true and proper popular will ; yes, it is possible,
but it must nevertheless be tested. By what is it then to be known ?
It is to be known precisely by this, that the goal that is
set embraces equally everything and everyone. To aim upward in such a way
that it looks as if the whole people were especially to
work for some higher institutions of learning, some
learned and cultivated persons, whose learning and cultivation, it does not
% --- p. 162 ---
\setcounter{footnote}{0}%
\opage{162}understand, is unhealthy. To aim downward, and in such a way
that it always looks as if it were only the material,
carnal well-being, temporal fatness, that was
the one thing needful, is likewise unhealthy. In both cases we may
be assured that it is a party will we have
before us, no true popular will. Where a society suffers
from the fundamental ill that a single party, be it a
minority or a majority party, rules one-sidedly, there
the popular will is paralyzed. Above all the small peoples must
guard against deifying the majority will; for by such
a deification they sign their own death sentence.
If a majority is right merely because it is a majority, what
follows from this? This follows from it, that, since the
small peoples are always a minority and the great a majority,
a small people has --- as a great power is fond of asserting
--- really no right to exist independently.


As the goal is, so must the means be; by the
choice of means and by the use of them we recognize again
the difference between party will and popular will. We remember
King Solomon's judgment, which ordered that the living child
be divided. But the true mother would rather leave
her child to the false one than suffer it to be
divided and killed. Show me for once a majority will which,
although it is conscious of its right and can testify to it
loudly, yet for the sake of life and concord, like the child's
true mother, exhibits the self-denial to yield
for a moment, so that nothing existing shall be violated
and no necessary intermediate step skipped over! Yes, when
a majority will is so high-minded, then it is sure
of itself; it wins the respect of its opponents, it
wins more than that, it wins the heart of the people: it
% --- p. 163 ---
\setcounter{footnote}{0}%
\opage{163}will surely appear that it is and remains the true, independent
popular will.

Today twenty-five years ago the spirit of the people held
festival in Denmark. Treason, which always
speculates in party spirit, wanted to split us apart and cut
a piece from our own flesh; but the King raised
his voice: \emph{it shall not be}! and the royal
will was a popular will. Alas, when that voice
fell silent in death, then there happened what we had long
feared. But the word from the royal lips is
not dead; it summons us to be at our post, it
still sounds in our ears, a voice of the spirit, a word of prophecy,
confirmed by God's Providence.

There is an independence that cannot be robbed;
there is something in a noble people that cannot be shot down
with cannon. If we still strive for the cause of popular life,
so that our watchword is still as before: with God, for
\emph{King and Fatherland}! then let only those who trust
confidently in their own strength set Right on the spear-shaft for long enough;
let them only think that the superior power of a majority is
enough to force a foreign tongue, a foreign nature
and a foreign soul into a people which has honestly resolved
to live its own life: every earthly superior power
has its limit after all; there is a word of Godhead
that says: thus far, and no farther! We have no right
above other peoples, whether great or small, to call
ourselves God's own people; but we have both the right and the obligation
to lead our life and fight for our independence
in such a way that we may with reason be said to
work for a sensible and many-sided development of noble and
true humanity. If our cause becomes in a deeper sense
a cause of humanity, then we may believe and
% --- p. 164 ---
\setcounter{footnote}{0}%
\opage{164}hope that the cause will prosper, that no superior power
will dare the inhuman. It is after all the cause of mankind
and the victory of humanity to which history
points; but the victory of humanity is inseparable
from an independent popular life. Away then with
anxious forebodings and heavy dreams! every threat against
our land and people shall be blown away; they shall not cow us,
``it shall not be''. We believe in God and in our good
right. Justice has a mighty voice, and the word
gives an echo here in these valleys; from field and forest, from
land and sea, from Danish hearts and from Danish lions,
from the tone of our song, from the voice in our breast, from
the vault of heaven and from the graves of the heroes of Isted echoes
this word: \emph{it shall not be}! With this faith, in this
hope we bless the memory of the fallen, we bring the still
living our thanks, and --- long live old Denmark!

\end{document}
